\quellseite{947}
\begin{schriftzitat}
V. 43. Jauchzet, ihr Heiden, mit seinem Volke; denn er wird das Blut seiner Knechte rächen, und wird sich an seinen Feinden rächen, und sich versöhnen lassen dem Lande seines Volks.
\end{schriftzitat}

77. Das ist, ob er wohl sein Volk hart strafen wird um ihrer Sünde willen, so werdet doch ihr Heiden euch noch freuen mit seinem Volke, der reichen Gnade, des Segens und Verheißung theilhaftig werden. Denn es wird noch ein Häuflein überbleiben, das recht Israel ist, recht gottesfürchtig und gläubig, als, die Propheten, Apostel und ihre Jünger, in denselbigen wird er nichtsdestoweniger (wenn solcher großer Sturm und Wetter solches schrecklichen Zorns vorüber ist) dennoch alle selige Verheißung der Gnaden erfüllen, und durch sie aus den Steinen, das ist, aus den abgöttischen Heiden, heiligen Abrahams Samen und Kinder machen [Matth. 3, 9.].

78. „Und wird das Blut seiner Knechte rächen“; das ist, aller frommen Propheten Blut, und hernach der Apostel, und alles unschuldige Blut, so dieses Volk, ihre falschen, abgöttischen Gottesdienste zu erhalten, mannigfaltig vergießen wird, wird GOtt schrecklich rächen, von Habels Blut an, bis auf den letzten unschuldigen Blutstropfen, der vergossen wird. Und wird doch „gnädig sein dem Lande seines Volks“, das ist, dem übrigen Haufen der Israeliten, die da glauben werden. Denn wird Gott ein neues und ewiges Reich auf dem Stuhl Davids und im Hause Israel erwecken, welches kein Ende haben wird, da der HErr und König sein wird der HErr Christus, welches nun durch das Evangelium offenbart ist, dem sei Lob und Preis in Ewigkeit, Amen.

Ende des Liedes Mose.

\part{V. Auslegungen über das zweite Buch Samuelis}

\chapter[Von den letzten Worten Davids, 2 Sam. 23, 1–7]{Von den letzten Worten Davids, 2 Sam. 23, 1–7.\footnote{In seiner zweiten großen Schrift wider die Juden, „vom Schem Hamphoras“, sagt Luther, daß sich die besten Uebersetzer der Bibel zu sehr von den Rabbinen haben beeinflussen lassen, und dadurch der Analogie des Glaubens gefehlt. Dann fährt er fort (Walch, St. Louiser Ausg., Bd. XX, 2107, Z. 19): „Denn ich auch selbst ihrer Translation etwa zu viel gefolgt, daß ich muß widerrufen, sonderlich 2 Sam. 23 in verbis novissimis David, wie ich bald thun will.“ Dies Versprechen erfüllte Luther sofort. Noch im Jahre 1543 erschien unsere Schrift zu Wittenberg bei Nickel Schirlentz unter dem Titel: „Von den Letzten Worten Davids. D. Mart. Luther.“ In derselben thut Luther auf das herrlichste dar, daß nur diejenigen das rechte Verständnis der heiligen Schrift haben können, welche dieselbe im Licht des Neuen Testamentes, das ist, im Glauben an den einigen GOttes Sohn, Jesum Christum, betrachten. D. Caspar Cruciger übersetzte unsere Schrift in das Lateinische. Während dieser Arbeit ward er krank, vollendete sie aber dennoch am 15. November 1549. Tags darauf starb er. (Witt. [1549], Tom. III, fol. 130b.) M. Georg Röser nahm sie zuerst in den 1549 erschienenen dritten Band der Wittenberger lateinischen Ausgabe auf, dann aber kam sie auch einzeln heraus unter dem Titel: De novissimis verbis Davidis Commentatio pia et sancta, Martini Lutheri Theologiae Doctoris, Germanice explicata ab autore, et in Latinum sermonem conversa, studio et diligentia Casparis Crucigeri, D. Theologiae. Nunc primum seorsim edita, Lipsiae in officina Valentini Papae. Anno M.D.L. Cal. Maii. In den Sammlungen findet sich die deutsche Schrift: in der Wittenberger (1556), Bd. V, Bl. 536; in der Jenaer (1562), Bd. VIII, Bl. 137; in der Altenburger, Bd. VIII, S. 304; in der Leipziger, Bd. IV, S. 302 und in der Erlanger, Bd. 37, S. 1. Die lateinische Uebersetzung findet sich in der lateinischen Wittenberger (1549), Tom. III, fol. 91 b. Wir geben den Text nach der Erlanger, welche der ersten Edition folgt, unter Vergleichung der Wittenberger, der Jenaer und des Lateinischen.}}

1543.

\section[Vorrede]{Vorrede.\footnote{Diese Ueberschrift ist in der Wittenberger und in der Jenaer Ausgabe.}}

1. St. Hieronymus schreibt, es habe ihm begegnet, die Biblia aus dem Hebräischen aufs neue zu dolmetschen ins Latein, daß die Feinde Christi, die Juden, uns Christen verspotteten, als hätten wir nicht die rechte Biblia, welche dazumal im Brauch durch die ganze Christenheit ging, weil es an etlichen Worten und Buchstaben mangelte, die viel anders im Hebräischen stünden; welches
\quellseite{948}
zuvor auch andere mehr beweget hat, als, Aquila, Theodotion, Origenes etc., bis daß man zu der Zeit bei schierlei Dolmetschungen hatte, welche sie Hexapla nannten. Also ist jetzt, zu unserer Zeit auch so viel Dolmetschens, daß mit der Zeit vielleicht (wie dazumal geschah) so viel Biblien kommen werden, so mancher Klügel in der hebräischen Sprache Meister sein wird, daß da kein Ende sein wird.

2. Und das muß auch zuletzt geschehen, wenn man sich daran kehren will, was die Juden von unserer Bibel sagen oder urtheilen, welche untereinander selbst nicht eins, die Bibel also zerdehnen und zerreißen mit ihrer Grammatica, daß man freilich (wo man ihnen sollte folgen) nimmermehr zur einträchtigen Bibel kommen kann, weil ein jeglicher Rabbi besser denn der andere sein will; auch sie allein bekennen müssen, daß sie an manchen Orten die Worte nicht verstehen, viel weniger einträchtiglich, ohne allen Mangel, eine reine, gewisse hebräische Bibel haben, auch der Grammatica nach zu reden, geschweige der Theologie, darinnen sie doch zu gar nichts sind.

3. Darum sieht mich solch der Juden Gespötte nichts an, und um ihres Urtheilens willen wollte ich nicht einen Buchstaben kennen lernen in der hebräischen Sprache. Ursache ist die: Wir Christen haben den Sinn und Verstand der Bibel, weil wir das Neue Testament, das ist, Jesum Christum haben, welcher im Alten Testament verheißen, und hernach kommen, mit sich das Licht und Verstand der Schrift bracht hat; wie er spricht Joh. 5, 46.: „Mose hat von mir geschrieben; wo ihr Mosi glaubtet, so würdet ihr mir auch glauben.“ Item, Luc. 24, 44. 45.: „Es muß erfüllet werden, was im Gesetz, Propheten und Psalmen von mir geschrieben ist. Und öffnete ihnen den Sinn, daß sie konnten die Schrift verstehen.“

4. Denn da steht’s, da liegt’s, da bleibt’s: Wer diesen Mann, der da heißt Jesus Christus, GOttes Sohn, den wir Christen predigen, nicht recht und rein hat, noch haben will, der lasse die Bibel zufrieden, das rathe ich, er stößt sich gewißlich, und wird, je mehr er studirt, je blinder und toller, er sei Jude, Tartar, Türke, Christ, oder wie er sich rühmen will. Siehe an, was hat bei uns Christen den Ketern Arianis, Manichäis, Pelagianis\footnote{Erlanger: Pelagianus.} und Unzähligen mehr gefehlt? Was hat dem Papst gefehlt? Haben sie nicht die gewissen, hellen, gewaltigen Worte des Neuen Testaments gehabt? Was fehlt unsern Rotten zu dieser Zeit? Haben sie nicht das Neue Testament klar und gewiß genug? Sollte man einem jeglichen solchen tollen Teufelskopfe nach das Neue Testament verdeutschen, wie viel müßten wir wohl Neue Testamente haben?

5. Wenn es nun sollte Wünschens und Wählens gelten, entweder, daß ich St. Augustini und der lieben Väter, das ist, der Apostel Verstand in der Schrift sollte haben mit dem Mangel, daß St. Augustinus zuweilen nicht die rechten Buchstaben oder Worte im Hebräischen hat, wie die Juden spotten; oder sollte der Juden gewisse Buchstaben und Worte (die sie dennoch nicht durch und durch, allenthalben haben) ohne St. Augustini und der Väter Verstand, das ist, mit der Juden Verstand haben, ist gut zu rechnen, wozu ich wählen würde. Ich liebe die Juden mit ihrem Verstande und Buchstaben zum Teufel fahren, und führe mit St. Augustini Verstande, ohne ihre Buchstaben, zum Himmel. Denn ob Augustinus nicht kann, wie die Juden, sagen Kitajon, da er cucurbita sagt, Jona 4, 6., item, nicht kann sagen, veniet Hemdath, da er veniet Desideratus sagt, Hagg. 2, 8., und dergleichen viel, so bricht damit sein Glaube nicht den Hals noch Bein, weil er den rechten Mann kennt, der da heißt Weg, Wahrheit und Leben, von welchem die Propheten weissagen und zeugen; wie gesagt.

6. Wiederum, die Juden, weil sie diesen Christum nicht annehmen, können sie nicht wissen noch verstehen, was Moses, die Propheten und Psalmen sagen, was rechter Glaube ist, was die zehn Gebote wollen, was die Exempel und Historien lehren und geben, sondern die Schrift muß ihnen sein (nach Jesaia 29, 12. Weissagung) wie ein Brief dem, der nicht lesen kann, welcher sieht die Buchstaben sehr wohl, weiß aber nicht, was sie geben, wie das deutsche Sprüchwort sagt: Ein weiß Feld, darin ist schwarze Saat, manch Mann vorüber gaht, der nicht weiß, was da steht. Wer aber des Lesens kundig und fertig ist, der läuft überhin, faßt den Sinn, ungeacht, ob er etliche Buchstaben oder Worte nicht eigentlich ansieht. Ehe der andere ein Wort buchstabet, hat er den ganzen Brief ausgelesen. Also, ein Musicus hat ein Lied ausgesungen, ehe der andere sucht und findet, ob es ein Sol oder Fa im Clave sei.

\quellseite{949}
7. Man sehe den feinen Mann Lyra an, der ein guter Hebräist und treuer Christ, wie macht er so gute Arbeit, wo er sich wider den jüdischen Verstand legt, nach dem Neuen Testament. Aber wo er seinem Rabbi Salomo sonst folgt, wie kalt und faul geht es ihm ab, daß es weder Hände noch Füße hat; ob er wohl die Worte und Buchstaben gewiß hat, noch ist er ja besser und reiner, denn sie alle, beide, alte und neue Hebräisten, die zu gar sehr den Rabbinen folgen. Fürwahr, man darf den Fleiß nicht vornehmen mit Dolmetschen und Glossiren, wie man der Rabbinen und Grammatisten Verstand unter uns Christen bringe, er klebt ohne das von ihm selber allzu gerne an, wie Pech und Leim, wenn man sich gleich will vorsätzlich davor hüten. Denn die Buchstaben und Exempel der andern blenden die Augen, daß man den Sinn Christi zuweilen fahren läßt, da es nicht sein sollte, damit der jüdische Verstand also unversehens herein schleicht, wie allen Dolmetschern geschehen ist, keinen ausgenommen, mich auch nicht.

8. Summa, wenn wir unsern Fleiß nicht dahin kehren, daß wir die hebräische Bibel, wo es immer sich leiden will, zum Verstand des Neuen Testaments ziehen, wider den Verstand der Rabbinen, so wäre es besser bei der alten Dolmetschung geblieben (die doch das Beste und Meiste heraus hat, durch das Neue Testament), denn daß man so viel Dolmetschung haben soll (um etlicher weniger Wörter willen, die anders, oder noch unverdolmetscht sein sollen),\footnote{Diese Parenthese übersetzt Cruciger: propter pauca quaedam loca, ubi aut Hebraica lectio alia est, aut quorundam vocum interpretatio nondum certa est.} die dem Leser die Memorie irre macht, und sein Studium hindert, und ungewisser, denn vorhin, macht.

9. Darum habe ich zum Exempel vor mich genommen, die letzten Worte Davids auszulegen; nicht wie sie verdeutscht sind, da ich den andern allen gefolgt habe, damit ich nicht alleine klug wäre; jetzt will ich eigensinnig sein, und niemand folgen, denn meinem Geist. Wem es nicht gefällt, der lasse es fahren. Es ist nicht das erstemal, daß ich schreibe, was andern nicht gefällt, bin es (GOtt Lob) wohl gewohnt. Denn ich wiederum mir auch nicht alles lasse gefallen, was andere schreiben. Ein jeglicher sehe zu, wie er auf den Grund baue, Gold oder Holz, Silber oder Heu, Edelsteine oder Stroh. Des HErrn Tag wird es offenbaren [1 Cor. 3, 12. 13.].

\section{[Auslegung der letzten Worte Davids, 2 Sam. 23, 1–7.]}

\begin{schriftzitat}
V. 1. Dies sind die letzten Worte Davids.
\end{schriftzitat}

1. Solche letzten Worte meint er, darauf er sterben und hinfahren will, wie man spricht: Dabei ich es lassen, und ewig bleiben will. Denn es sind nicht seines Lebens oder Regiments letzte Worte, sondern sein Testament, und das wir Deutschen heißen: Seelrecht, darauf einer sterben will, daß es nach seinem Tode also geschehen und bleiben soll. Ultima voluntas heißt es bei den Juristen; da kann einer lange Zeit nach leben, viel reden, thun und leiden, bleibt gleichwohl sein Testament und letzter Wille immer fest. Also sind dies auch Davids letzte Worte, das ist, seiner Seele Testament, ob er darnach wohl viel geredet, gethan und gelitten hat; wie im Text folgt, als er das Volk zählen ließ, und darum gestraft ward, seinen Sohn zum Könige einsetzte, den Tempel auf dem Berge Moria zu bauen befahl, auch noch eine junge Metze, Abisag, eine\footnote{Die Worte: „Abisag, eine“ fehlen in der Wittenberger, in der Jenaer und bei Walch. Auch im Lateinischen fehlt der Name Abisag.} Sunamitin, zur Ehe nahm, die ihn wärmen sollte, weil der andern Weiber keine mehr zu ihm durften, nachdem sie von Absalom geschändet waren etc. [1 Kön. 1, 3.]

\begin{schriftzitat}
Es sprach David, der Sohn Isai
\end{schriftzitat}

2. Wie demüthiglich fähret er an; rühmt nicht seine Beschneidung noch Heiligkeit, auch nicht sein Königreich, sondern schlecht: „Der Sohn Isai“; schämt sich nicht seines geringen Geschlechts, daß er ein Schäferhirte gewesen ist; ja, das viel mehr ist, er bekennt seine Geburt, da er voll Sünden, und des Todes werth, hergekommen
\quellseite{950}
ist, wie alle Menschen. Denn er will von andern Sachen reden, die so hoch sind, da kein Adel noch Heiligkeit nütze zu ist, und kein Elend, weder Sünde noch Tod, schaden soll.

\begin{schriftzitat}
Es sprach der Mann, der versichert ist von dem Messia des GOttes Jakob, lieblich mit Psalmen Israel.
\end{schriftzitat}

3. Da fährt er heraus, und rühmt sich über alle Maße hoch, doch mit Wahrheit, ohne allen Hochmuth. Hier ist David ein anderer Mann denn Isai Sohn. Das hat er nicht von seiner Geburt geerbt, noch von seinem Vater gelernt, noch durch seine königliche Gewalt oder Weisheit erlangt. Von oben herab ist es ihm gegeben, ohne allen seinen Verdienst, deß ist er fröhlich, lobt und dankt so herzlich. Was ist es denn, das er rühmt? Erstlich ist es (spricht er), daß ich der Mann bin, dem GOtt verheißen hat den Messiam des GOttes Jakob; daß derselbe von mir, von meinem Blut, Stamm und Hause kommen wird, deß bin ich sicher und gewiß. Nicht allein daher, daß mir es GOtt zugesagt hat, der in seinen Worten gewiß und fest ist, und mir gewißlich nicht lügen wird, sondern auch ich dasselbe festiglich glaube, sicher und unbeweglich darauf stehe, als der ich in solchem Glauben nicht fehlen kann, und mich tröstlich verlasse auf GOttes Wort mit aller Zuversicht. Darum ich fröhlich bin, will gerne leben und sterben, wie und wann GOtt will. Ich weiß, wo ich oder meine Seele bleiben wird, und wo ich sie lasse. Sie soll mir nicht in der Irre oder in Zweifel schweben, noch übel von hinnen fahren. Ich habe gewisse Zusage GOttes von seinem Messia, so habe ich auch deß einen festen, gewissen Glauben.

4. Das Wort, Hebräisch \hebraeisch{אמן}, ist nicht wohl zu geben mit einem Worte. Constitutum est, sagt St. Hieronymus, ist nahe genug dabei. Stabilitus, certificatus, firmatus, gefestiget, wollte ich gerne sagen. Aber ich bin neuen Wörtern gram, so lautet es auch hier nicht wohl, gefestiget vom Messia etc. Ich halte, die Epistel Hebr. 11, 1. sehe hieher auf das Wort \hebraeisch{אמן}, da sie spricht: fides est substantia; graece ὑπόστασις, welches mir verdeutscht haben: „Der Glaube ist eine gewisse Zuversicht.“ Anders kann man es einem Deutschen nicht sagen, so er es verstehen soll. Denn der Glaube ist, und soll auch sein, ein Standfest des Herzens, der nicht wankt, wackelt, bebt, zappelt noch zweifelt, sondern fest steht, und keiner Sache gewiß ist. Desselben Wortes Gleichen hat man auch in dem Spruch Jes. 40, 8.: „GOttes Wort bleibet ewiglich.“ „Bleibet“, das ist, es hält fest, ist gewiß, weicht nicht, zuckt nicht, sinkt nicht, fällt nicht, läßt nicht fehlen. Wo nun dieses Wort in das Herz kommt mit rechtem Glauben, da macht es das Herz ihm gleich, auch fest, gewiß und sicher, daß es so steif, aufrecht und hart wird wider alle Anfechtung, Teufel, Tod, und wie es heißen mag, daß es trotzig und hochmüthiglich alles verachtet und spottet, was zweifeln, zagen, böse und zornig sein will, denn es weiß, daß ihm GOttes Wort nicht lügen kann. Solches ist ein \hebraeisch{אמן}, stabilitus, substantiatus, constantius, substantificatus,\footnote{In der Wittenberger, in der Jenaer und bei Walch: stantificatus; im Lateinischen fehlt dies Wort.} hypostaticus, certus passive. sicut verbum Domini certum active; wie Paulus 2 Tim. 1, 12.: „Ich weiß und bin gewiß“ etc., 2 Petr. 1, 10.: „Machet euren Beruf gewiß.“

5. Also ist David ein \hebraeisch{אמן}, der es gewiß hat in der Verheißung, und glaubt es dazu auch fest, daß Messia, den GOtt dem Patriarchen Jakob verheißen hat (1 Mos. 49, 10.: „Es soll das Scepter von Juda nicht wegkommen, bis daß der Schilo kommt“), gewißlich von seinem Blute kommen sollte; und ist hiermit die Verheißung vom Messia, Jakob geschehen, in David erneuet und klarer gegeben (wie folgen wird), daß man nun hinfort den Stamm Juda zu beiden Seiten gehen läßt, und allein auf Davids Haus sehen muß, als aus welchem, und keinem andern Hause im Stamme Juda, Messia kommen muß aufs allergewisseste. Doch, obwohl die zwei Stücke, Verheißung und Glaube, bei einander sein müssen; denn wo nicht Verheißung ist, da kann kein Glaube sein, und wo nicht Glaube ist, da wird die Verheißung zunichte; der Glaube aber nicht allezeit gleich fest, sondern zuweilen angefochten und schwach wird, die Verheißung aber, als das ewige GOttes Wort, immer für und für gleich fest und gewiß bleibt; darum heißt David vornehmlich daher \hebraeisch{אמן}, gefestiget, daß er die feste Verheißung hat, ob er gleich dieselbige ohne Glauben nicht halten oder fassen kann, er muß auch da sein. Das ist vom ersten gesagt.

6. Zum andern rühmt er sich: „Lieblich mit
\quellseite{951}
„Psalmen Israel“; das ist, er hat diese gewisse Verheißung vom Messia nicht bei sich, oder allein für sich behalten. Denn der Glaube ruht und feiert nicht, er fährt heraus, redet und predigt von solcher Verheißung und Gnade GOttes, daß andere Leute auch dazu kommen und der theilhaftig werden. Ja, vor großer Freude fährt er an, dichtet schöne, süße Psalmen, singt liebliche, lustige Lieder, damit zugleich GOtt fröhlich zu loben und zu danken, und auch die Menschen nützlich zu reizen und zu lehren. Also rühmt hier David auch, daß er habe viel schöner, süßer, lieblicher Psalmen von dem verheißenen Messia gemacht, die man zu Lobe GOtt in Israel singen sollte, und auch gesungen hat, darinnen zugleich auch treffliche Weissagung und hoher Verstand dem Volke Israel gepredigt und gegeben ist. Und da solch Dichten der Psalmen David anginge, und in Schwang brachte, wurden dadurch viel andere erleuchtet und zu Propheten erweckt, die auch dazu halfen und schöne Psalmen machten, als, die Kinder Korah, Heman, Asaph etc.

7. Denn er meint nicht allein die Lieblichkeit und Süßigkeit der Psalmen nach der Grammatica und Musica, da die Worte zierlich und künstlich gestellt sind, und der Gesang oder Ton süß und lieblich lautet, das da heißt schöner Tert und schöne Noten, sondern vielmehr nach der Theologia, nach dem geistlichen Verstande, da sind die Psalmen recht lieblich und süß, denn sie tröstlich allen betrübten, elenden Gewissen sind, die in der Sünden Angst, und Todes Marter und Furcht, und allerlei Noth und Jammer stecken. Solchen Herzen ist der Psalter, weil er den Messiam singt und predigt, ein süßer, tröstlicher, lieblicher Gesang, wenn man gleich die bloßen Worte, ohne Noten, daher liest oder sagt. Doch hilft die Musica oder Noten, als eine wunderliche Creatur und Gabe GOttes, sehr wohl dazu, sonderlich wo der Haufe mitsingt, und fein ernstlich zugeht. Denn so lesen wir von Propheten Elisa, 2 Kön. 3, 15., daß er durch das Psalterspiel (da man freilich Psalter aufgespielt hat, nach der Ordnung Davids) den Geist der Weissagung in sich erweckt; wie auch David mit seinem Psalterspiel oft den bösen Geist Sauls verjagt, oder doch hindert oder schwächt, lesen wir 1 Sam. 16, 23. Denn dem bösen Geiste ist nicht wohl dabei, wo man GOttes Wort im rechten Glauben singt oder predigt. Er ist ein Geist der Traurigkeit, und kann nicht bleiben, wo ein Herz geistlich (das ist, in GOtt und seinem Wort) fröhlich ist; davon auch St. Antonius sagt, daß geistliche Freude dem Teufel wehe thue.

8. Er nennt aber seine Psalmen „Israels Psalmen“, und will sie nicht sein eigen oder allein den Ruhm davon haben, sondern Israel soll sie bestätigen, und für die seinen urtheilen und erkennen. Denn es liegt daran, daß der Haufe GOttes oder GOttes Volk ein Wort oder Lied annehme, und für recht erkenne, weil der Geist GOttes in solchem Volke sein muß, der in seinem Volke will und soll geehrt sein. Also reden wir Christen von unsern Psalmisten. St. Ambrosius hat viel schöner hymnos ecclesiae gemacht, heißen Kirchengesänge darum, daß sie die Kirche angenommen hat und braucht, als hätte sie dieselben gemacht, und wären ihre Lieder. Daher spricht man nicht: Es singt Ambrosius, Gregorius, Prudentius, Sedulius, sondern, so singt die christliche Kirche. Denn es sind nun der Kirche Gesänge, die Ambrosius, Sedulius etc. mit der Kirche sungen, und die Kirche mit ihnen, und wenn sie sterben, so bleibt die Kirche, die immerfort ihre Lieder singt. Also will David seine Psalmen „Israels Psalmen“, das ist, der Kirche Psalmen heißen, welche denselben Geist hat, der sie durch David gemacht hat, und dieselben immerfort singt, auch nach Davids Tode. Er hat gefühlt im Geist, daß seine Psalmen für und für bleiben würden, so lange Israel oder GOttes Volk bleiben würde, das ist, bis an der Welt Ende; wie denn bisher geschehen ist, und geschehen wird. Darum sollen es „Israels Psalmen“ heißen.

\begin{schriftzitat}
V. 2. Der Geist des HErrn hat durch mich geredet. Und seine Rede ist durch meine Zunge geschehen.
\end{schriftzitat}

9. Hier will David mir zu wunderlich werden, und zu hoch fahren; GOtt gebe, daß ich es doch ein wenig erlangen möge; denn er fährt hier an von der hohen heiligen Dreifaltigkeit göttliches Wesens zu reden. Erstlich nennt er den Heiligen Geist; denn gibt er alles, was die Propheten weissagen. Und auf diesen und dergleichen Sprüche sieht St. Petrus 2. Ep. 1, 21.: „Es ist noch nie keine Weissagung aus menschlichem Willen hervorgebracht, sondern die heiligen
\quellseite{952}
Menschen Gottes haben geredet aus Eingebung des Heiligen Geistes.“ Daher singt man in dem Artickel des Glaubens\footnote{Das ist, in dem Nicänischen Glaubensbekenntniß.} von dem Heiligen Geist also: „Der durch die Propheten geredet hat.“ Also gibt man nun dem Heiligen Geiste die ganze heilige Schrift, und das äußerliche Wort und Sacrament, so unsere äußerlichen Ohren und Sinne rühren oder bewegen; denn auch unser HErr Christus selbst seine Worte dem Heiligen Geist gibt, da er spricht Luc. 4, 18. aus Jes. 61, 1.: „Der Geist des HErrn ist auf mir, darum hat er mich gesalbet“ etc. und Matth. 12, 18. aus Jes. 42, 1.: „Siehe, das ist mein Knecht, den ich erwählet habe. Ich will meinen Geist auf ihn legen“; und Luc. 1, 35.: „Der Heilige Geist überschattet Mariam“, rühret sie, nimmt ihr Blut und macht sie schwanger, daß der HErr heißt: „empfangen vom Heiligen Geist.“

10. Welch ein herrlicher, hochmüthiger Hochmuth ist das! Wer sich rühmen darf, daß der Geist des HErrn durch ihn redet, und seine Zunge des Heiligen Geistes Wort rede, der muß freilich seiner Sachen sehr gewiß sein. Das wird nicht sein David, Isai Sohn, in Sünden geboren, sondern der zum Propheten durch GOttes Verheißung erweckt ist. Sollte der nicht liebliche Psalmen machen, der solchen Meister hat, der ihn lehrt und durch ihn redet? Höre nun, wer Ohren hat zu hören! Meine Rede sind nicht meine Rede, sondern, wer mich höret, der höret GOtt, wer mich verachtet, der verachtet GOtt [Luc. 10, 16.]. Denn ich sehe, daß meiner Nachkommen viel werden meine Worte nicht hören, zu ihrem großen Schaden. Solchen Ruhm dürfen [weder] wir noch niemand führen, der nicht ein Prophet ist. Das mögen wir thun, sofern wir auch heilig, und den Heiligen Geist haben, daß wir Catechumenen und Schüler der Propheten uns rühmen, als die wir nachsagen und predigen, was wir von den Propheten und Aposteln gehört und gelernt, und auch gewiß sind, daß es die Propheten gelehrt haben. Das heißen in dem Alten Testament der „Propheten Kinder“, die nichts Eigenes noch Neues setzen, wie die Propheten thun; sondern lehren, daß sie von den Propheten haben, und sind „Israel“, wie David sagt, denn er die Psalmen macht.

\begin{schriftzitat}
V. 3. Es hat der GOtt Israel zu mir gesprochen, der HErr Israel hat geredet, der gerechte Herrscher unter den Menschen, der Herrscher in der Furcht GOttes.
\end{schriftzitat}

11. Nun haben wir drei Redner. Droben [V. 2.] sagt David, der Geist des HErrn habe durch seine Zunge geredet; damit ist die Person des Heiligen Geistes uns Christen klärlich angezeigt. Was Türken, Juden und andere Gottlosen glauben, achten wir nichts. So haben wir gehört, daß dem Heiligen Geist zugeneigt wird in der Schrift und in unserm Glauben die äußerliche Wirkung, da er durch die Propheten, Apostel und Kirchendiener mit uns leiblich redet, tauft und regiert. Darum sind diese Worte Davids auch des Heiligen Geistes, die er durch seine Zunge redet von zweien andern Rednern. Was redet er denn von denselben? Er redet erstlich von dem GOtt Israel, der zu David gesprochen, das ist, ihm verheißen habe. Wer nun GOtt, dieser Sprecher, sei, wissen wir Christen aus dem Evangelio Johannis, nämlich, es ist der Vater, der im Anfang sprach, 1 Mos. 1, 3.: „Es werde Licht“; und sein Wort ist die Person des Sohns, durch welch Wort alles gemacht ist, Joh. 1, 3. Denselben Sohn nennt der Geist durch David hier nur den „Hort Israel, und gerechten Herrscher unter den Menschen“. Der redet auch, das ist, der Heilige Geist führt den Hort Israel ein, daß er auch rede. Also reden alle drei Personen, und ist doch Ein Redner, Ein Verheißer, Eine Verheißung, wie es ein einiger GOtt ist.

12. Gleich aber, wie dem Heiligen Geist zugeeignet wird die äußerliche Wirkung GOttes an den Menschen, also ist es des Sohnes Eigenschaft, daß er Mensch worden, ein HErr und Richter über alle Menschen, und über alle Creatur gesetzt sei, wie Psalm 8, 5–7. singt: „Was ist der Mensch, daß du sein gedenkest, und des Menschen Sohn, daß du ihn heimsuchest. Du hast ihn eine kleine Zeit GOttes mangeln lassen, aber mit Ehren und Schmuck hast du ihn gekrönet. Du hast ihn zum Herrn gesetzt über deiner Hände Werk, alles hast du unter seine Füße gethan.“ Noch sind es nicht dreierlei Herrschaft, oder drei Herrscher, sondern Ein Herrscher und Eine Herrschaft, die der Vater dem Sohne, ja, dem Menschen und Menschenkinde gegeben hat, ohne Zweifel nicht also, daß er sich selbst und den Heiligen Geist davon ausgeschlossen
\quellseite{953}
habe. Dennoch ist es und heißt es des Menschen Herrschaft, die ihm GOtt gibt. Darum muß derselbe Mensch, der hier Herrscher heißt, rechter GOtt sein, weil er GOttes Reich besitzt, und also GOtt gleich in der eigenen Herrschaft ist. Denn GOtt läßt und gibt keinem andern seine Ehre oder sein eigen Reich; wie er spricht 2 Mos. 20, 3.: „Du sollst keine anderen Götter haben vor mir“, und Jes. 42, 8.: „Ich will meine Ehre keinem andern geben, noch meinen Ruhm den Götzen.“ Nun er aber diesem Menschen und des Menschen Sohn seine Ehre und Reich gibt, nämlich alles, was genannt ist, ihm unterwirft, wie er es selber unter sich hat: so muß derselbe Mensch kein anderer Gott noch Götze sein, sondern der rechte, natürliche Gott mit dem Vater und Heiligem Geist. Hiervon, so wir Zeit und Gnade haben, wollen wir hernach vergleichen mehr Sprüche, sonderlich im Psalter, handeln.

13. Jetzt lasset uns diese Worte Davids zuvor abarbeiten, darinnen er fein bekennt die zweien höchsten Artikel, daß in GOtt drei unterschiedliche Personen seien, und die eine, der Sohn, Mensch werden sollte, und das Reich und Ehre von dem Vater über alles empfangen, und der Heilige Geist solches in der Menschen Herzen durch den Glauben schreiben, der es zuvor auch verkündigt, hat durch den leiblichen Mund und Zunge der Propheten. Welches auch nichts denn eitel rechte göttlicher Majestät Werke sind; denn es ist nicht Menschen noch Engel Werk, solches zuvor verheißen, und hernach den Glauben in der Menschen Herz schaffen. Es ist GOttes Gabe (spricht St. Paulus Eph. 2, 8.) solcher Glaube, den der Heilige Geist wirkt und gibt.

14. Es ist auch nicht eines jeden, in der Schrift und Psalter die göttlichen drei Personen unterschiedlich also zu merken und zu lesen; denn wo ein fleischlicher Verstand über diese Worte kommt, der liest sie nach einander daher, wie sie da stehen: „Der Geist des HErrn hat durch mich geredet, der GOtt Israel hat zu mir gesprochen, der Hort Israel hat geredet, der gerechte Herrscher unter den Menschen“ etc., und denkt nicht anders, es sei alles von GOtt, als von Einer Person geredet, mit vielen übrigen Worten; oder fällt in die jüdische Blindheit, da sie David zu solchem gerechten Herrscher und Herrscher in der Furcht GOttes machen, und die Verheißung in (Gebot und Gesetze verkehren, daß, wer über Menschen will herrschen, der solle gerecht und gottesfürchtig sein, so doch David so andächtig und herzlich rühmt, es seien Worte der Verheißung, von dem Messia des GOttes Jakobs, und nicht Gebote über die weltlichen Herren.

15. Eben also würde es auch ihm gehen im andern Psalm, da auch die drei Personen unterschiedlich, als dreierlei Redner, reden. GOtt der Vater spricht [Ps. 2, 6.]: „Ich habe meinen König eingesetzt auf meinem heiligen Berge.“ Dieser König ist gewißlich eine andere Person von dem, der ihn zum Könige setzt. Nun folgt flugs darauf: „Ich will predigen von der Weise“; das lautet, als sei es noch der Vater, und so würde es lesen die Vernunft, so es doch der König, der Sohn ist, wie folgt [V. 7.]: „Der HErr hat zu mir gesagt: Du bist mein Sohn, heute habe ich dich gezeuget.“ Daß derselbe sei ein Mensch, ist gewiß; denn er soll predigen, und Messia sein, wie er sagt im andern Vers: „Sie toben wider den HErrn und seinen Messiam.“ Daß er aber GOtt sei, beweist, daß der Vater sagt: „Du bist mein Sohn, heute habe ich dich gezeuget“; wie wir Christen das wohl wissen. Item, daß er ihm gibt den ganzen Erdboden zu eigen, mit Heiden, und was darinnen ist, welches eben so viel ist, als GOttes eigenes Reich. Item, er befiehlt, V. 12.: „Sie sollen den Sohn küssen“ oder anbeten, ihm dienen mit Furcht, und schleußt, daß „selig sind alle, die ihm vertrauen“, welches allein GOtt zugehört. Ob sie ihm nicht alle gehorchen nach dem Evangelio, bricht seiner Herrschaft über alle Creaturen nichts ab. Wer nicht unter ihm sein will mit Gnaden, der muß unter ihm sein mit Ungnaden. Wer nicht mit ihm regieren will, der muß (wie seine Feinde) seiner Füße Schemel sein. Er ist Richter über Lebendige und Todte. Meinst du, ob der Türke, Papst, Juden und der ganze böse Haufe der Welt und Teufel, seiner Gnaden nicht wollen, sondern darwider toben, sie werden darum seiner Gewalt entgehen? Das werden sie wohl erfahren. Denn „GOtt spottet ihrer im Himmel, und wird in seinem Zorn mit ihnen reden“ [V. 4.]. Summa, er ist HErr, und bleibt HErr, so weit GOtt selbst HErr ist, denn er hat ihm die Herrschaft über alles gegeben; die Gewalt ist gewiß, und bleibet wohl. Wehe dem, der sie mit Gnaden nicht annimmt, der wird sie finden mit Zorn ewiglich.

\quellseite{954}
16. So haben wir nun hier abermal zwo unterschiedliche Personen, den Vater und den Sohn; so ist der Heilige Geist ohne das da, der solchen Psalm\footnote{Im Original: „solchen Psalmen“, das ist: solchen Psalm.} vom Vater und Sohn, mit ihren Worten eingeführt,\footnote{„mit ihren Worten eingeführt“ das heißt: in welchem Vater und Sohn redend eingeführt werden. — Cruciger bietet: hos psalmos, doch ist hier nur vom zweiten Psalm die Rede.} gemacht, und geredet hat. Also ist die unterschiedliche Dreifaltigkeit der Personen in Einem unzertrennlichen göttlichen Wesen, und daß der Sohn Mensch und Messias sei, bekannt, gleichwie es in den letzten Worten Davids bekannt ist. Ein fleischlich Herz läuft überhin, oder denkt, David habe es als ein frommer Mann gemacht, von sich selbst oder andern; wie die blinden Juden thun; aber David will es nicht leiden, daß man sollte ihm die Worte zuschreiben. Es sind lustige, liebliche „Psalmen Israel“ (spricht er), aber ich habe sie nicht gemacht, sondern „der Geist des HErrn hat durch mich geredet“. Und wie könnte auch Fleisch und Blut, Vernunft und menschliche Weisheit von solchen hohen, unbegreiflichen Sachen reden? Es ist eitel Narrheit und Aergerniß vor ihr.

17. Zu bestätigen aber diese Meinung Davids, daß er also, wie gesagt, geglaubt, und darauf gestorben ist, wollen wir die Worte vor uns nehmen, darauf er sich gründet, und aus welchen er solche lieblichen Psalmen gemacht hat, die stehen 2 Sam. 7, 12. ff. 1 Chron. 18, 10. ff. und lauten also:

\begin{quote}
[1 Chron. 18, 10–14.] Und der HErr verkündiget dir, daß der HErr dir will ein Haus bauen; und wenn deine Tage aus sind, daß du hingehest zu deinen Vätern, will ich deinen Samen nach dir erwecken, der deiner Söhne einer sein soll. Dem will ich sein Königreich bestätigen, der soll mir ein Haus bauen, und ich will seinen Stuhl bestätigen ewiglich. Ich will sein Vater sein, und er soll mein Sohn sein. Und ich will meine Barmherzigkeit nicht von ihm wenden, wie ich sie von dem gewandt habe, der vor dir war; sondern ich will ihn setzen in mein Haus und in mein Königreich ewiglich, daß sein Stuhl beständig sei ewiglich.
\end{quote}

18. Das erste Stück, da er sagt: „Der HErr verkündiget dir, daß der HErr dir will ein Haus bauen“, ist klarlich vom Hause Davids gesagt, daß seine Kinder sollen den Scepter Juda besitzen, bis auf Messia; wie davon genug gesagt ist in jenem Büchlein von den Juden.\footnote{Das ist: in der Schrift „Von den Juden und ihren Lügen“. Walch, St. Louiser Ausgabe, Bd. XX, 1912, S. 96 ff.} Und hier finden wir abermal die drei Personen in GOtt; erstlich des Heiligen Geistes, der durch den Propheten Nathan redet, wie wir droben [§ 9] gehört, daß die heilige Schrift durch den Heiligen Geist gesprochen ist, nach dem Spruch David: „Der Geist des HErrn hat durch mich geredet.“ Eben also redet er durch alle Propheten. Weiter führt der Heilige Geist ein die Person des Vaters in seiner Rede, da er spricht: „Der HErr verkündiget dir.“ Und flugs darauf die Person des Sohnes, da er spricht: „Daß dir der HErr ein Haus bauen will.“ Und ist doch ein einiger GOtt und HErr, der da redet durch Nathan, und David verkündigt, und sein Haus bauet; alle drei Ein Redner, Ein Verkündiger, Ein Bauherr. Ob solchen Unterschied der Personen nicht jedermanns Vernunft sieht in der Schrift, da liegt nichts an. Ich weiß wohl, wie sie hier und dergleichen Oerter mehr glossiren, die Nasenweisen, so den Heiligen Geist meistern.

19. Aber wo du in der Schrift findest, daß GOtt von GOtt, als wären es zwo Personen, redet, da magst du füglich drauf gründen, daß daselbst drei Personen in der Gottheit angezeigt werden. Also, hier an diesem Orte spricht der HErr, daß der HErr will David ein Haus bauen. Item, 1 Mos. 19, 24.: „Der HErr ließ regnen vom HErrn Feuer und Schwefel“ etc. Denn der Heilige Geist ist kein Narr noch Trunkenbold, der Ein Tüttel, gleichwie Ein Wort sollte vergeblich reden. Regnet' nun der HErr (das ist der Sohn) vom HErrn (das ist vom Vater) Feuer und Schwefel, so ist zugleich da der Heilige Geist, der solches durch Abraham, oder wer es gewesen ist, von den zweien HErren redet. Dennoch sind sie alle drei Ein HErr, Ein GOtt, der Feuer und Schwefel regnet. Solcher Exempel wollen wir hernach mehr hören.

20. Das andere Stück, da er sagt: „Wenn deine Tage aus sind, daß du hingehest zu deinen Vätern, will ich deinen Samen nach dir ermeden
\quellseite{955}
“ etc. Hier geht der rechte Text an von Messia. Denn solches kann von Salomo nicht gesagt sein, viel weniger von einem andern Sohne Davids; es muß der rechte einige Sohn Davids sein, Messia, der nach dem Scepter Juda kommen sollte. „Der soll mir ein Haus bauen“, spricht er, „und ich will sein Königreich bestätigen ewiglich.“ Dies Haus kann nicht sein der Tempel Salomonis. Denn hart davor [1 Chron. 18, 4. 5.] spricht er: „Du sollst mir nicht ein Haus bauen zur Wohnung. Habe ich doch in keinem Hause gewohnet, sint der Zeit ich die Kinder Israel aus Egypten geführet habe.“ Und Salomo selbst 1 Kön. 8, 27.: „Meinest du, daß GOtt auf Erden wohnen möge? Siehe, der Himmel und aller Himmel Himmel mögen dich nicht begreifen, wie sollte es denn thun dies Haus, das ich gebauet habe?“ Aber noch gewaltiger sagt solches Jesaia, Cap. 66, 1.: „So spricht der HErr: Der Himmel ist mein Stuhl, und die Erde meine Fußbank. Was ist es denn für ein Haus, das ihr mir bauen wollet? oder welches ist die Stätte, da ich ruhen soll?“

21. Hier verwirft GOtt mit ausgedrückten Worten der Juden tolle Andacht, die sich des Tempels hoch rühmten, als baueten sie GOtt ein Haus, und thäten GOtt damit großen Dienst, wurden drüber stolze, halsstarrige Prophetenmörder, so doch GOtt hier bekennt, er wolle den Tempel nicht ansehen, sondern einen demüthigen, betrübten Geist, der sich vor seinem Wort fürchtet [Jes. 66, 2.], ja, der soll sein Tempel und Altar sein; verwirft auch daselbst alle Opfer und Gottesdienste des Tempels, da er spricht, V. 3.: „Wer einen Ochsen schlachtet, ist eben als der einen Mann erschlüge. Wer ein Schaf opfert, ist als der einem Hunde den Hals bricht. Wer Speisopfer bringt, ist als der Säublut opfert. Wer des Weihrauchs gedenket, ist als der Unrecht“, das ist, Götzendienst, „lobet.“ Denn GOtt hatte den Tempel nicht darum lassen bauen, daß sie sein Wort stolziglich sollten verachten, und dafür viel Opfers treiben, damit sich selbst zu heiligen, sondern „daß sein Name“, nicht er selbst, „da wohnen sollte“ (wie die Schrift allenthalben sagt), das ist, sie sollten daselbst sein Wort hören, und ihn anrufen, dadurch er geehrt würde. So wollten sie den Ruhm und die Ehre davon haben, daß sie solchen Tempel hätten, und schlugen die Propheten um GOttes Worts willen zu Tode.

22. Darum muß dies Haus GOttes, so Messia, Davids und GOttes Sohn, bauen soll, viel ein anderes, größeres und herrlicher Haus sein. Denn, rechne du selbst, soll GOtt in diesem Hause wohnen, so muß es viel größer und herrlicher sein, neben Himmel und Erde, weil er so groß ist, daß der Himmel sein Stuhl, darauf er sitzt, und die Erde seine Fußbank ist. Was will hier noch Raum her gehören, da sein Haupt, Brust, und Arm wohnen mögen? Daß es wohl gesagt ist: Was wollet ihr mir für ein Haus von Stein und Holz bauen zur Wohnung, so mir Himmel und Erde viel, viel zu enge ist? Von diesem Hause lehrt uns nun die heilige Schrift, sonderlich das Neue Testament, und heißt die heilige christliche Kirche, so weit die Welt ist; nicht allein das, sondern auch ein ewig Haus ist, das ewiglich bleibt und lebt, da GOtt ewiglich innen wohnt, regiert, und haushält. Das will ein Häuslein und Tempel werden!

23. Nun laßt uns den Zimmermann oder Baumeister dieses Hauses ansehen. Er soll ein Mensch und Davids Sohn sein, wie der Text da steht, „einer aus deinen Söhnen“; und soll doch ein Haus GOttes bauen, das besser und herrlicher sei, denn Himmel und Erde, dazu ewiglich bestehe. Wo will er die Kunst und Macht hernehmen? Das kann nicht weder Menschen noch Engel Kunst oder Macht sein. Denn Engel können nicht Himmel noch Erde schaffen, ja, nicht die geringste Creatur, viel weniger der Mensch. Darum muß dieses Hauses Zimmermann rechter wahrhaftiger GOtt sein, der die rechte Macht göttlicher Natur habe, zu schaffen Himmel und Erde, und noch viel Besseres, das ist, er muß ein allmächtiger GOtt sein, und ist doch nicht die Person, die von ihm sagt: „Ich will sein Vater, und er soll mein Sohn sein, und er soll mir ein Haus bauen.“ Hier sind die Personen gewaltiglich und klärlich unterschieden, als Vater und Sohn, und als Baumeister und Hausherr. Noch können es nicht zween Götter sein, oder der Sohn ein anderer und fremder Gott sein. Denn das erste Gebot leidet es nicht, da es sagt: „Du sollst nicht andere Götter neben mir haben“; und 5 Mos. 6, 4.: „Höre, Israel, der HErr, unser GOtt, ist ein einiger HErr“ oder GOtt.

24. So haben wir droben [§ 9. 16] gehört: Wo die Schrift von den zwo Personen des Vaters und Sohnes redet, da ist der Heilige Geist, die dritte Person, auch bei, der solches durch die
\quellseite{956}
Propheten redet, daß also an diesem Orte einem gläubigen Herzen gar gründlich und gewaltig beweiset und bezeuget\footnote{„und bezeuget“ fehlt in der Erlanger.} ist, daß GOtt, der allmächtige Schöpfer Himmels und der Erde, ist gewißlich ein einiger rechter GOtt, und außer ihm kein anderer Gott\footnote{„Gott“ fehlt in der Erlanger.} sein kann, und doch drei unterschiedliche Personen, der Vater, Sohn und Heiliger Geist ist; also, daß allein der Sohn sei Mensch und Davids Sohn worden. Und ist ohne Zweifel dem Volke Israel darum so hart geboten, daß sie ja nicht mehr denn Einen Gott sollten anbeten, auf daß sie sich nicht ärgern dürften, wenn Messia käme, und für GOtt gepredigt und geglaubt würde, gerade als wollte er wider Mosen mehr denn Einen, oder fremden Gott lehren, sondern die Ohren und Herzen stille halten, sich lehren lassen, wie das erste Gebot von Einem GOtt recht und gründlich zu verstehen wäre. Folgt:

\begin{quote}
[1 Chron. 18, 14.] Und ich will ihn setzen in mein Haus und in mein Königreich ewiglich.
\end{quote}

25. Was ist das? Hier lasset uns hören. Das Haus soll ewig sein und bleiben, wie droben [§ 22] auch gesagt ist, darum muß der Meister auch ewig [sein], und eine ewige göttliche Macht haben. Hier spricht er weiter zu David: Das Haus, das mir mein und dein Sohn bauen soll, da soll er auch selbst Hausherr innen sein, so wohl als ich. Er soll mir gleich sein in einerlei Hause. Ich will ihn hinein setzen, daß er es soll haben, gleichwie ich selbst. Nun haben wir gehört, daß dies Haus GOttes größer, besser, und herrlicher sei denn Himmel und Erde. Ist nun Davids Sohn, Messia, Meister und Hausherr dieses Hauses, so ist er gewißlich auch Meister und Hausherr Himmels und der Erde, und viel mehr und besser. Denn wer dieses Hauses HErr ist, wie GOtt selbst, der ist freilich viel mehr,\footnote{Cruciger: multo pluris ipsa domo. Die Wittenberger und die Jenaer haben es ebenso verstanden und deshalb ein Komma nach „viel mehr“. Die Erlanger bietet: „vielmehr auch“ etc.} auch Herr Himmels und der Erde, wie GOtt selbst; das kann aber niemand sein, denn allein der einige GOtt, Schöpfer Himmels und der Erde. Daraus folgt, daß Messia, Davids natürlicher Sohn, muß rechter wahrhaftiger GOtt, und kein fremder Gott sein. Denn, wie gesagt, GOtt läßt keinen fremden Gott in seinem eigenen Hause Herr sein, er muß und will seine Ehre und Gewalt selbst haben, und keinem andern lassen. Solches sollte ja deutlich genug so viel heißen, daß Messia, Davids Sohn, sei ein Herr und König in GOttes eigenem Reiche, oder GOtt gleich; denn gewiß ist es, daß GOtt daselbst vom Messia redet.

26. Will aber jemand judenzen, und diese Worte „mein Haus“, „mein Reich“, deuten auf den Tempel zu Jerusalem, und auf das Volk Israel, der mag es ohne mich für sich selbst thun, denn ich sehr wohl weiß, daß GOtt den Tempel sein Haus heißt, Jes. 56, 7.: „Mein Haus ist ein Bethaus“, wie Christus solches einführt Matth. 21, 13. Luc. 19, 46. etc., und Israel sein Königreich heißt, 2 Mos. 19, 6.: „Ihr sollt mein priesterlich Königreich sein“ (es steht aber dabei, „so ihr meiner Stimme gehorchet, und meinen Bund haltet“). Auch weiß ich, daß GOtt Ezech. 18, 4. spricht: „Des Vaters Seele ist mein sowohl, als des Sohnes Seele.“ Ja, ich weiß, daß der Wein, den ich trinke, und das Brod, das ich esse, heißt auch sein Brod und Wein; und was ist im Himmel und Erden, das nicht sein ist? wie er spricht Jes. 66, 1.: „Der Himmel ist mein Stuhl, die Erde meine Fußbank, was wollt ihr mir für ein Haus bauen? Hat nicht meine Hand das alles gemacht, und ist alles bereit da?“ das ist, hätte ich nicht zuvor Himmel und Erde geschaffen, wo wolltet ihr Kalk, Steine, Holz, Eisen, und was zum Bau gehört, nehmen? Ist es nicht zuvor alles mein? Wo habt ihr es her? Was habt ihr daran gearbeitet? Ja, wer und weß seid ihr selber? Bin ich nicht euer Schöpfer? Also bekennt auch David in seinem schönen Lobspruch, da er von dem Vorrath zum Tempel zu bauen sagt, 1 Chron. 30, 11.: „Es ist alles dein, HErr, und wir opfern dir, was wir von deinen Händen empfangen haben.“

27. Eben also redet er auch vom Opfer dieses Tempels, Pf. 50, 13.: Was wollt ihr mir opfern? „Meinest du, ich müsse Rindfleisch fressen, und Bocksblut trinken?“ Woher hast du deine Opfer, Rind, Schafe, Ziegen? Ist es nicht zuvor alles mein? Weß sind alle die Rinder, Schafe und alles, was auf\footnote{Wittenberger und Jenaer: an.} der Weide geht auf dem ganzen Erdboden? Habe ich sie
\quellseite{957}
nicht geschaffen, ohne alle dein Zuthun? Wo wolltest du sie nehmen mir zum Opfer, wenn ich es nicht dir zuvor gegeben hätte? Das ist so viel gesagt: Deiner Opfer darf ich nichts, und sollst nicht denken, daß du mir damit dienst, als müßte ich sie haben, und könnte ihrer nicht entbehren; sondern das ist die Meinung, daß du sollst erkennen und bekennen durch solch Opfer, wie du alles von mir hast, was du bist und hast, und also mich [als] deinen GOtt und Schöpfer ehren, loben und preisen. Ja, um der Ursache willen habe ich es eine Beilage meiner Opfer lassen nennen. Wo die Ursache aus ist, da ist das Opfer nichts mehr.

28. Gleichwie ich den Tempel habe mein Haus genannt, nicht, daß ich darinne wohnen müßte oder dürfte sein, daß ihr daran mir einen Dienst thätet, sondern um euretwillen heiße ich es mein Haus, daß ihr drinnen beten, mich loben und anrufen sollet; denn es soll nicht mein Wohnhaus, sondern euer Bethaus sein, ja, ein Bethaus soll es heißen. Nun darf es noch kann es um meinetwillen kein Bethaus heißen; denn ich habe niemand anzubeten noch anzurufen, als ein GOtt, der ich niemandes bedarf. Wo nun mein meines Hauses anders denn als eines Bethhauses braucht, so ist es eine Mördergrube, und nicht mein Haus; wie denn die thun, so da meinen, sie thun mir großen Dienst mit dem Werke, daß sie mir ein Haus bauen, wollen des gerühmt sein, und die Ehre selbst haben von solchem Werk und Gebäu, als verdieneten sie damit große Gnade bei mir. Da muß denn solch Haus zu Grunde, zunichte und zerstört werden als die allerschändlichste Mördergrube, die nicht mehr mein Haus, sondern des Teufels eigene Hölle ist.

29. Wer nun solchem Verstand nach will diese Worte „mein Haus“, „mein Reich“, wie [§ 26] gesagt, vom Tempel und Volk Israel an diesem Orte verstehen, der muß förter auch das auf sich nehmen, daß er beweise mit gutem, starkem Grunde, wie der Tempel zu Jerusalem und das Volk Israel im Lande Canaan noch jetzt und bis daher, sint der Zeit Davids, immer für und für geblieben sei, weil der Text hier klärlich ausspricht: Davids Haus solle ewiglich bleiben, und Davids Sohn, Messia, solle ewiglich in GOttes Hause und Reiche sitzen. Wir Christen müssen bekennen, daß wir solches nicht können beweisen, als die wir wissen, daß GOttes Haus, der Tempel zu Jerusalem, bei 1500 Jahr in der Asche liegt, Davids Haus und Königreich und das Volk Israel auch bei 1500 Jahren nichts gewesen, noch Regiment oder Reich im Lande Canaan gehabt hat, und müssen bei unsern vorigen Verstande bleiben, daß die Worte „mein Haus“ und „mein Reich“ müssen heißen das ewige Reich GOttes, da er ewiglich inne wohnen wolle und regieren, welches sein und Davids Sohn, Messia, bauen sollte durch seine allmächtige göttliche Macht und Weisheit. Deß lässet uns David selbst hören, wie er diese Worte verstanden hat. So steht 1 Chron. 18, 16.:

Da Nathan alle diese Worte (wie droben angezeiget) mit David geredt hatte, kam der König David, und setzte sich vor den HErrn, und sprach: Was bin ich, HErr GOtt? Und was ist mein Haus, daß du mich bis daher bringest?

30. Hier zeigt David an, daß er die Worte wohl verstanden hat, da GOtt durch Nathan ihm verhieß: „Ich will sein Vater sein, und er soll mein Sohn sein; ich will ihn setzen in mein Haus, und in mein Königreich ewiglich.“ Darum spricht er: „Was bin ich? Was ist mein Haus, daß du mich bis dahin bringest?“ Es ist ja zu hoch und zu herrlich Ding, das du mir verheißest, daß mein Haus, ich, mein Sohn, solle dahin kommen, daß er in deinem eigenen\footnote{„eigenen“ fehlt in der Erlanger. Cruciger: in illo tuo solius sempiterno regno.} ewigen Reiche sitzen, HErr und König sein soll? HErr GOtt, wohin bringst du mich? Er kann es vor großem Wundern nicht sagen, und nennt es „bis dahin“. Wohin, wohin, mein lieber GOtt? Soll ich (das ist, mein Fleisch und Blut) dir gleich sitzen in deinem ewigen Reiche, so wird mein Fleisch und Blut, mein Sohn und dein Sohn, rechter wahrhaftiger Gott sein müssen, der dir gleich sitzt. Ach GOtt, wo bringst du mich hin? Folgt:

\begin{schriftzitat}
V. 17. Du hast angesehen mich, als in der Gestalt eines Menschen, der in der Höhe GOtt, der HErr, ist.
\end{schriftzitat}

31. Fast alle anderen Hebräischen dolmetschen hier viel anders; etliche aber, als Bernhard Ziegler,\footnote{M. Bernhard Ziegler, hebräischer Professor zu Leipzig. Walch, St. Louiser Ausg., Bd. XX, 2106, § 189. Am 10. October 1543 wurde er zum Doctor promovirt. De Wette, Bd. VI, S. 349.} geben mir Kundschaft, daß es kann
\quellseite{958}
und solle nach der Grammatica also verdolmetscht werden, wie ich es jetzt verdeutscht habe. Hiermit bekennt nun David klarlich, daß sein Sohn, Messia, solle gewißlich ein rechter Mensch sein, in aller Gestalt, Weise und Maße wie ein anderer Mensch, Phil. 2, 7., und doch überwärts und in der Höhe, da keine Menschen-Weise, sondern allein GOtt ist und regiert, soll er GOtt, der HErr, sein. Das ist, sage ich, klarlich die Meinung Davids, mit dürren Worten herausgesagt. Darum er droben [V. 14.] sagt: Wohin, wohin, lieber GOtt, bringst du mich? Und hier: Wofür siehst du doch mich unwürdigen Menschen an, daß mein Sohn soll in deinem ewigen Reiche König sein? Er versteht es wohl, daß in GOttes ewigem Reiche König sein niemand gebühren mag, denn der ein rechter GOtt ist. Weil nun der Sohn Davids Mensch ist, und eine andere Person denn der Vater, so ihn in sein Königreich setzt, und doch nicht können zween, oder mehr denn ein einiger GOtt sein; so schließt hiermit David, daß sein Sohn, Messia, muß rechter, natürlicher GOtt sein, und doch kein anderer GOtt denn der Vater, sondern eine andere Person in derselbigen einigen unzertrennten Gottheit, und der Heilige Geist, der solches durch Nathan und David, als rechter GOtt, vom Vater und Sohn redet, die dritte Person in derselbigen einigen Gottheit ist.

32. Das ist die Lehre und Glaube des Neuen Testaments, nämlich, daß Jesus Christus von Nazareth, Davids und der Jungfrauen Marien Sohn, rechter Mensch sei, GOttes natürlicher, ewiger Sohn, mit dem Vater und Heiligen Geist ein einiger GOtt, und drei unterschiedliche Personen. Weil nun Davids Worte an diesem Orte solchen Verstand gerne geben, nach aller Art hebräischer Sprache, sollen wir Christen keinen andern Verstand drinnen suchen noch achten, sondern diesen [für] den einigen allein rechten Verstand, alle andere Deutung für menschlichen nichtigen Dünkel, halten. Das Neue Testament kann nicht fehlen, also daß Alte Testament auch nicht, wo es sich reimt, und dem Neuen ähnlich ist.

33. Möchtest du hier fragen: Gehen die Worte Davids und Nathans so klarlich den Artikel von der Gottheit Christi, wie geht es zu, daß weder die heiligen Väter, noch kein anderer Lehrer, solches gesehen oder jemals gerührt haben, und ihr neuen, jungen Hebräisten habt es nun erst ersehen? warum ersehen es die Rabbinen der Juden nicht? Antwort: Die hebräische Sprache ist nach der Apostel Zeit wenig und übel bekannt gewest. [Es] haben sich die lieben Väter und Lehrer am Neuen Testament genügen lassen, darinnen sie solches und alles reichlich und überflüssig haben. Die Propheten aber und Apostel haben es wohl gesehen; wie wir hernach hören werden. Daß aber die Rabbinen nicht sehen, da geht es ganz recht zu.\footnote{Cruciger: id sane recte et merito fit = daran geschiehet ihnen freilich ganz recht.} Denn wer blind ist, der soll nichts sehen, wie Jesaias, Cap. 6, 10., von ihnen sagt: „Mit sehenden Augen werdet ihr blind sein.“ Wer von denen lernen muß, der wird auch blind, gewißlich. Zwar wir selbst würden es auch nicht sehen, wo wir nicht, durch das Neue Testament erleuchtet, dem Alten recht unter die Augen sehen könnten. Denn ohne das Neue Testament ist das Alte verdeckt, 2 Cor. 4, 3. 4.

34. Siehe unsere Zeit an, da wir predigen von der Gnade Christi wider die vermeinten eigenen Werke und Heiligkeit, wie viel deren sind, die es sehen, oder mit Ernst annehmen? Woran fehlt es? Es wird ja so helle gepredigt, gelehrt, gelesen, geschrieben, gesungen, gemalt, und auf alle Weise getrieben, daß es schier sollten Holz und Stein verstehen, wenn sie ein wenig Vernunft hätten: noch sehen es nicht Pabst, Könige, Fürsten, Bischöfe, Gelehrte, Herren, Adel, Bürger, Bauern, sondern gehen vorüber, mit sehenden Augen blind, mit hörenden Ohren taub; denn ihr Herz ist nicht daheime, und steht anderswo hin. Also haben die Propheten zu ihrer Zeit von Christo auch klarlich genug geweissagt, daß er sollte GOtt und HErr über alles sein, wie David hier thut. Aber wenig haben es geglaubt oder geachtet, die andern sind blind und taub dagegen gewesen, haben ihrem Herzen nachgegangen, und ihrem Dünkel gefolgt. Es heißt mysterium, Geheimniß, und bleibt mysterium. Wer es versteht und mit Ernst\footnote{So die Wittenberger und die Jenaer. Erlanger: ernst.} meint, der dankt GOtt, und kehre sich nicht an den andern großen Haufen der Verächter.

35. Meinst du nicht, daß Jesaias diesen Text Davids mit Fleiß gelesen habe, da er spricht, Cap. 9, 6. 7.: „Ein Kind ist uns geboren, ein Sohn ist uns gegeben, deß Herrschaft ist auf
\quellseite{959}
jener Schulter, und sein Name heißt Wunderbar, Rath, GOtt, Held, ewig Vater, Friedefürst, daß seine Herrschaft groß werde, und des Friedens kein Ende, auf dem Stuhl Davids und seinem Königreich, daß er es zurichte und stärke mit Gericht und Gerechtigkeit, von nun an bis in Ewigkeit“? Hier nimmt Jesaias dem Nathan das Wort aus dem Munde, da er von Messia weissagt, daß er ein ewiger König und Vater sein soll in GOttes Reiche, und er nennt ihn auch GOtt: Denn das Wort El heißt wohl, nach dem Buchstaben, Kraft; aber wenn es ein Eigenname ist (wie hier), so heißt es GOtt, durch die ganze Schrift, der allein Kraft hat; wie das beide, Juden und Hebräiten, bekennen müssen.

36. So stimmt nun Jesaia mit David und dem Neuen Testament, daß Christus sei ein ewiger König, und rechter GOtt; darum muß sein Reich ein göttliches, ewiges Reich sein auf dem Stuhle Davids etc. Denn er hat sonderlich das Wort „ewiges Reich“ bewogen, da GOtt spricht durch Nathan zu David: „Ich will deinen Sohn in mein Reich setzen ewiglich“, und fühlt wohl, daß solches geredet ist (wie David spricht) von einem Menschen, der droben in der Höhe El, das ist, GOtt sein müsse. Denn ewiges Reich GOttes besitzen, und König drinnen sein, kann nicht eines schlechten Menschen sein, auch nicht ein vergänglich, zeitlich, irdisch Reich sein, welches ein Ende hat, und der König sterben muß mit seinen Kindern nach ihm; aber hier soll der Sohn Davids ein ewiger König im ewigen Reiche GOttes sein, und, wie Jesaias mitstimmt, des Friedens kein Ende, und er, der Sohn Davids, das Kind, so uns geboren und gegeben ist, soll ein ewiger Vater und Friedefürst sein, von nun an bis in Ewigkeit. Darum muß er GOtt oder El sein, der durch seine göttliche Macht solchen ewigen Frieden geben und erhalten könne.

37. Solche Ewigkeit des Reichs Messia zeigt Jesaias an mehr Oertern, als Cap. 51, 4. 5.: „Merke auf mich,\footnote{„mich“ fehlt in der Wittenberger, der Jenaer und der Erlanger. Im Lateinischen ist es ausgedrückt.} mein Volk; höret mich, meine Leute. Denn von mir soll ein Gesetz ausgehen, und mein Recht will ich zum Licht der Völker stellen. Meine Gerechtigkeit ist nahe, und mein Heil ist ausgegangen.“ Und bald hernach, V. 6.: „Mein Heil soll ewig bleiben, und meine Gerechtigkeit soll kein Ende haben.“ Dies ist die ewige Gerechtigkeit, davon auch Daniel sagt, Cap. 9, 24.: „Siebenzig Wochen sind bestimmt, daß die ewige Gerechtigkeit komme“, und ist Messias, wie es alle alten Hebräer verstanden haben. „Ewige Gerechtigkeit“ aber und Heil kann kein schlechter Mensch noch Engel sein, sondern muß GOtt selber sein; und doch Davids Sohn, natürlicher Mensch, und eine andere Person von dem, der von ihm redet, und ihn nennt „mein Heil, meine Gerechtigkeit“. Die dritte Person ist der Heilige Geist, der solches redet von den beiden. Also nennt ihn das Neue Testament auch, 1 Cor. 1, 30.: „Jesus Christus ist uns von GOtt worden eine Weisheit, Gerechtigkeit, Heiligung und Erlösung.“ Das reimt sich mit\footnote{„mit“ fehlt in der Erlanger.} Jesaia, und Jesaia mit Paulo.

38. Item, Jesaia 60, 19. 20.: „Die Sonne soll dir nicht mehr des Tages scheinen, und der Glanz des Monden soll dir nicht leuchten; sondern der HErr wird dir ein ewiges Licht sein, und dein GOtt wird deine Herrlichkeit sein. Deine Sonne wird nicht mehr untergehen, und dein Mond wird nicht mehr den Schein verlieren, denn der HErr wird dir ein ewiges Licht sein, und die Tage deines Leides sollen ein Ende haben.“ Hier steht klar, daß unser ewiges Licht solle der HErr und unser GOtt selbst sein, und redet ein HErr von dem andern; ja, durchs ganze Capitel redet nicht Jesaias, sondern der HErr, und spricht allhier: „Der HErr wird dein ewiges Licht sein.“ Wer ist der HErr, der solches spricht? Ohne Zweifel GOtt der Vater. Wer ist der HErr, von dem er spricht: „Der HErr wird dein ewiges Licht sein“? Ohne Zweifel GOtt der Sohn, Jesus Christus. Denn hier steht der große Name GOttes, Jehovah, den wir in unserer Bibel mit diesen großen Buchstaben schreiben, HERR, zum Unterschied der andern Namen. Wer ist es denn, der solches durch die Zunge Jesaia redet? Ohne Zweifel GOtt der Heilige Geist, der durch die Propheten redet, und führt ein die Person des Vaters, redend von dem ewigen Licht, das ist, von seinem Sohne, Jesu von Nazareth, Davids und Mariä Sohn.

39. Solch ewig Licht, ja, solcher Herr, kann
\quellseite{960}
auch kein schlechter Engel noch Mensch sein. Und hierin stimmt Jesaias Weissagung mit dem Neuen Testament, da unser HErr Jesus Christus sich selbst oft ein Licht nennt, Joh. 1, 4. 5.: „Das Leben war das Licht der Menschen, und das Licht leuchtete in die Finsterniß, aber die Finsterniß begriffen es nicht.“ Weil nun solches sich mit dem Neuen Testament reimt, soll man Jesaiä Weissagung fröhlich, und nicht anders denn von Jesu Christo verstehen, der uns nicht ein Reich unter dieser Sonne und Mond bereitet hat, das da vergeht, sondern will selber unser ewiges Licht, Sonne und Mond, Leben und Heil sein, wie er droben sagt, Cap. 51, 6.: „Hebet eure Augen auf gen Himmel, und sehet herunter auf Erden! Der Himmel wird wie ein Rauch vergehen, und die Erde wie ein Kleid veralten, und die drauf wohnen werden dahin sterben wie ein Nichts; aber mein Heil soll ewig bleiben, und meine Gerechtigkeit soll kein Ende haben.“

40. Lieber, sage mir, kann auch dieser und desgleichen Sprüche neben sich leiden der Juden tollen Verstand von ihrem Messia, der ein sterblicher, weltlicher König sein soll zu Jerusalem, auf Erden? so GOtt hier gegen einander hält seinen Messiam gegen Himmel und Erde, und spricht: „[Der] Himmel soll wie ein Rauch (das wird ohne Feuer nicht geschehen, 2 Petr. 3, 12.) vergehen, die Erde wie ein alt Kleid verweisen, die Menschen wie ein Nichts dahin sterben. Aber sein Heil, das nahe ist (spricht er), seine Gerechtigkeit“, die ausgegangen, „soll ewig bleiben“, und ein ewiges Licht sein; denn er ist der HErr selbst, und dein GOtt. Hier siehst du, ob Jesaias die Worte Nathans, 1 Chron. 18, 13., verstanden habe, da er GOtt einführt: „Ich will sein Vater sein, und er soll mein Sohn sein, ich will ihn setzen in mein Königreich ewiglich.“ Und Davids Worte, 1 Chron. 18, 17., da er spricht: „Du hast mich angesehen, als in der Gestalt eines Menschen, der droben in der Höhe GOtt der HErr ist“; und 2 Sam. 7, 19.: „Das ist eine Weise eines Menschen, der der HErr HErr ist“, ubi Latinus\footnote{Latinus, das ist, die Vulgata.} habet: Ista est lex Adam, Domine Deus; et nihil significat [wo der Lateiner hat: „Dies ist das Gesetz Adams, o HErr GOtt“, und keinen Verstand gibt].

41. Lasset uns Daniel auch sehen, der spricht Cap. 7, 13. 14.: „Ich sahe im Gesichte des Nachts, und siehe, es kam einer in Wolken des Himmels, wie eines Menschen Sohn, bis zu dem Alten, und ward vor denselben gebracht, der gab ihm Gewalt, Ehre und Reich, daß ihm alle Völker, Leute und Sprachen dienen sollten. Seine Gewalt ist ewig, die nicht vergehet, und sein Königreich hat kein Ende.“ Diesen Spruch können und verstehen die Christen wohl, doch wollen wir auch sehen, wie er sich mit dem Neuen Testament reimt. Ein Menschenkind sieht er in den Wolken; ohne Zweifel, daß sein Reich nicht irdisch, vergänglich, noch zeitlich, sondern himmlisch und ewig sein soll, wie er spricht: „Der Alte“, das ist, GOtt der Vater, „gab ihm Gewalt über alles, und seine Gewalt soll ewig sein, und kein Ende haben.“ Diese Ewigkeit oder ewiges Reich kann keiner schlechten Creatur, weder Engel noch Menschen, gegeben werden, denn es ist göttliche Gewalt, und GOttes eigene Gewalt. Was würde oder könnte GOtt behalten oder haben, wenn er die ewige Gewalt und ewiges Reich von sich gäbe? Nichts überall behielte er, und würde selbst zunicht, weil ein anderer da ist, der die ewige Gewalt hat. Es kann ja über und außer der ewigen Gewalt nichts sein. Ewige Gewalt begreift es alles, und läßt nichts Anderes oder Größeres über sich, oder außer sich sein; es muß GOtt selber sein, und nichts Anderes.

42. Hier gibt der Text Danielis auch gewaltiglich den Artikel von der Gottheit in drei Personen, und von der Menschheit des Sohnes; denn es muß eine andere Person sein, die da gibt, und eine andere, die es empfähet. Nämlich: der Vater gibt die ewige Gewalt dem Sohne, und der Sohn hat sie vom Vater, und das alles von Ewigkeit her, sonst wäre es nicht eine ewige Gewalt; so ist der Heilige Geist da, der es durch Daniel redet. Denn solch hoch, heimlich Ding könnte niemand wissen, wo es der Heilige Geist nicht durch die Propheten offenbarte; wie droben [§ 9. 11.] oft gesagt, daß die heilige Schrift durch den Heiligen Geist gesprochen ist. Daneben ist der Sohn gleichwohl auch ein Menschenkind, das ist, ein rechter Mensch und Davids Sohn, dem solch ewige Gewalt gegeben wird. Also sehen wir, wie die Propheten das Wort „ewig“ wohl angesehen und verstanden haben, da GOtt durch Nathan zu David spricht: „Ich will meinen und deinen Sohn in mein ewiges Reich setzen.“

\quellseite{961}
43. Hier stößt sich nun Frau Klüglingin,\footnote{Wittenberger und Erlanger: „klüglingne“; Jenaer: „klüglinge“.} die Vernunft, die gemeinhin weiser ist denn GOtt selbst, und fragt: Wie kann GOtt seine ewige Gewalt von sich einem andern geben? Was behielte er selbst? Ist doch droben [§ 12] gesagt, daß GOtt Jes. 42, 8. spricht: „Ich will meine Ehre keinem andern geben, noch mein Lob den Götzen“; sonderlich einem Menschen kann er sie nicht geben, der nicht von ewig her gewesen ist, wie GOtt, sondern hier zeitlich angefangen, geboren, und sterblich ist; wie wir Christen von Jesu, Davids und Marien Sohne, bekennen und predigen. Solcher Art hochverständige Leute sind auch die Juden, Mahomet, Türken und Tartaren, die können das unbegreifliche Wesen GOttes in den Löffel oder Mus-Schale ihrer Vernunft fassen, und sagen, GOtt habe kein Weib, darum könne er keinen Sohn haben. Pfui, pfui, pfui dich an, Teufel, mit Juden und Mahomet, und alle die, so der blinden, thörichten, elenden Vernunft Schüler sind in diesen hohen Sachen, die niemand versteht, denn GOtt allein, und wie\footnote{Wittenberger und Erlanger: „wo viel“; Jenaer: „wie viel“. Cruciger: quantum.} viel der Heilige Geist uns davon durch die Propheten offenbart hat.

44. Wir Christen, aus dem Neuen Testament erleuchtet, können hierzu richtig, deutlich und fein antworten, also: Christus, unser HErr, hat zwo Geburten, oder zwo Naturen in Einer unzertrennten Person; denn er ist Ein Christus, nicht (wie der tolle Geist Nestori narrt) zween Christi. Nach der ersten Geburt hat er, nicht zeitlich, sondern von Ewigkeit her, vom Vater empfangen die ewige Gewalt oder Gottheit, und der Vater hat sie ihm gegeben ganz und völlig, wie er sie selbst hat von Ewigkeit. Nicht hat er sie ihm also gegeben, daß er sich derselbigen beraubt oder entledigt habe, sondern dieselbe Gewalt, und keine andere, die er von Ewigkeit ganz und völlig gehabt und in Ewigkeit behält, hat er dem Sohne gegeben. Denn es sind nicht zwo Gottheiten, sondern beider Personen ist eine einzige Gottheit, und bleibt recht geredet, Jes. 42, 8.: „Ich will meine Ehre keinem andern geben, noch mein Lob den Götzen.“ Denn der Sohn ist kein anderer Gott, noch Götze, sondern mit dem Vater ein einiger, rechter, ewiger Gott.

45. Hievon spricht er selbst Joh. 16, 15.: „Alles, was der Vater hat, das ist mein.“ Spricht nicht: Der Vater hat nichts mehr, ich habe es alles allein; oder: Der Vater hat es alles allein, ich habe nichts; sondern: Der Vater hat es alles; aber dasselbe „Alles“, das er hat, das ist mein. Das ist ja klarlich so viel gesagt, daß der Vater und Sohn eine einige Gottheit haben, und von demselben Alles des Vaters, das des Sohnes ist, hat es der Heilige Geist auch; wie er daselbst spricht: „Er wird es von dem Meinen nehmen.“ Von welchem „Meinen“? Ohne Zweifel von dem Meinen, das der Vater hat. Also nimmt der Heilige Geist von beiden, dem Vater und Sohne, dieselbige völlige ganze Gottheit von Ewigkeit her. Item, Joh. 5, 26.: „Wie der Vater das Leben hat in ihm selber, also hat er dem Sohne gegeben das Leben zu haben in ihm selber“; und, V. 21. 23.: „Wie der Vater [die] Todten auferwecket, und lebendig macht, also auch der Sohn machet lebendig, welche er will; auf daß sie alle den Sohn ehren, wie sie den Vater ehren.“ Das alles ist von der ersten, ewigen, göttlichen Geburt gesagt.

46. Nach der andern, zeitlichen, menschlichen Geburt, ist ihm auch die ewige Gewalt GOttes gegeben, doch zeitlich, und nicht von Ewigkeit her. Denn die Menschheit Christi ist nicht von Ewigkeit gewesen, wie die Gottheit; sondern, wie man zählt und schreibt, ist Jesus, Marien Sohn, dies Jahr 1543 Jahre alt. Aber von dem Augenblicke an, da Gottheit und Menschheit ist vereinigt in Einer Person, da ist und heißt der Mensch, Marien Sohn, allmächtiger, ewiger GOtt, der ewige Gewalt hat, und alles geschaffen hat und erhält, per communicationem idiomatum, darum, daß er mit der Gottheit Eine Person, und auch rechter GOtt ist. Davon redet er Matth. 11, 27.: „Alles ist mir vom Vater gegeben“, Matth. am letzten, V. 18.: „Mir ist alle Gewalt gegeben im Himmel und [auf] Erden.“ Welchem „Mir“? Mir, Jesu von Nazareth, Marien Sohn und Menschen geboren. Von Ewigkeit habe ich sie vom Vater, ehe ich Mensch ward; aber da ich Mensch ward, habe ich sie zeitlich empfangen nach der Menschheit, und heimlich gehalten bis auf mein Auferstehen und Auffahrt, da es hat sollen offenbar und verklärt werden; wie St. Paulus Röm. 1, 4. spricht: Er ist verkläret, oder „erweiset ein Sohn GOttes kräftiglich“. Johannes nennt es „verkläret
\quellseite{962}
“ (Cap. 7, 39.: „Der Heilige Geist war noch nicht [da], denn Jesus war noch nicht verkläret.“

47. Siehe nun, ob nicht Daniel fast auf den Schlag redet, wie Jesaias, von dem Menschensohne, der das ewige Reich von GOtt empfähet, und wie Nathan und David reden, daß GOtt wolle Davids Sohn in sein ewiges Reich zum Könige setzen; und daß solches geredet ist (spricht David) als von einem Menschen, der droben in der Höhe GOtt der HErr ist.

48. Ach! daß wir Christen solche unaussprechliche Gnade so reichlich, beide im Neuen und Alten Testament, erkennen, und nicht fröhlich sind, und danken, wie sich es gebühret. Wunder wäre es nicht, wenn ein Christenherz recht gründlich bedächte und ergriffe, daß es vor Freuden stürbe, und vor Freuden wieder lebendig würde. Was ist es doch groß Wunderes, daß Gott Mensch ist, mit uns Menschen redet, lebt, dazu für uns stirbt! David verstummet und erstarret vor Freuden, kann nicht mehr sagen, denn: „Was bin ich? Was ist mein Haus, daß du mich dahin bringest?“

49. Das ist noch alles von dem Texte 1 Chron. 18. geredet, auf welchen, wie droben [§ 17] gesagt, sich die letzten Worte Davids gründen, daß Christus müsse rechter GOtt und Mensch sein. Und was mehr aus solchem Texte quillt, wollen wir hernach (mit GOttes Hülfe) weiter sehen. Denn die Propheten nach David, und er selbst auch, viel daraus genommen haben von Christo, daß er GOtt und Mensch sei. Als der 110. Psalm, V. 1.: „Der HErr hat gesagt zu meinem Herrn: „Siehe dich zu meiner Rechten, bis ich deine Feinde lege zu deinem Fußschemel.“ Was kann\footnote{Erlanger: „heißen kann“.} „zu meiner Rechten sitzen“ anders heißen, denn GOtt gleich sitzen, das ist, in GOttes ewigem Reiche sitzen? Denn er sitzt ihm nicht zum Häupten, noch zum Füßen, weder höher noch niedriger; sondern zur Rechten, ihm gleich, daß der Himmel ebenso wohl sein Stuhl, und die Erde seine Fußbank ist; wie er spricht Matth. am letzten, V. 18.: „Mir ist gegeben alle Gewalt im Himmel und Erden“; und Marc. am 16. [V. 19.]: „Er ist aufgenommen gen Himmel, und sitzet zur rechten Hand GOttes.“ Und da er Matth. 22, 43. f. die Pharisäer fragt: „Ist Christus Davids Sohn, wie nennt er ihn denn im Geist (das ist, der Geist redet durch ihn) seinen HErrn? da er spricht: Der HErr hat gesagt zu meinem HErrn: „Setze dich zu meiner Rechten“, konnten sie ihm nichts antworten.

50. So können die Juden\footnote{Zu der Erlanger: „Sie können“ statt: So können die Juden.} auch heutiges Tages bis in Ewigkeit nicht darauf antworten; sie lästern wohl sehr, allein zuwider uns Christen, mit ungeschickter, boshafter Gaukelei, auch wider ihre eigenen alten Rabbinen und Lehrer. Davon magst du fromm daselbst lesen. Wir haben aber das Neue Testament, welches sich nicht allein reimt mit diesem Psalm, und der Psalm ihm ähnlich ist; daran wir Christen (wie nun oft gesagt [§ 32]) genug haben, das Alte Testament auf unsern Verstand zu ziehen, und kann uns nicht fehlen, sondern aller anderer Verstand muß fehlen. So ist Christus selbst da mit seinen Aposteln, die zeugen und zeigen uns diesen Verstand mit reichen Worten und Werken, und ist dieser Psalm auch der Lehrer einer, da die drei Personen unterschiedlich, in Einer Gottheit, verkündigt sind; welches allein der Propheten und Christen Verstand ist, durch den Heiligen Geist gegeben. Juden, Mahomet, Vernunft sollen hiervon nichts wissen. Da ist der Vater, der spricht: „Setze dich zu meiner Rechten“; und ist Davids Sohn, Christus, zu dem er solches spricht.

51. Nun ist der Vater nicht Christus, oder Davids Sohn, und Christus ist nicht der Vater, und soll doch dem Vater zur rechten Hand gleich sitzen, einerlei Reich, Gewalt, Ehre und alles haben. Denn GOtt leidet keinen Gleichen, der neben ihm in gleicher Gewalt und Ehre sitze; darum muß Christus, Davids Sohn, rechter und mit dem Vater ein einiger GOtt sein, der ihm gleich sitze; sintemal nicht mehr, denn ein einiger GOtt sein muß; wie das erste Gebot sagt: „Du sollt keine anderen Götter neben mir haben.“ So ist der Heilige Geist auch da, als der rechte einige GOtt, der durch David und alle Propheten mit uns Mensch redet, und alle Wahrheit von der Gottheit uns offenbaret und lehrt; wie David spricht: „Der Geist des HErrn hat durch mich geredet“, und Christus selbst Matth. 22, 45.: „Wie nennt David, durch den Geist, Christum seinen HErrn, so er sein Sohn ist?“ Freilich, ohne den Geist würde er
\quellseite{963}
es weder nennen noch kennen, wie Christus sein Sohn und HErr wäre. Der Heilige Geist aber ist nicht Christus der Sohn, noch der Vater, und kann doch kein anderer GOtt sein. Also schließt sich’s gewaltiglich, daß ein einiger GOtt, und doch unterschiedliche drei Personen sind, Vater, Sohn, Heiliger Geist, von Ewigkeit in Ewigkeit.

52. Hier möchte jemand bewegen, warum David spricht: „Was bin ich? Was ist mein Haus?“ Nem: „Du hast mich angesehen, wie einen Menschen, der in der Höhe GOtt der HErr ist.“ Hat doch GOtt nicht zu ihm gesagt: Du sollst mein Sohn sein, ich will dich in mein ewiges Reich setzen; sondern so sagt er: „Dein Sohn soll mein Sohn sein, demselbigen will ich in mein ewiges Reich setzen.“ Wie ändert denn David die Worte GOttes, und deutet sie auf sich, als sei er der Mensch, der in der Höhe GOtt der HErr ist? Wohlan, David ist der Vater dieses Sohnes (wie du hörst), und aus seinem Hause und von seinem Fleisch und Blute soll er kommen. Nun ist es also in der Natur, daß sich ein Vater der Herrlichkeit seines Sohnes ja so hoch, und wohl höher freuet denn der Sohn selbst, denn er alle Ehre und Gut gönnt, und viel mehr, denn ihm selbst. Wiederum, der Sohn und Schmach, dem Sohn angelegt, viel höher ihn verdrießt, denn so es ihm selbst widerführe. Darum rühmt sich nicht allein David, sondern auch sein ganzes Haus (wie er spricht: „Was ist mein Haus?“) der Herrlichkeit, daß aus ihrem Fleisch und Blute ein solcher Sohn kommen soll, der zur rechten Hand GOttes sitzen werde.

53. Also liest man in den Historien, daß ein Vater, den sie Chilo nennen, vor Freuden gestorben sei, da er gehört, daß sein Sohn den Sieg behalten hatte in Olympiis.\footnote{Das ist, in den olympischen Spielen.} Und eine Römerin, da sie nicht anders wußte, ihr Sohn wäre im Kriege mit andern von Hannibal bei Cannae erschlagen, und er plötzlich gesund heimkommt, in dem Blick, da sie ihn ansieht, fällt sie vor Freuden dahin, und stirbt. Also ist David hier so voller Freuden und Geistes, daß er schier nicht weiß, wie, und was er reden soll, und nimmt sich seines Sohnes, seines Fleisches und Blutes Ehre nicht anders an, denn als wäre sie ihm selbst geschehen.

54. Zudem so steht solcher Sohn Davids noch zur Zeit in seinem Fleisch und Blute, und ist nichts davon vorhanden denn er selbst, David, in seiner Person und mit seinem Fleisch und Blute, daher künftig der Sohn kommen sollte. Denn diese Geschichte und Worte sind geschehen zeitlich, ehe denn Nathan, Davids Sohn (von welches Schnur Christus gekommen ist, Luc. 3, 31.), geboren ward; ja, seine Mutter Bathseba, war noch nicht Davids Weib worden, sondern noch Uriä Weib, eine gute Zeit vor dem Fall und Ehebruch. Darum ist [es] nicht eine unförmliche Rede von einem Vater, wenn er von seines Sohnes Herrlichkeit GOtt lobt und dankt, mit diesen oder dergleichen Worten: Ach du lieber GOtt! Was bin ich? wofür siehst du mich doch an, daß du mich so hoch ehrst, und aus meinem Fleisch und Blut einen solchen Herrn machst? Ich bin es ja, des die Ehre und Freude ist; denn es ist ja mein Fleisch und Blut, das noch jetzt in mir und an mir ist, künftig [aber]\footnote{Dieser Zusatz, welcher das Verständnis erleichtert, findet sich in der Wittenberger.} soll geboren werden.

55. Daher wird auch unser HErr Christus oft in den Propheten mit dem Namen seines Vaters David genannt, Hos. 3, 5.: „Darnach werden sich die Kinder Israel bekehren, und den HErrn, ihren GOtt, und David, ihren König, suchen, und den HErrn und seine Gnade ehren, in der letzten Zeit.“ Hier heißt David unser HErr Christus, und wird in gleiche Ehre mit GOtt gesetzt, und HErr genannt, den sie suchen und ehren werden. Es ist einerlei Suchen, damit sie GOtt und ihren König suchen und ehren sollen; gleiche mit einerlei Glauben den Vater und Sohn ehren, nicht mit einem andern Glauben den Vater, und mit einem andern den Sohn ehren. Und ist hier die dritte Person, der Heilige Geist, der solches durch Hoseam mündlich redet, und uns glauben lehrt.

56. Jer. Ezech. 34, 23. 24.: „Ich will ihnen einen einigen Hirten erwecken, der sie weiden soll, nämlich, meinen Knecht David, der wird sie weiden, und ihr Hirte sein; ich will ihr GOtt sein, und mein Knecht David wird Fürst unter ihnen sein.“ Hier heißt Christus David und GOttes Knecht, wie er auch Jes. 52, 13. „GOttes Knecht“ heißt, und au\footnote{„au“ fehlt in der Jenaer und in der Erlanger.} viel Ehren mehr. Dazu Paulus Phil. 2. macht einen Knecht aus
\quellseite{964}
ihm, der ihn\footnote{So die Wittenberger. Jenaer und Erlanger: ihnen.} doch immer und immer einen rechten GOtt predigt, da er spricht, V. 5–7.: „Ein jeglicher sei gesinnet wie Jesus Christus, welcher, ob er wohl in göttlicher Gestalt war, hielt er es nicht für einen Raub, GOtt gleich sein, sondern äußerte sich selbst, und nahm Knechtlich Gestalt an“ etc. Laß uns den Apostel fragen, wie er so ungeschickt mag reden. Ist Christus GOtt gleich, wie kann er ein Knecht, und in Knechts Gestalt sein? Ist er ein Knecht, wie kann er GOtt gleich, und in GOttes Gestalt sein? Aber wir Christen verstehen und wissen solches alles wohl; aber die Juden verstocken sich mit diesem Spruche Ezechiels getrost, und wollen gewiß sein ihres Sinnes (ich wollte sagen, Wahnsinnes); die laß fahren.

57. Item, Jeremia 30, 8. 9.: „Zu der Zeit, spricht der HErr, will ich das Joch von deinem Halse zerbrechen, und deine Bande zerreißen, und sollen nicht mehr drinnen den Fremden (oder wie Fremde) dienen; sondern werden dienen dem HErrn, ihrem GOtt, und ihrem Könige David, den ich ihnen erwecken will.“ Hier heißt Christus auch David, wie die Juden, alte und neue, diesen Spruch müssen von Messia verstehen, ohne daß sie das Joch und die Bande nicht recht verstehen, wollen das Gefängniß zu Babylon daraus machen. Aber die ganzen drei Capitel nach einander reden stark von der Erlösung, so Messia thun soll, das ist, vom Tod und Sünden, so das Gesetz auf uns treibt, davon Juden und Vernunft nichts wissen. Es ist ein Verstand für die Christen, und was christlich gewesen ist von Anfang der Welt etc.

58. Gleichwohl macht Jeremia mit diesem Spruch diesen seinen König David zum rechten GOtt, da er GOtt und diesen David zusammen setzt in einerlei Dienst und Ehre, die ihm das Volk Israel thun solle. Denn wo dieser David nicht rechter GOtt wäre, würde ihn GOtt nicht neben sich setzen und sprechen: „Sie sollen dienen ihrem GOtt, und ihrem Könige David“; denn es heißt: „Du sollst keinem andern Gott dienen, sondern den Herrn, deinem Gott, sollst du fürchten, und ihm allein dienen“, 5 Mos. 6, 13. und 10, 20. Und stimmen die Worte Jeremiä mit dem Spruche 1 Chron. 18, 11. 17.: „Ich will deinen Sohn in mein eigenes Reich setzen, der ein Mensch ist, und zugleich droben überwärts, in der Höhe, GOtt der HErr ist, der mit dem Vater gleich geehrt, und ihm gleich gedient wird. Hierbei muß der Heilige Geist die dritte Person sein, der solches durch Jeremia redet, und uns lehrt, daß wir es glauben und verstehen, und ist ein einiger GOtt, neben dem wir keinen andern Gott ehren, noch dienen.

59. Hiermit sollte es nun wohl genug sein von dem Text 1 Chron. 18., daraus sich Davids letzte Worte gründen, nämlich, daß Christus sei GOtt und Mensch, von David geboren, daß wir nun wohl möchten wieder zu den letzten Worten Davids uns machen, dieselben zu enden, darinnen er Christum seinen Sohn bekennt, und rühmt für seinen GOtt, nach dem Spruch [1 Chron. 18, 17.]: „Du hast mich angesehen, wie einen Menschen“, qui superne vel in supernis, vel in excelsis est Dominus Deus, „der hoch droben, oder in der Höhe, GOtt der HErr ist.“ Denn daß unser lateinischer Text in vocativo sagt, Domine Deus, gibt keinen Verstand, so wenig als 2 Sam. 7, 19.: Ista est lex Adam, Domine Deus; das besser wäre: Ista est lex vel forma hominis, Domini Dei; seu, qui sit Dominus Deus etc. Davon jetzt genug gesagt. Aber weil die Materie so gut ist, und, leider, wir der geringste Haufe mit den Propheten und Aposteln sind, die sich um den Christum, den gekreuzigten David und ewigen GOtt, annehmen, wollen wir, ehe wir die letzten Worte Davids lesen\footnote{lesen = verabschieden.} und aus Ende bringen, weiter davon reden, zu stärken uns in unserm Glauben, zu Verdrieß allen Teufeln, Juden, Mahometisten, Papisten, und was diesem Sohne Davids feind ist.

60. Und erstlich wollen wir Mosen, den Hauptbrunn, Quelle, Vater und Meister aller Propheten, vor uns nehmen; versuchen, ob er sich wolle lassen einen Christen machen, und uns bekennen, weil Christus selbst ihn tauft, Joh. 5, 46., und spricht: „Moses hat von mir geschrieben.“ Hat er von ihm geschrieben, so hat er freilich von ihm auch geweissagt, gepredigt, und befohlen allen Propheten nach ihm, von Christo zu schreiben und zu predigen; wie sie mit allem Fleiß gethan, also, daß auch alle Juden, jung und alt, wissen zu sagen, daß ein Messias hat kommen sollen. Aber Mose ist ihnen verborgen, daß sie nicht wissen können, wo er liegt. Darum wollen wir zween treue, gewisse Legaten oder Botschafter
\quellseite{965}
ausrichten und abfertigen, die ihn suchen, finden, aufwecken und herbringen sollen; die heißen Johannes Evangelista und Paulus Apostolus. Was gilt’s, sie werden treffen, und nicht fehlen! Doch, daß du nicht vergessenest, was ich droben [Vorrede § 8] gesagt habe: ich wolle diesmal das vornehmen, wo der hebräische Text sich gerne gibt und reimt mit dem Neuen Testament, daß solches sei und sein solle der einige rechte Verstand der Schrift; alles andere, was Juden, Hebräisten, und wer sie sind, nach ihrer zerstipten\footnote{zerstipt = durch Punkte zerrissen. Cruciger: ex ipsorum punctis et apicibus misere contorta. — Im Englischen ist noch das Stammwort to stipple — in Punkten ausführen. — Erlanger: zerstieben. — In der Wittenberger und in der Jenaer: „zerstipten“ und „zemarterten“.} und zernatterten, gezwungenen Grammatica davider plaudern, soll uns gewißlich eitel Lügen sein.

61. Wohlan, St. Johannes fähet sein Evangelium also an: „Im Anfang war das Wort, und das Wort war bei GOtt, und GOtt war das Wort. Daselbige war im Anfange bei GOtt. Alle Dinge sind durch daselbige gemacht, und ohne daselbige ist nichts gemacht, was gemacht ist.“ Dies sind St. Johannis, oder vielmehr des Heiligen Geistes Reden, der alle Dinge lebendig macht; nun laß sehen, ob er Mosen hiermit finden, und von Todten auferwecken könne. Denn Moses hört sehr leise, und kommt frei daher und spricht: Hier bin ich. Denn eben wie du, Johannes, vom Worte redest, so habe ich auch geredet, und rede noch also, und du nimmst mir die Worte aus meinem Munde. Denn also habe ich auch geschrieben, im Anfange meines Buchs, vom Anfange der Creatur [1 Mos. 1, 3. ff.]: „GOtt sprach: Es werde Licht; und es ward Licht. GOtt sprach: Es werde eine Feste zwischen den Wassern. GOtt sprach: Es sammle sich das Wasser unter dem Himmel an sonderer Oerter, daß man das Trockene sehe. GOtt sprach: Es lasse die Erde aufgehen Gras und Kraut. GOtt sprach: Es werden Lichter an der Feste, und scheiden Tag und Nacht“; und so fortan.

62. Hier stimmt Mose mit Johanne überein, daß im Anfange der Creatur ein Wort sei gewesen, durch welches GOtt alles gesprochen, das ist, geschaffen und gemacht hat. Und Mose lispet\footnote{In den alten Ausgaben: liebet.} oder stammelt hier nicht, sind auch nicht dunkle, ungewisse Reden; die Grammatica ist auch gewiß, daß, wo ein Sprecher ist, da ist ein λόγος, Wort oder Rede. Ob nun die Juden, Ketzer, Mahomet ihre eigene Deutung hier erträumen, dem christlichen Glauben zuwider, da fragen wir nichts nach. Wir haben den Text und Grammaticam Mosis für uns, der steht dürre und klar da, daß im Anfange und außer allen Creaturen habe GOtt gesprochen, und sei ein Wort da, durch welches GOtt alles spricht; wie uns St. Johannes schreibt, Joh. 1, 1. ff. Welches Evangelium die lieben Väter Hilarius, Augustinus, Cyrillus etc. reichlich und gewaltiglich haben gehandelt, daß nicht noth ist, davon jetzt weiter zu schreiben, ihre Bücher sind vorhanden. Genug ist es diesmal, daß wir selbst hier sehen und hören, wie Mose ungezwungen und ungedenket, von sich selbst, so gleich und eben mit Johanne sich reimt, daß es auch die blinde Vernunft nicht leugnen kann, sondern nach der gewissen Grammatica bekennen muß, daß sie beide einerlei Weise und Rede führen vom Worte, durch welches Gott im Anfange alles geschaffen und gemacht hat.

63. Denn sie wollen alle beide, Mose und Johannes, anzeigen, womit, und durch welch Werkzeug, oder woraus GOtt solch groß Werk, die ganze Welt, gemacht habe. Aber, da ist kein Rüstzeug, kein Holz noch Stein, lauter Nichts ist da, daraus die Welt geschaffen ist; sondern allein durchs Wort ist alles gemacht. Das Wort aber ist nicht gemacht, sondern es ist bei GOtt im Anfang, da er alles machte, wie Mose hier sagt: „GOtt sprach: Es werde“ dies und das etc. Durchs Wort, sagt er, ist alles worden. Nun kann außer der Creatur bei GOtt nichts sein, das nicht GOtt selber ist; darum muß das Wort GOtt selber sein, so groß und mächtig als GOtt selbst, weil alles Ding durch dasselbige gemacht wird; und kann doch nicht die Person sein, die das Wort spricht. Es muß der Sprecher und das Wort zweierlei sein. Wiederum müssen es nicht zween Götter sein, weil nicht mehr denn ein einiger, rechter, wahrhaftiger GOtt ist, und sein muß ein einiger Schöpfer Himmels und der Erde, nicht zween oder drei Schöpfer oder Götter. Also zeuget Mose mit St. Johanne, daß GOtt und das Wort gewißlich zwo unterschiedliche Personen, und gleichwohl alle beide ein einiger Schöpfer und GOtt sind, in dem einigen göttlichen Wesen unzertrennt.

\quellseite{966}
64. Also hat David Mosen gelesen und verstanden, da er spricht Ps. 33, 6.: „Der Himmel ist durchs Wort des HErrn gemacht, und alle sein Heer durch den Geist seines Mundes.“ „Gemacht“ ist der Himmel und alles, das drinnen und dran ist (spricht er). Lieber, woraus? Aus nichts. Durch was? Durch sein Wort, und durch den Geist seines Mundes. Redet hier nicht David auch gleich wie Mose, und will auch schier mit gleicher Rede sagen: (G)Ott sprach: Es werde der Himmel, und es ward der Himmel? Ist aber der Himmel,\footnote{Die Worte: „und es ward der Himmel? Ist aber der Himmel“ fehlen in der Originalausgabe.} mit allem, das darinnen ist, durch (G)ottes Sprechen oder Wort worden und gemacht, so ist ohne Zweifel die Erde mit allem, das darinnen, auch durch dasselbige Wort worden und gemacht. Nun ist das Wort nicht der Himmel, noch Erde, noch etwas, das darinnen, oder mit denselben durchs Wort gemacht ist. Darum muß es GOtt selber sein, und doch ja eine andere Person von dem Sprecher, der durchs Wort alles macht, in einigem, unzertrennten Wesen, göttlicher Gewalt, Macht und Wirkung. Haben wir aber das Wort, so wollen wir die dritte Person wohl finden, da David hier sagt: „Und durch den Geist seines Mundes alle sein Heer.“

65. Er sagt nur einmal „gemacht“: „Der Himmel ist durchs Wort des HErrn gemacht, und alle sein Heer durch den Geist seines Mundes.“\footnote{In der Wittenberger, Jenaer und Erlanger: „durch den Mund seines Geistes“. Im Lateinischen fehlt hier das Schriftzitat.} Zwei Personen nennt er unterschiedlich, nämlich den HErrn, sein Wort und seinen Geist; und jetzt doch nicht mehr, denn einen Macher, und jetzt doch nicht mehr, denn Einen Macher, ohne allem Unterschied. Gemacht ist es alles. Von wem? Von Einem Macher, der ist der HErr, Wort, Geist. Nicht macht der HErr sein eigen Werk besonders, das Wort macht nicht sein eigen Werk besonders, der Geist macht nicht sein eigen Werk besonders; es sind alle drei unterschiedliche Personen ein einiger Macher eines jeglichen Werks; und ein jeglich Werk ist aller drei Personen, als eines einigen Machers oder Meisters Werk. Denn wie der HErr den Himmel macht, so macht das Wort auch denselben und keinen andern Himmel, so macht der Heilige Geist auch denselben und keinen andern Himmel; ein einiger ist, der es macht, und ein einiges Werk ist, das alle drei Personen machen.

66. Darum ist hier einem Christen wohl zu merken, daß er, wie Athanasius singt in seinem Symbolo, nicht die Personen in Eine Personenmenge, oder das einige göttliche Wesen in drei Personen theile oder trenne. Denn wo ich einer jeglichen Person von außen, in der Creatur, ein sonderlich Werk gäbe, da die andern zwo nicht mit zu thun haben sollten, so habe ich die einige Gottheit zertrennt, und drei Götter oder Schöpfer gemacht; das ist falsch. Wiederum, wo ich einer jeglichen Person inwendig\footnote{inwendig = innerhalb.} der Gottheit, oder außer und über der Creatur, nicht ein sonderlich Unterschied\footnote{„ein sonderlich Unterschied“ = eine besondere Eigenthümlichkeit. (Erlanger: si non discernas singulas personas singulis suis proprietatibus.} gebe, die den andern zween nicht gebührt, so habe ich die Personen in Eine Person gemengt; das ist auch falsch. Hierher gehört die Regel St. Augustini: Opera Trinitatis ad extra sunt indivisa, die Werke, so von GOtt auswendig der Gottheit gemacht, sind nicht [zu] zertheilen,\footnote{In der Jenaer und in der Erlanger: „zertheilen“; in der Wittenberger: „zu teilen“. Die Vorsilbe „ge“ ist unsrer. Vgl. §60: „gestipt“, „gemartert“.} das ist, man soll die Personen nicht theilen in die Werke, einer jeglichen von außen ihr unterschiedlich Werk zueignen; sondern die Person soll man inwendig der Gottheit unterscheiden, und doch allen dreien auswendig ein jeglich Werk ohne Unterschied zueignen.

67. Als, daß ich Exempel gebe: Der Vater ist mein und dein GOtt und Schöpfer, der mich und dich gemacht hat; eben dasselbe Werk, das ich und du sind, hat auch der Sohn gemacht, ist gleichwohl mein und dein GOtt und Schöpfer als der Vater. Also, der Heilige Geist hat eben dasselbige Werk, das ich und du sind, gemacht, und ist mein und dein GOtt und Schöpfer, gleichwohl als der Vater und Sohn. Noch sind es\footnote{„es“ fehlt in der Erlanger.} nicht drei Götter oder Schöpfer, sondern ein einiger GOtt und Schöpfer unser aller beide. Hier, mit diesem Glauben, verwahre ich mich
\quellseite{967}
vor der Ketzerei Arii, und seines Gleichen, daß ich das einige göttliche Wesen nicht zertrenne in drei Götter oder Schöpfer; sondern behalte im rechten christlichen Glauben nicht mehr, denn den einigen GOtt und Schöpfer aller Creaturen.

68. Wiederum, wenn ich nun über und außer der Schöpfung oder Creatur gehe in das inwendige, unbegreifliche Wesen göttlicher Natur, so finde ich, wie mich die Schrift lehrt (denn Vernunft ist hier nichts), daß der Vater eine andere unterschiedliche Person ist von dem Sohn, in der einigen unzertrennten einigen Gottheit; sein Unterschied ist, daß er Vater ist, und die Gottheit nicht vom Sohne, noch von jemand hat; der Sohn eine unterschiedliche Person ist vom Vater, in derselbigen einigen väterlichen Gottheit; sein Unterschied ist, daß er Sohn ist, und die Gottheit nicht von sich selbst, noch von jemand, sondern allein vom Vater hat, als ewiglich vom Vater geboren; der Heilige Geist eine unterschiedliche Person ist vom Vater und Sohn, in derselbigen einigen Gottheit; sein Unterschied ist, daß er der Heilige Geist ist, der vom Vater und Sohn zugleich ausgeht ewiglich, und die Gottheit nicht von sich selbst, noch von jemand hat, sondern beide vom Vater und Sohn zugleich; und das alles von Ewigkeit in Ewigkeit. Hier, mit diesem Glauben, verwahre ich mich vor der Ketzerei Sabellii und seines Gleichen, vor Juden, Mahomet, und wer sie mehr sind, die klüger sind denn GOtt selbst, und menge die Personen nicht in eine einige Person, sondern behalte im rechten christlichen Glauben drei unterschiedliche Personen in dem einigen, göttlichen, ewigen Wesen, die doch alle drei gegen uns und die Creaturen ein einiger GOtt, Schöpfer und Wirker ist aller Dinge.

69. Dies alles ist vielleicht uns Deutschen scharf oder subtil, und sollte billiger in den Schulen bleiben. Aber weil der Teufel den Schwantz regt in dieser letzten Zeit, als wollte er gerne allerlei Ketzerei wieder aufwecken, und die Welt ohne das lüstern und toll worden ist, Neues zu hören, und überdrüssig der heilsamen Lehre (wie St. Paulus 2 Tim. 4, 3. weissagt), damit dem Teufel die Thüren aufgesperrt sind, hinein zu führen, was er will, so ist es nütze und noth, daß doch etliche, beide Laien und Gelehrte, sonderlich Pfarrherren, Prediger und Schulmeister, von solchen nöthigen Artikeln unseres Glaubens auch lernen denken, und deutsch reden. Wenn es aber zu schwer ist, der bleibe mit den Kindern bei dem Catechismo, und bete wider den Teufel und seine Ketzerei, wider Juden und Mahomet, damit er nicht geführet werde in Anfechtung. Darum, weil wir darauf kommen sind, wollen wir denen, so es gerne haben, mehr Gleichniß geben von dem Artikel, daß die einige Gottheit nicht zertrennt, noch die Personen in einander gemenget werden sollen, unsern Glauben zu stärken und zu bekennen.

70. Am Jordan, da St. Johannes den HErrn taufte, that sich der Himmel auf, und fuhr der Heilige Geist hernieder leiblich, in einer Tauben Gestalt, und des Vaters Stimme ward gehört: „Dies ist mein lieber Sohn, an dem ich Wohlgefallen habe“, Luc. 3, 22. Hier ist die Taube eine Creatur, welche nicht allein der Heilige Geist geschaffen hat, sondern auch der Vater und Sohn; wie gesagt, daß opera Trinitatis ad extra sunt indivisa: was Creatur ist, das hat GOtt Vater, Sohn und Heiliger Geist zugleich, als ein einiger GOtt, gemacht; noch heißt die Taube allein der Heilige Geist, oder, wie Lucas sagt, ist allein der Heilige Geist hernieder gefahren in der Tauben Gestalt, und würde in seinem Weg der christliche Glaube leiden, daß du wolltest sagen von der Taube: Das ist GOtt der Vater, oder: Das ist GOtt der Sohn; sondern mußt sagen: Das ist GOtt der Heilige Geist; ob wohl GOtt der Vater, Sohn und Heiliger Geist ein einiger GOtt ist; daß du ganz recht sagest von der Taube: Das ist GOtt, und ist kein anderer GOtt mehr; und doch unrecht sagest: Das ist GOtt der Vater, oder GOtt der Sohn; sondern mußt sagen: Das ist GOtt der Heilige Geist.

71. Also ist da die Stimme, so da spricht: „Dies ist mein lieber Sohn“ etc., eine Creatur, die nicht allein der Vater geschaffen hat, sondern auch der Sohn und Heiliger Geist, wie gesagt, opera Trinitatis etc., daß außer der Gottheit alle Creaturen zugleich sind von allen dreien Personen, als von einem einigen GOtt, geschaffen, und gegen der Creatur alle drei Personen ein einiger GOtt ist. Und wiederum, die Creatur gegen die drei Personen einerlei, und nicht dreierlei Werke sind, noch heißt und ist dieselbe Stimme allein des Vaters; und könnt, als ein Christ, hier von der Stimme nicht sagen: Das ist GOtt der Heilige Geist, oder: Das ist GOtt der Sohn, sondern mußt sagen: Das ist GOtt der Vater; ob wohl GOtt
\quellseite{968}
der Heilige Geist und GOtt der Sohn und GOtt der Vater ein einiger GOtt ist; daß du ganz recht sagest von der Stimme: Das ist GOtt, und ist kein anderer GOtt mehr; doch unrecht sagest: Das ist GOtt der Sohn, oder GOtt der Heilige Geist; sondern mußt sagen: Das ist GOtt der Vater.

72. Gleich dem ist zu reden von der Menschheit Christi: Die ist an sich selbst eine rechte Creatur, geschaffen zugleich vom Vater, Sohn und Heiligen Geist; und ist nicht zu leiden im Glauben, daß der Vater allein, oder der Sohn allein, oder der Heilige Geist allein diese Creatur oder Menschheit geschaffen habe; sondern ist ein opus indivisum Trinitatis, ein Werk, welches alle drei Personen als ein einiger GOtt und Schöpfer einerlei Werks geschaffen hat; wie der Engel Gabriel zu der Jungfrau Maria sagt, Luc. 1, 35.: „Der Heilige Geist wird über dich kommen, und die Kraft des Allerhöchsten wird dich überschatten.“ Nicht allein ist der Heilige Geist da über dir (spricht er), sondern auch „der Allerhöchste“; das ist, der Vater wird dich überschatten mit seiner Kraft, das ist, durch seinen Sohn oder Wort; auch so soll „das in dir geboren wird, des Allerhöchsten Sohn sein und heißen“, daß also die ganze Dreifaltigkeit als ein einiger Schöpfer hier ist, und das einige Werk, die Menschheit, geschaffen und gemacht hat, und doch die Person des Sohnes allein damit vereinigt, und Mensch worden, nicht der Vater, noch Heiliger Geist.

73. Und kann man von diesem Menschen nicht sagen: Das ist GOtt der Vater; oder: Das ist GOtt der Heilige Geist; sondern mußt sagen: Das ist GOtt der Sohn; obwohl Gott der\footnote{„der“ fehlt in der Erlanger.} Vater, Sohn und Heiliger Geist ein einiger GOtt ist, daß du ganz recht sagest von dem Menschen: Das ist GOtt, und ist kein anderer GOtt mehr; doch unrecht sagest: Das ist GOtt der Vater; oder: Das ist GOtt der Heilige Geist; sondern mußt sagen: Das ist GOtt der Sohn; wie St. Paulus Col. 2, 9. sagt: „Denn in Christo wohnet die ganze Fülle der Gottheit“, und ist doch damit der Vater und Heilige Geist derselben Gottheit nicht beraubt, sondern mit dem Sohn und Menschen Christo ein einiger GOtt. Hieraus siehst du, wie die drei göttlichen Personen unterschiedlich, inwendig der Gottheit, zu glauben, und nicht in Eine Person zu mengen sind, und doch das göttliche einige Wesen nicht zu trennen, oder drei Götter zu machen; sondern äußerlich, gegen die Creatur, ein einiger Schöpfer sei, so gar einig, daß auch die Creatur, so die Personen unterschiedlich an sich nehmen, aller drei Personen als einiges GOttes einerlei Werk sind.

74. Solch hoch Ding etlichermaßen zu begreifen, geben die Doctores, sonderlich Bonaventura, ein grob Gleichniß. Als: Wenn drei Jungfrauen Einer unter sich ein Kleid anzögen, da sie alle drei das Kleid angriffe, und der dritten anzögen, und die dritte selbst auch mit gleich zugriffe: da ziehen alle drei das Kleid der dritten an, und wird doch allein die dritte mit dem Kleide angezogen, und nicht die andern zwo. Also soll man hier verstehen, daß alle drei Personen, als ein einiger GOtt, die einige Menschheit geschaffen, und mit dem Sohne vereinigt habe in seine Person, daß allein der Sohn Mensch sei, und nicht der Vater, noch Heiliger Geist. Eben so soll man auch verstehen die Taube, so des Heiligen Geistes Person an sich nimmt, und die Stimme, so des Vaters Person an sich nimmt. Item, die feurigen Zungen am Pfingsttage, darin des Heiligen Geistes Person offenbar wird. Item, der Wind, und was man mehr vom Heiligen Geist predigt, daß er thue in der Christenheit oder heiligen Schrift.

75. Hier fragt sich's billig: Warum sprechen wir denn, oder vielmehr, warum lehrt uns denn die Schrift also sagen: Ich glaube an GOtt Vater, Schöpfer Himmels und der Erde, und nicht auch den Sohn Schöpfer nennen? Item: An Jesum Christum, der empfangen ist vom Heiligen Geiste. Item, daß der Heilige Geist lebendig mache, und durch die Propheten geredet habe. Hier werden ja äußerlich den Personen ihre besonderen unterschiedlichen Werke zugeeignet, wie sie selbst unterschieden sind. Dies ist einfältigen Christen vielleicht auch zu scharf, die mögen bei ihrem einfältigen Glauben bleiben, daß GOtt Vater, Sohn und Heiliger Geist Ein Gott sei etc. Doch muß man in der Christenheit hiervon reden, und lernen verstehen, dem Teufel und seinen Werken zu widerstehen. Erstlich ist es\footnote{„es“ fehlt in der Erlanger.} gewiß, daß GOtt will von uns erkannt sein hier im Glauben, denn
\quellseite{969}
ewiglich im Schauen, wie er sei ein einiger GOtt, und doch drei Personen; das ist unser ewiges Leben, Joh. 17, 3. Hierzu hat er uns sein Wort und die heilige Schrift gegeben, mit großen Wunderzeichen und Werken bestätigt, daß wir darinnen lernen sollen. Denn, sollten wir ihn also erkennen, müßte\footnote{So die Jenaer und das Lateinische. Wittenberger und Erlanger: müßte.} er es wahrlich uns lehren, und sich gegen uns offenbaren und erscheinen, von uns selber würden wir nicht in Himmel steigen, und finden, was GOtt sei, oder wie sein göttlich Wesen gethan ist. Nun, hierzu braucht er sichtbarerlich seiner Creatur, wie die Schrift uns lehrt, auf daß wir es ergreifen mögen, denn unsichtbare Creaturen bewegen unsere Sinne nicht.

76. Demnach mußt du nun die Creatur [auf] zweierlei Weise ansehen. Aufs erste als eine Creatur oder Werk an sich selbst, absolute, so und so von GOtt geschaffen, oder gemacht. Auf diese Weise sind alle Creaturen GOttes, das ist, aller drei Personen gleich einerlei Werk, ohne allen Unterschied; wie gesagt ist. Denn sie geben uns nach solcher Weise keine unterschiedliche Offenbarung der drei Personen, weil sie alle gleich einerlei Werk sind der drei Personen, als des einigen GOttes. Zum andernmal mußt du sie ansehen, nicht an ihr selbst, absolute, sondern relative, nach ihrem Brauch, wie GOtt derselben braucht gegen uns. Hier nimmt GOtt sein Geschöpf, das alle drei Personen geschaffen, als ein einiger GOtt, hat, und braucht desselben zum Bilde oder Form, oder Gestalt, darin er sich offenbart, und erscheint. Hier werden unterschiedliche Bilder, Gestalt, oder Offenbarung der drei unterschiedlichen Personen. Also braucht er der Taube, daß sie soll ein Bild oder Offenbarung sein, darin sich der Heilige Geist offenbart; und ist ein unterschiedlich Bild, das nicht den Vater noch Sohn uns zeigt, sondern allein den Heiligen Geist, unterschiedlich. Denn der Vater, Sohn und Heiliger Geist will, daß die Taube soll unterschiedlich uns allein die Person des Heiligen Geistes zeigen und offenbaren, damit wir gewiß werden, daß GOttes einiges Wesen gewißlich drei unterschiedliche Personen seien von Ewigkeit. Darum spricht Lucas Cap. 3, 22.: „Der Heilige Geist fuhr hernieder in leiblicher Gestalt, wie eine Taube.“

77. Eben so reden wir von dem Sohne, daß er uns ist offenbart in der Menschheit, oder, wie St. Paulus Phil. 2, 7. redet, in knechtlicher Gestalt, geberdet wie ein anderer Mensch. Und diese Gestalt, oder Menschheit, ist nicht des Vaters oder Heiligen Geistes Bild oder Offenbarung, ob sie wohl beide, des Vaters, Sohnes und Heiligen Geistes gleiche einerlei Creatur ist, sondern ist unterschiedlich eine Gestalt und Offenbarung allein des Sohnes. Denn so hat es GOtt, das ist, dem Vater, Sohne und Heiligen Geiste gefallen, daß der Sohn durch diese Gestalt oder Form der Menschheit unter den Menschen offenbart und erkannt würde als eine unterschiedliche Person vom Vater und Heiligen Geiste, in einem einigen einigen Wesen göttlicher Natur. Dem gleich soll man vom Vater reden, daß er uns offenbart ist in der Stimme. Diese Gestalt oder Form ist nicht des Sohnes oder Heiligen Geistes Form oder Offenbarung, sondern allein des Vaters, der in solcher unterschiedlichen Form hat uns wollen bekannt werden als eine unterschiedliche Person vom Sohne und Heiligen Geist, in einem unzertrennten göttlichen Wesen.

78. Nimm dir auch aus der Grammatica ein groß Exempel. Wenn der Priester tauft, oder absolvirt, und spricht: „Im Namen des Vaters, und des Sohnes, und des Heiligen Geistes.“ Diese Worte allesamt sind Gottes Geschöpf und Werk in unserm Munde (so wohl als wir selbst, und was wir haben), und ist keines unterschiedlich des Vaters allein, oder des Sohnes allein, oder des Heiligen Geistes allein, sondern aller drei Personen, des einigen GOttes einerlei Geschöpf. Aber nach der Deutung oder Offenbarung mußt du nicht sagen, daß dies Wort „des Vaters“ bedeute alle drei Personen, sondern unterschiedlich allein den Vater; das Wort „des Sohnes“ unterschiedlich allein den Sohn; das Wort „des Heiligen Geistes“ unterschiedlich allein den Heiligen Geist, in einer einigen Gottheit, die uns durch solche Worte oder Deutung offenbart wird, daß drei unterschiedliche Personen sind in der einigen Gottheit. Denn er nicht spricht: In den Namen, als vieler, oder als hätte eine jegliche Person einen sonderlichen Namen und Wesen; sondern „im Namen“ (spricht er) als in Einem Namen Eines Wesens, und doch drei unterschiedlicher Personen.

79. Also siehst du, daß die Creatur [auf] zweierlei Weise anzusehen ist, ut res et signum, daß sie etwas für sich selbst ist, von Gott geschaffen
\quellseite{970}
, und auch gebraucht wird, etwas Anderes zu zeigen oder zu lehren, das sie selbst nicht ist. Der Rauch ist eine res, ein Ding für sich selbst, und doch auch ein Zeichen eines andern Dinges, das er nicht ist, sondern zeigt und offenbart dasselbe, nämlich das Feuer. Davon schreibt St. Augustinus viel de doctrina christiana; aber hier ist es, in dieser hohen Sache, etwas mehr. Denn die Menschheit Christi ist nicht ein schlecht Zeichen, oder ledige Gestalt, gleichwie die Taube auch nicht eine ledige Gestalt, und die Stimme nicht eine ledige Gestalt oder Bild; sondern die Menschheit, darin GOttes Sohn unterschiedlich offenbar wird, ist voll, und mit GOtt in Eine Person vereinigt, die ewig sitzen wird zur Rechten GOttes in seinem Reiche, wie droben David verheißen ist, 1 Chron. 18, 12. Die Taube ist eine Gestalt, vom Heiligen Geist eine Zeitlang angenommen, sich zu offenbaren; nicht mit ihm in Eine Person vereinigt ewiglich, sondern wieder verlassen; wie die Engel Menschengestalt annehmen, darinnen erscheinen, und wieder verlassen. Also ist es auch gethan mit der Stimme des Vaters. Denn da ist keine Verheißung, daß [es] so sollte ewiglich bleiben, sondern ist eine zeitliche Offenbarung.

80. Wenn wir nun im Kinderglauben sprechen: „Ich glaube an GOtt den Vater, allmächtigen Schöpfer Himmels und der Erde“, ist nicht die Meinung, daß allein die Person des Vaters sollte allmächtig, Schöpfer und Vater sein, sondern, der Sohn ist ebenso wohl allmächtig, Schöpfer und Vater; der Heilige Geist auch so wohl allmächtig, Schöpfer und Vater; und doch nicht drei Allmächtige, Schöpfer, Väter, sondern ein einiger Allmächtiger, Schöpfer, Vater Himmels und der Erde, und unser aller, gleichwie der Vater unser Heiland und Erlöser, der Sohn unser Heiland und Erlöser, der Heilige Geist unser Heiland und Erlöser, und doch nicht drei Heilande noch Erlöser sind, sondern ein einiger Heiland und Erlöser ist. Gleichwie der Vater unser GOtt, der Sohn unser GOtt, der Heilige Geist unser GOtt, doch nicht drei Götter, sondern ein einiger GOtt ist: also heiligt der Heilige Geist die Christenheit, der Vater auch, der Sohn auch, und sind doch nicht drei Heiliger oder Heiligmacher, sondern ein einiger Heiligmacher etc. Opera Trinitatis ad extra sunt indivisa.

81. Es ist aber alles geredet darum, daß wir unterschiedlich drei Personen in der einigen Gottheit glauben und erkennen, und ja nicht die Personen mengen, noch das Wesen trennen. Der Unterschied des Vaters (wie [§ 68] gehört) ist, daß er die Gottheit von niemand hat, sondern sie von Ewigkeit, durch die ewige Geburt, dem Sohne gegeben hat. Darum ist der Sohn GOtt und Schöpfer, gleich dem Vater; aber das hat er alles vom Vater, nicht wiederum, der Vater vom Sohne. Denn, daß der Vater GOtt und Schöpfer ist, das hat er nicht vom Sohne; sondern, daß der Sohn GOtt und Schöpfer ist, das hat er vom Vater. Also hat der Vater oder Sohn vom Heiligen Geiste nicht, daß er GOtt und Schöpfer ist; sondern, daß der Heilige Geist GOtt und Schöpfer ist, das hat er vom Vater und Sohn. Also steht nun das Wort „GOtt allmächtig, Schöpfer“, bei dem Vater, und nicht bei dem Sohne und Heiligen Geiste, zu merken den Unterschied des Vaters vom Sohne und Heiligen Geiste in der Gottheit. Wiederum, der Unterschied des Sohnes vom Vater und Heiligen Geiste,\footnote{Die Worte: „in der Gottheit. Wiederum der Unterschied des Sohnes vom Vater und Heiligen Geist“ fehlen in der Originalausgabe.} des Heiligen Geistes vom Vater und Sohne; nämlich, daß der Vater ist der Ursprung oder Quelle (so man es so nennen sollte, wie die Väter thun) der Gottheit, von welchem sie der Sohn hat, und der Heilige Geist vom Sohne und Vater in Ewigkeit; und nicht wiederum.

82. Ueber solchen innerlichen der Personen Unterschied ist nun der äußerliche Unterschied, da der Sohn und Heiliger Geist innen offenbar ist. Der Sohn in der Menschheit; denn der Sohn ist allein Mensch worden, vom Heiligen Geist empfangen, von Maria der Jungfrau geboren, für uns gelitten, gestorben etc., wie der Glaube weiter lehrt. Doch, daß gleichwohl recht heißt: GOtt ist für uns gestorben. Denn der Sohn ist GOtt, und ist kein anderer GOtt mehr, sondern mehr Personen in derselben Gottheit. Der Heilige Geist ist allein unterschiedlich offenbar in den feurigen Zungen, Gaben, mancherlei Sprachen und Wunderzeichen etc., obwohl die Menschheit von allen dreien Personen gemacht, und die feurigen Zungen, die Gaben des Heiligen Geistes, aller dreier Personen
\quellseite{971}
Geschöpf und Werk sind; wie nun genugsam auf diesemal gesagt ist. Man hat hiervon köstliche Bücher, St. Augustini, Hilarii, Cyrilli, und ist solcher Artikel im Pabstthum und bei den Schultheologen rein geblieben, daß wir mit ihnen darüber keinen Zank haben.

83. Hier bekümmern sich etliche, ob sie die Person des Vaters nennen, wenn sie beten „Vater unser“, oder das göttliche Wesen? Wunder ist es nicht, daß einem Menschen in diesem über und über wunderlichen, unbegreiflichen Artikel wunderliche Gedanken einfallen, deren zuweilen einer mißlingt, oder ein Wort mißbräth. Aber wo der Grund des Glaubens fest bleibt, werden uns solche Splitter, Spänlein oder Strohhalme nicht schaden. Der Grund aber des Glaubens (wie [§ 66 ff.] gehört) ist, daß du glaubest, es sind drei Personen in der einigen Gottheit, und eine jegliche Person ist derselbige einige vollkommene GOtt; daß also die Personen nicht gemenget, das Wesen nicht zertrennt werde, sondern Unterschied der Personen, und Einigkeit des Wesens bleibe. Denn das ist es, daß sich die Engel in Ewigkeit nicht satt können (wie St. Petrus [1. Ep. 1, 12.] sagt) sehen und wundern, und darüber ewig selig sind; und wo sie es zu Ende ansehen könnten, würde ihre Seligkeit auch aus sein und ein Ende haben; wie wir auch solches sehen werden, und dadurch ewig selig sein, wie der HErr spricht Joh. 17, 3.: „Das ist das ewige Leben, daß sie dich, und den du gesandt hast, Jesum Christum, erkennen.“ Indes muß der Glaube sich am Worte halten; Vernunft kann nichts hier thun, denn sprechen, es sei unmöglich, und wider sich selbst, daß drei Personen, eine jegliche vollkommener GOtt, und doch nicht mehr denn ein einiger GOtt sei, und allein der Sohn Mensch sei; wer aber den Vater und Sohn hat, dem wird der Heilige Geist wohl bekannt werden, vom Vater und Sohne.

84. So hast du droben [§§ 67. 80] gehört, daß der Vater ist unser aller GOtt und Vater, der Sohn ist unser aller GOtt und Vater, der Heilige Geist ist unser aller GOtt und Vater; und ist doch nicht mehr denn ein einiger GOtt unser Vater. Denn das Wesen ist unzertrennt; darum, welche Person du nennst, so hast du den rechten einigen GOtt in dreien Personen genannt, weil eine jegliche Person derselbige einige vollkömmliche GOtt ist, und kannst hierin nicht irren noch fehlen. Denn Jesus Christus ist kein anderer GOtt oder Vater oder Schöpfer, denn der Vater und Heiliger Geist ist, ob er gleich eine andere Person ist. Eben so ist der Vater und Heiliger Geist auch.

85. Demnach ist es nicht allein falsch, sondern auch unmöglich und nichtig, daß du die Person des Vaters, als die unterschiedliche Person, wolltest Vater nennen, und nicht den Sohn und Heiligen Geist zugleich mit Vater nennen; denn das hieße das göttliche Wesen zertrennen, und den Sohn und Heiligen Geist ausgeschlossen. Das ist nichts. Denn nach solcher Weise der persönlichen Vaterschaft hat der Vater nicht mehr denn Einen Sohn, und der Sohn nicht mehr denn Einen Vater. Solcher Vater ist er dir nicht, und du nicht solcher sein Sohn; sondern das ist der einige Sohn vom Vater in Ewigkeit, wie der 2. Psalm, V. 7., spricht: „Der HErr sprach zu mir: Du bist mein Sohn, heute habe ich dich geboren“; aber du bist, nach deinem Alter, dreißig, vierzig, fünfzig Jahr, so lange du geschaffen und getauft bist gewesen, ein zeitlicher Sohn aller drei Personen, Eines GOttes. Quia opera Trinitatis ad extra sunt indivisa; sic cultus Trinitatis ad extra est indivisus; was GOtt gegen die Creatur thut, das thun alle drei Personen ohne Unterschied.

86. Denn es ist ein einig göttlich Wesen aller dreier Personen, und was wir oder die Creatur gegen eine jegliche Person thun, das thun wir gegen dem einigen GOtt, und allen dreien Personen ohne Unterschied. Denn er ist gegen uns ein einiger GOtt, und in sich selbst drei Personen unterschiedlich; wie der HErr Christus selbst spricht Joh. 14, 9. 10.: „Philippe, wer mich siehet, der siehet den Vater; wie sprichst du denn: Zeige uns den Vater? Glaubst du nicht, daß ich im Vater bin, und der Vater in mir?“ Joh. 5, 23.: „Sie sollen den Sohn ehren, gleich wie den Vater.“ Joh. 10, 30. 31.: „Ich und der Vater sind Eins“; das sprechen wir: Ein Ding, Ein Wesen, Ein GOtt, Ein HErr. Hier huben die Juden Steine auf, und wollten ihn steinigen. Joh. 5, 17.: „Mein Vater wirket, oder schafft bisher, und ich wirke auch“; darum trachten die Juden vielmehr ihn zu tödten, daß er nicht allein den Sabbath brach, sondern sagte auch, GOtt sei sein Vater, und machte sich selbst Gott gleich etc., Joh. 10, 33. Davon will ich jetzt aufhören; denn ich hätte Willen ein Buch zu schreiben, so bin ich ins Predigen kommen.

\quellseite{972}
Lies das Evangelium Johannis, das lehrt uns solches alles reichlich.

87. Nun wir haben Mosen, daß er mit St. Johanne stimmt, es sei ein Wort im Anfang gewesen, durch welches alles gemacht ist, und daß solch Wort nicht könne eine Creatur oder gemacht sein, und doch etwas Anderes, oder eine andere Person sein denn GOtt, deß solch Wort ist. Denn, weil es nicht gemacht ist, sondern alle Dinge durch dasselbe gemacht sind, muß es GOtt, Schöpfer sein aller Creatur, nachdem es gewiß ist, daß außer der Creatur, die gemacht ist, nichts sein kann denn GOtt, der sie macht; und doch das Wort, der GOtt und Schöpfer, durch den alles gemacht ist, ein Anderes ist von dem Sprecher, oder der solch Wort spricht. Hiermit ist nun Moses unser Zeuge und ein Christ worden, lehrt eben, daß wir Christen lehren, nämlich, daß GOtt ein Wort habe im Anfang, durch welches alles gemacht ist, gleichwie Johannes schreibt.

88. Nun laßt uns den andern Legaten, Sanct Paulum, auch kürzlich hören, wie er Mosen grüßt und ruft. Col. 1, 15. ff. redet er von unserm HErrn Jesu Christo also: „Welcher ist das Ebenbild des unsichtbaren GOttes, der Erstgeborne vor allen Creaturen. Denn durch ihn ist alles geschaffen, das im Himmel und auf Erden ist, das Sichtbare und Unsichtbare; beide die Thronen und Herrschaften, und Fürstenthümer und Obrigkeiten; es ist alles durch ihn, und zu ihm geschaffen, und er ist vor allen, und es bestehet alles in ihm.“ Diese Worte können nicht von Christo nach der Menschheit geredet sein, das ist gewiß; denn er ist nicht Mensch gewesen vor allen Creaturen, sondern sind heuer 1543 Jahr, daß er Mensch worden ist. Und ist fürwahr ein gewaltiger, klarer Spruch, daß Christus ein ewiger GOtt, Schöpfer Himmels und Erden sei, und noch heutiges Tages und immerfort alles durch ihn bestehe, erhalten oder gemacht werde, auch alles, was doch ist im Himmel und Erden, Engel und Geister, sichtbar und unsichtbar. Hierin stimmt er mit Johanne gleich, da er Cap. 1, 3. spricht: „Alles ist durch ihn gemacht, und ohne ihn ist nichts gemacht, was gemacht ist.“ Hört nun Mose und bekennt die Worte Johannis, so hört und bekennt er gewißlich auch diese Worte Pauli, und spricht: Ja, mein lieber Paule, eben wie du und Johannes sagst, so habe ich auch geschrieben, daß alle Dinge sind durchs Wort geschaffen, 1 Mos. 1, 3. ff.

89. Weiter spricht Paulus 1 Cor. 10, 4.: „Sie tranken von dem geistlichen Fels, der mit ihnen zog, der Fels aber war Christus.“ Ist Christus zur selbigen Zeit gewesen, der mit den Kindern Israel gezogen ist, und von dem sie geistlich getrunken und geistlich getauft sind, das ist, mit uns einerlei Glauben an den zukünftigen Christum, der uns nun erschienen ist, gehabt, so muß Christus rechter ewiger GOtt sein. Denn an die Engel kann man nicht glauben, welches GOtt allein gebührt, auch sie nicht unsere geistliche Speise sein können, GOtt muß [es] selber sein. Item, daselbst 1 Cor. 10. spricht er, V. 9.: „Laßt uns Christum nicht versuchen, wie jener etliche [ihn] versuchten, und wurden von den Schlangen umgebracht.“ Was will hier werden? Schreibt doch Mose allenthalben, es sei der HErr, Jehovah, der rechte einige GOtt, den die Kinder Israel versuchten, 2 Mos. 17, 2.: „Warum versuchet ihr den HErrn?“ 4 Mos. 14, 22. spricht der HErr: „Sie haben mich wohl zehmal versucht.“ Ist es der HErr, wie Mose schreibt, wie kann es Christus sein, wie Paulus schreibt? Nun müssen sie beide recht schreiben; denn der Heilige Geist ist nicht wider sich selbst.

90. Hieraus folgt gewaltiglich und unwidersprechlich, daß der GOtt, der das Volk Israel aus Egypten und durchs rothe Meer geführt, in der Wüste durch die Wolkensäule und Feuersäule geleitet, mit Himmelbrod genährt, und alle die Wunder gethan, so Mose in seinen Büchern beschreibt; item, der sie ins Land Canaan gebracht, und darinnen Könige und Priesterthum, und alles gegeben hat, sei eben der GOtt, und kein anderer denn Jesus von Nazareth, Marien der Jungfrau Sohn, den wir Christen unsern GOtt und HErrn nennen, den die Juden gekreuzigt haben, und noch heutiges Tages lästern und fluchen, wie Jesaia Cap. 8, 21. sagt: „Sie werden in ihrer Angst fluchen ihrem Könige und ihrem GOtt.“ Item, er ist es, der auf dem Berge Sinai Mosi die zehn Gebote gibt, und spricht 2 Mos. 20, 2. 3.: „Ich, der HErr, bin dein GOtt, der dich aus Egypten geführet hat, du sollst vor mir keine andere Götter haben.“ Ja, Jesus Nazarenus, am Kreuz für uns gestorben, ist der GOtt, der in dem ersten Gebote spricht: „Ich, der HErr, bin dein GOtt.“ Wenn solches die Juden und Mahomet hören sollten, wie sollten sie toben! Dennoch ist es wahr, und
\quellseite{973}
muß wahr bleiben ewiglich, und soll ewiglich dafür zittern und brennen, wer es nicht glaubt.

91. Denn da steht Mose klar, und spricht, daß durch den Spruch, oder durch das Wort GOttes, sei alles geschaffen; und David Ps. 33, 6.: „Der Himmel ist durch GOttes Wort gemacht.“ Ist der Himmel durchs Wort gemacht, so ist auch alle andere Creatur dadurch gemacht; denn wer Eine Creatur macht, der macht sie alle; wer sie nicht alle macht, der wird keine machen können. Und stimmt also Mose und David mit Johanne und Paulo, die auch beide mit ihnen gleich sagen: „Alle Dinge sind durch das Wort“, oder durch „Christum, geschaffen und gemacht.“ Ist nun alles durch ihn gemacht, und ohne ihn ist nichts gemacht, wie der Text aller viere, Mosis, Davids, Johannis, und Pauli, da steht; so muß unter dem, daß sie „Alles“ heißen, begriffen, und nicht ausgeschlossen sein die Ausführung aus Egypten, und was mehr in dem Volke Israel geschehen ist; ja, alles, was allenthalben von Anfang der Creaturen geschehen ist, noch immer geschiehet, und hinfort geschehen wird. Denn es sind große, wichtige Worte, da sie sagen: „Alles ist durch ihn gemacht.“ Und wie Moses redet: „GOtt sprach: Es werde; und es ward also.“ Ob nun Mose nicht nennt den Sohn, oder Christum, nach der Grammatica, so nennt und bekennt er doch den Spruch, oder Wort, durch welches alles gemacht ist; damit er anzeigt, daß in GOtt ein anderer ist, der da spricht, und ein anderer, der das Wort ist, und doch ein einiger Schöpfer aller Creaturen ist. Denn es mußte dem Neuen Testamente auch etwas vorbehalten sein, darinnen der Vater und der Sohn und der Heilige Geist klarlicher genannt würde, welche\footnote{sc. Dreifaltigkeit. Cruciger: qui, bezogen auf das vorher von ihm eingefügte Deus.} im Alten GOtt der Sprecher, das Wort, und der Geist des HErrn genannt ist.

92. Darum hilft die Juden, Türken, Ketzer nichts, daß sie sehr große Andacht vorgeben, und rühmen wider uns Christen, wie sie glauben an den einigen GOtt, Schöpfer Himmels und der Erde, nennen ihn auch Vater mit großem Ernst, und ist doch nichts denn eitel vergebliche, unnütze Worte, darin sie den Namen GOttes unnützlich führen und mißbrauchen wider das andere Gebot; wie Christus spricht Joh. 8, 54. zu den Juden: „Es ist mein Vater, der ehret mich, welchen ihr sprechet, er sei euer GOtt, und kennet ihn nicht.“ Fürwahr, das reimt sich sehr übel, den Vater GOtt nennen, und nicht wissen, wer er ist. Denn so du solltest einen solchen großen heiligen\footnote{Cruciger: valde sanctus Judaeus etc. In den deutschen Ausgaben: „solchen großen Heiligen, Juden“ etc.} Juden, Türken, Ketzer, fragen: ob er auch glaube,\footnote{In den deutschen Ausgaben: „glaubet“; Cruciger: credatne.} daß derselbige einige GOtt, Schöpfer Himmels und der Erde (des Namen sie so andächtig rühmen und ihn Vater nennen, wiewohl alles fälschlich), auch ein Vater sei, und einen Sohn habe außer der Creatur in der Gottheit? so würde er vor großer Heiligkeit erschrecken, und solches für eine greuliche Lästerung halten. Fragst du weiter: ob derselbige einige GOtt, Schöpfer, Vater (den sie also nennen mit ihrem Lügenmaul), auch ein Sohn sei, und einen Vater habe in der Gottheit? da würde er vor großer Andacht die Ohren zustopfen, die Zähne zubeißen,\footnote{Cruciger: dentibus infrendet = er würde mit den Zähnen knirschen.} und sorgen, die Erde möchte dich und ihn verschlingen. Fragst du weiter: ob derselbige einige GOtt, Schöpfer, Vater (wie sie rühmen), auch ein Heiliger Geist sei, und habe den Vater und Sohn, von welchen er sein göttlich Wesen habe? da würde der allerheiligste Mann von dir laufen, als wärest du der ärgste Teufel aus der Hölle heraus.

93. Hier siehst du, daß sie nicht wissen, was GOtt ist; und wenn sie ihn nennen: GOtt, Schöpfer, Vater, wissen sie nicht, was sie sagen. Denn wo GOtt nicht soll sein ein solcher GOtt (wie uns die Schrift lehrt), der ein natürlicher Vater ist, einen natürlichen Sohn, und beide einen natürlichen Heiligen Geist haben, in dem einigen göttlichen Wesen, da ist GOtt nichts, und gar kein GOtt. Darum haben sie keinen GOtt, ohne daß sie GOttes Namen mit Sünden und Schanden mißbrauchen, und erdichten ihnen einen eigenen GOtt und Schöpfer, der ihr Vater, und sie seine Kinder sein sollen; nehmen ihm seine natürliche Vaterschaft, seinen natürlichen einigen Sohn, und den natürlichen Heiligen Geist, das ist, die ganze rechte Gottheit, und geben ihm dafür ihren nichtigen Traum und Lügen von GOtt, Schöpfer, Vater; ja, solchen heiligen Namen GOttes geben sie ihrem Traum und Lügen, das ist, dem Teufel, derselbe ist ihr Gott und Vater, ein Vater aller Lügen; wollen
\quellseite{974}
gleichwohl die liebsten Kindlein und größesten Heiligen sein.

94. Denn es ist beschlossen, und also hat sich GOtt selbst uns offenbart, daß er sei ein einiger GOtt, Schöpfer und Vater Himmels und der Erde; und derselbige einige GOtt, Schöpfer und Vater aller Welt, sei ein natürlicher Vater eines einigen Sohnes in der Gottheit; und derselbige einige GOtt, Schöpfer und Vater aller Welt, sei ein einiger natürlicher Sohn des Vaters in der Gottheit; und derselbige einige GOtt, Schöpfer und Vater aller Welt, sei ein Heiliger Geist, vom Vater und Sohne in der Gottheit. Denn die drei unterschiedlichen Personen sind ein einiger GOtt, Schöpfer und Vater aller Welt; und eine jegliche Person ist derselbige völlige einige GOtt, Schöpfer und Vater aller Welt. Und wenn du JEsum Christum anrufst, und sprichst: O mein lieber HErr GOtt, mein Schöpfer und Vater, JEsu Christe, du einiger ewiger GOtt! darfst du nicht sorgen,\footnote{Walch und die Erlanger: „sagen“.} daß der Vater und Heiliger Geist darum zürne; sondern [sie] erkennen,\footnote{Wir haben mit Cruciger „Vater, Sohn und Heiliger Geist“ als Subject zu „erkennen“ angenommen. Er übersetzt: et totam divinitatem agnoscere et velut attestari.} daß, welche Person du anrufest, gleich alle drei Personen, und den einigen GOtt anrufest; denn du kannst keine Person ohne die andern anrufen, sintemal da ist ein einiges unzertrennetes göttliches Wesen in allen, und in einer jeglichen Person. Wiederum, kannst du keine Person insonderheit verleugnen, es sind alle drei und der einige GOtt ganz und gar verleugnet; wie 1 Joh. 2, 23. sagt: „Wer den Sohn verleugnet, der hat auch den Vater nicht.“

95. Ja, sage ich, unrecht ist es nicht, sondern wohl gethan, wenn du JEsum Christum also anrufst, gleichwie die Kirche singt auch vom Heiligen Geist: Veni Pater pauperum, komm, du Vater der Elenden. Doch ist es feiner, daß man die Ordnung der Personen halte, und nicht verachte, wie die Apostel thun, und die Kirche nach ihrem Exempel thut, da sie die Person des Vaters im Anrufen oder Gebet nennen, wie im Vater-Unser c. Denn er ist der Ursprung oder Brunn (wie man es nennen kann) der Gottheit im Sohne und Heiligen Geist, und kann der Sohn (wenn der Vater genannt ist) nicht abgesondert, sondern muß zugleich mit genannt und gemeint sein. Also auch, der Heilige Geist muß mit dem Vater und Sohne genannt und gemeint sein, weil keine Person außer der andern ein sonderlicher GOtt sein kann. Also redet St. Paulus [2 Cor. 1, 3.] und St. Petrus [1. Ep. 1, 3.]: „Gelobet sei GOtt, der Vater unseres HErrn JEsu Christi, der Vater aller Gnaden“; und er selbst im Evangelio immer den Vater vorzeucht, und ihm alles zuschreibt, und doch sagt, Joh. 5, 23., sie sollen den Sohn ehren, wie den Vater; item, Joh. 16, 15.: „Alles, was der Vater hat, ist mein“; ohne daß der Vater die erste Person ist, von dem es der Sohn hat; und nicht wiederum. Daß aber etliche Sünde unterschiedlich wider den Vater, wider den Sohn, und in den Heiligen Geist geschehen, gehört zur Offenbarung der Personen, nicht zur Trennung des Wesens, davon droben [§ 66 ff.] ein wenig, und anderswo das und Mehreres genugsam gehandelt ist.

96. Wie wollen wir aber hier thun, da St. Johannes weiter schreibt von dem Worte, und spricht: „Das Wort ist Fleisch worden.“ Das wird sich ja mit dem Worte nicht reimen, davon Mose schreibt: „GOtt sprach: Es werde Licht“; oder da David von sagt: „Der Himmel ist durchs Wort des HErrn gemacht.“ So gebeut Moses, oder das Wort selbst (wie wir glauben), auf dem Berge Sinai, 2 Mos. 20, 4., man solle ihm kein Bild noch Gleichniß machen, weder im Himmel noch auf Erden; und Johannes macht nicht ein Bild, sondern eine Creatur und Menschen daraus, und spricht [Joh. 1, 14.]: „Das Wort ist Mensch worden.“ Dasselbe thut Paulus auch, und spricht [Röm. 1, 3. Gal. 4, 4.]: „er sei Davids Sohn, oder Saame, von einem Weibsbilde geboren.“ Darum muß Mose von einem andern Worte reden, durch welches alles geschaffen ist. Durch einen Menschen, der selbst ein Geschöpf ist, kann nichts geschaffen werden. Auch ist Paulus und Johannes wider sich selbst, daß sie ihn zum Menschen machen, und doch sagen: es sei alles durch ihn geschaffen.

97. Lasset uns versuchen, ob Mose wolle sich finden lassen, daß er auch solches sage. 1 Mos. 3, 15. schreibt er also, daß GOtt zur Schlange spricht: „Ich will Feindschaft setzen zwischen dir und dem Weibe, zwischen deinem Samen und ihrem Samen, derselbe soll dir den Kopf zertreten, und du wirst ihm in die Ferse stechen.“ Offenbar ist es, daß GOtt hier nicht redet von gemeiner Schlange, die im Grase oder Häuser
\quellseite{975}
läuft und junge Frösche frißt, sondern von der Schlange, die zu der Zeit ein schön Thier gewesen, und einen hohen Geist bei sich gehabt, der nicht allein hat reden können, sondern auch von hohen, göttlichen Sachen und Geboten disputiren, als hätte er es im Himmel gelernt; welches keiner Creatur, ohne den Engeln und Menschen, gegeben ist; und thut dadurch solchen Schaden, daß er den Menschen fället in die Sünde und in den ewigen Tod, durch den schönsten Schein des göttlichen Namens. Das ist nicht ein schlecht, alber Schlänglein, das Fröschlein frißt, sondern die ganze Welt verschlingt. Es ist der leidige Teufel in der Schlange, der den Tod in die Welt durch die Sünde gebracht hat. Von diesem Todtschläger und Sündenmeister und Weltfresser redet GOtt, daß ihm soll sein Kopf zertreten werden, das ist, seine Macht, der Tod und Sünde sollen zerstört, das Leben und Gerechtigkeit sollen wiedergebracht werden.

98. Und das soll thun des Weibes Same. Und wie er durch ein Weib, das vom Manne, ohne Weib, gekommen ist, den Menschen gefället hat, so soll ihn der Same, der vom Weibe, ohne Mann, kommen wird, wiederum fällen. Dieser Same des Weibes wird müssen ein Mensch und Sohn sein; denn in der Schrift heißt „Menschen Same“ die Frucht, so ein Sohn eines Menschen ist (wie man weiß), und ist das gar ein Sonderliches an diesem Orte, daß dies Kind oder Mensch eines „Weibes Same“ heißt. Sonst heißt allenthalben „Same“ des Mannes oder Vaters Same, als, Abrahams Same, Davids Same, und so fort, durch und durch, im Mose und Propheten heißt „Same“ des Mannes Same; daß Moise an diesem Ort eben lautet wie Lucas und Matthäus, daß dies Weib soll eine Jungfrau sein, die ohne Mann eine Mutter werden solle ihres eigenen allein selbst Samen. Und weil solches sich mit dem Neuen Testament reimt, sollen wir Christen, nach vorgenommener Regel, weder Juden noch Teufel keinen andern Verstand gestatten.

99. Summa, dieser Weibes Same soll ein Mensch sein, das ist gewiß; darüber muß er wahrlich auch GOtt sein, oder Mose wird ein abgöttischer Teufelsprophet sein. Denn er gibt diesem Samen die Macht, die allein GOtt, und keiner Creatur gebühret, nämlich, daß er solle den Tod und Todtschläger, Sünde und GOttes Zorn wegthun, Gerechtigkeit und Leben wiederbringen. Lieber, das wird kein Engel, noch alle Engel zusammen nicht thun; es muß ein höherer, mächtigerer Mann sein, denn alle Engel und alle Creaturen sind. Ein verdammter abgöttischer Prophet (sage ich noch eins) muß Moses sein, so er solche Werke, den Tod und Sünde erwürgen und überwinden,\footnote{Erlanger: überwunden.} lebendig und gerecht machen, einem Weibes Samen gibt, der eine bloße Creatur, und nicht selbst der einige GOtt ist, der allein lebendig machen kann,\footnote{Im Original fehlt: kann.} wie Joh. 1, 4. vom Worte sagt: „In ihm war das Leben.“

100. Es muß ja auch die Vernunft selbst bekennen, daß, wer den Tod kann untertreten, der kann auch das Leben wiedergeben; wer Sünde kann untertreten, der kann Gerechtigkeit wiedergeben; sintemal Tod wegnehmen nichts anders ist, denn das Leben wiedergeben; Sünde wegnehmen nichts anders ist, denn Gerechtigkeit wiedergeben, davon die Schlange (der Teufel in der Schlange) Adam und Eva, sammt allen ihren Nachkommen und Menschenkindern, böslich gebracht hat, und durch seine Lügen in die Sünde und Tod gefället hat; wie der Text klar da steht 1 Mos. 2, 17.: „Du sollst nicht essen vom Baume, oder du mußt sterben.“ Dawider sagt der Lügner und Mörder 1 Mos. 3, 4.: Ihr mögt wohl davon essen, und werdet nicht sterben, sondern GOtt gleich werden, und alles wissen. Das ist (wie droben [§ 97] gesagt) alles von Sünden und Tod geredet, so die Schlange gestiftet und angerichtet hat. Darum muß das Untertreten der Schlange nichts anders sein, denn sein Werk und Macht zerstören; wie Paulus redet, daß Christus den Tod zerstört hat, und das Leben ans Licht gebracht, 2 Tim. 1, 10. Was Juden, Mahomet und andere hier sudeln mit ihren Glossen, achten wir nichts. Mose stimmt hier mit dem Neuen Testament, das ist uns genug.

101. Solchen Verstand, daß des Weibes Same müsse GOtt sein, der dem Teufel den Kopf zertreten soll, hat auch Adam und Eva gehabt. Denn 1 Mos. 4, 1., da Eva Cain geboren hatte, dachte sie, vielleicht weil\footnote{In den Ausgaben ist so interpungirt, wie wir dafürhalten, falsch: „dachte sie vielleicht, weil“, denn Luther sagt: „Solchen Verstand hat Adam und Eva gehabt“, nicht: vielleicht gehabt. Im Lateinischen ist „vielleicht“ nicht ausgedrückt.} das der

i
\quellseite{976}
erstgeborne Mensch auf Erden wäre, er würde der beste sein, und meinte, er sollte der Same des Weibes sein, sie aber die Mutter, oder das Weib; darum spricht sie: „Ich habe den Mann, den HErrn“; als wollte sie sagen: Das wird der Mann, der HErr sein, davon GOtt geredet hat, des Weibes Same etc. Nennt das Kind einen Mann und HErrn, oder GOtt; denn hier steht der große und eigene Name GOttes, Jehovah, der nichts Anderes bedeutet denn allein GOtt selbst in seiner Natur oder Wesen; und was, welches, wo es allein steht, ohne ein Weib, heißt es nicht schlecht ein Mannsbild, wie alle Männer sind; sondern einen Ausbund und vornehmlichen Mann; wie wir Deutschen auch sagen: Das ist ein Mann! Das will ein Mann werden! Also will Eva hier: Ich habe einen Sohn geboren, der wird ein Mann werden, ja, er ist der Mann, GOtt selbst, der es thun soll, und die Schlange zertreten, wie GOtt uns geredet hat. Wie ist es möglich? wie sollt es ihr einfallen, von ihrem Kinde also zu reden: Ich habe den Mann, den HErrn, wo sie nicht den Spruch also verstanden hätte, daß des Weibes Same müßte GOtt sein, der es thun sollte, was GOtt geredet hatte?

102. Sie wird es auch nicht allein so verstanden haben, sondern Adam wird mit ihr lange zuvor davon geredet, und diesen Spruch untereinander wohl geübt, und sich getröstet haben wider die Sünde und Tod, so durch diesen Samen sollte weggenommen, dafür die verlorne Unschuld und Leben wiedergebracht werden; sonst wären sie verzweifelt. So ist auch GOttes Wille noch Weise nicht, daß sein äußerlich Wort (wie dies ist) solle vergeblich geredet, und von niemand verstanden werden; wie er spricht Jes. 55, 11.: „Mein Wort soll nicht leer wieder zu mir kommen, sondern ausrichten, dazu ich es sende. Nun waren hier allein zwei Menschen, Adam und Eva, die es verstehen konnten, darum haben sie es müssen fruchtbarlich, seliglich und recht verstanden haben, allerdings wie mir es Christen, und zuvor alle Propheten verstanden haben.

103. Darin fehlt die liebe elende Mutter Eva, daß sie wähnt, sie sei das Weib, weil kein ander Weib auf Erden ist, denn sie allein; und vor großer Begierde und Verlangen hofft sie, ihr Sohn solle der Same, der Mann Jehova sein. Das ist zu früh und zu sehr geeilt; aber [sie] ist nicht darum zu verdenken, daß sie der Sünden und Todes, das ist, des Teufels gerne halb los wäre. Etc. aber hatte nicht zu ihr gesagt: Dein Same soll es thun; auch nicht zu Adam: Deines Weibes Same soll es sein, sondern ließ allen beiden ihren guten Text, den alle Menschenkinder noch fühlen, bis an der Welt Ende. Zur Schlange aber kehrt er sich, und spricht: Ich will dir wiederum einen Kopftreter schaffen, der soll eines Weibes Same sein; will dich hochmüthigen, mächtigen, bösen Geist durch eines Menschen Sohn fällen, daß wiederum alle Menschen über dich laufen und mit Füßen treten sollen, wie du jetzt Adam und Eva unter dich getreten hast. Das hat gethan, thut noch, und will es immerfort thun, unser lieber HErr Christus, mit GOtt dem Vater Ein Jehova, Amen.

104. Hier möchte jemand sagen: Wie geht es zu, daß solches kein Christ noch Jude an diesem Orte gesehen hat. Denn die Dolmetscher alle machen es anders. Der lateinische also: „Ich habe einen Menschen bekommen durch GOtt.“\footnote{So die alten Rationalisten (z. B. Gesenius) und etliche neuere lutherische Theologen, z. B. Köstlin, Martin Luther (3. Aufl.), Bd. II, S. 602, welcher Luthers Uebersetzung und Auslegung eine „exegetische Verirrung“ nennt.} Die andern Hebräisten also: „Ich habe den Mann kriegt von dem HErrn. Da frage ich jetzt nicht nach. Ich habe droben [Vorrede § 9] oft bedingt, ich wolle diesmal keinen Meister haben, sondern meine Meinung in Dolmetschung anzeigen. Gefällt es niemand, so ist es genug, daß [es] doch mir allein gefällt. Das hebräische Wörtlein \hebraeisch{את} heißt „den“ oder „die“, und ist ein Artikel Accusativi, wie das alle Grammatici bekennen müssen. Als, da Mose Cap. 1, 1. spricht: „Im Anfang schuf GOtt, \hebraeisch{את} Himmel, und \hebraeisch{את} Erde“: das heißt deutlich: den Himmel und die Erde, und immer so fort, in dem und folgenden Capitel. Als „Adam erkannte \hebraeisch{את} Heva, sein Weib“: „Heva gebar \hebraeisch{את} Kain.“ Item, weiter gebar sie \hebraeisch{את} Habel, seinen Bruder. Item, Adam zeugete \hebraeisch{את} Seth; Seth zeugete \hebraeisch{את} Enos; und so fort an. Eben der Weise nach spricht hier Eva, da sie Kain geboren hatte: \hebraeisch{קניתי איש את־יהוה}: Ich habe den Mann kriegt, den HErrn. Denn sie höfft (wie gesagt), Kain solle der Same sein, der von GOtt verheißen war, der Schlange den Kopf zu zertreten.

105. Und ich weiß fürwahr, wenn die ärg
\quellseite{977}
Juden, die Christum gekreuziget haben, oder noch ärger wären, als die, so ihn noch jetzt gerne viel greulicher kreuzigen wollten; wie man sagt von denen, so in Ungarn zu Ofen neulich sammt den Fürsten eine Katze gekreuziget und umgetragen haben, zu Hohn und Spott GOtt, unserm HErrn JEsu Christo, mit viel schändlichen Lästerworten: solche bösen, giftigen GOttes- und Katzenkreuziger, wenn sie glauben könnten, oder müßten (ohne Glauben) die Wahrheit der Sprache sonst bekennen, so würden sie also sagen: Ja, ihr verfluchten Gojim, wenn das wahr wäre, daß des Weibes Same GOtt und Mensch wäre, so müßten wir selbst wohl, daß der Text sich aus der Maßen fein drauf reimt, da Eva spricht: „Ich habe den Mann kriegt, den Jehova“, und bekennen frei, daß die Sprache gerne und fein gibt, daß dieser Sohn, der Mann und GOtt, der HErr, wäre. Was man aber anders hier deutet, als: Ich habe den Mann kriegt, durch den HErrn, oder von dem HErrn, oder mit dem HErrn, das ist genöthigt, gezwungen, unartig Ding, und nicht die rechte Art und Natur der Sprache, kann es auch niemand anders beweisen. Ja, auf die Weise müßten die bösen Leute bekennen. Aber nun sie nicht leiden können, daß GOtt Mensch sei geboren von einem Weibsbild, muß dieser Text und die ganze Schrift unrecht haben, oder von ihnen eine andere Nase machen lassen.

106. Eben also müßten auch alle anderen Hebräischen bekennen, wenn sie den Text recht ansehen, und hielten, daß dieser Weibes Same Jehova, das ist GOtt und Mensch, wäre. Denn daß dies Wörtlein \hebraeisch{את} heiße „den“ oder „die“, und eine nota accusativi sei, das ist überwiesen, überzeigt, bekannt von allen Hebräisten, Juden und Christen in allen Grammatiken. Daß es aber auch sollte heißen, ad, de, vel cum, von, oder mit, oder durch, das ist noch unbeweisen, und soll wohl unbeweisen bleiben. Denn auf die Exempel, die sie führen aus Rabbi Kimchi, oder aus der Schrift, kann man leichtlich sagen, daß die hebräische Sprache noch nie wieder aufgekommen ist, und die Juden nicht wissen können virtutem omnium vocabulorum, sicut res ostendit; viel weniger müssen sie vim phrasis, figurarum et idiotismorum; sondern sie zweifeln, äquivociren, tappen und suchen, wie ein ungelehrter Organist die Claves oder Drechselpfeifen sucht, und fragt: Bist du’s, bist du’s?

107. Denn, wie der lateinischen Sprache Lehrer schreiben, ist es gar viel ein ander Ding, lateinisch reden, und grammatisch reden: also ist es auch gar viel ein ander Ding, hebräisch reden, und grammatisch reden. Grammatikisch mögen sie reden, doch sehr unfertig; aber hebräisch, rein, gut, und fertig reden, ist nunmehr nicht möglich. Es lernt ein jedermann gar viel besser deutsch, oder andere Sprachen, aus der mündlichen Rede, im Hause, auf dem Markt und in der Predigt, denn aus den Büchern. Die Buchstaben sind todte Wörter, die mündliche Rede sind lebendige Wörter, die geben sich nicht so eigentlich und gut in die Schrift, als sie der Geist oder Seele des Menschen durch den Mund gibt. Wie St. Hieronymus von Demosthene und Aeschine schreibt in prologo, und von Livio: Habet, nescio quid, latentis energiae viva vox. Sonderlich taugt das nichts, da sie vorgeben, \hebraeisch{אֶת} möge heißen de, a, ab, das ist, von; als: Ich habe den Mann kriegt vom HErrn; denn die Exempel 1 Mos. 44, 4. 2 Mos. 9, 29.: Egressi ex urbe, und dergleichen, thun es nicht; denn man sagt recht wohl: Egressus urbem, aedificavit lapides in altare, aber appositive.

108. Daß aber auch Mose schreibt, 1 Mos. 5, 22. und 6, 9.: Henoch ambulavit \hebraeisch{עם} Deum, Noha ambulavit \hebraeisch{את} Deum; daß deuten sie: Henoch und Noah wandelten mit GOtt; das taugt nicht, und lautet auch nicht. Wohin haben sie mit GOtt gewandelt? gegen Morgen oder gegen Abend? Es soll heißen: ambulavit Deum, in accusativo; wie die Lateinischen auch so reden: Vixit Sardanapalum; qui Curios\footnote{Curios Dentatus, der Besieger des Pyrrhus, bekannt durch seine große Mäßigkeit.} simulant et bacchanalia vivunt; exuit patrem; sic: Noha ambulavit Deum, id est, divinam viam, duxit vitam divinam; gessit et fecit opus Dei. So redet auch St. Paulus Gal. 1, 10.: An suadeo Deum, vel homines? id est: An doceo divina vel humana? Ibidem Cap. 2, 20.: Quod vivo. Röm. 6, 10.: Quod vivit, vivit Deo. Idem 1 Petr. 4, 6. Das und alles mehr befehl ich den Hebräisten; also, 1 Mos. 39, 2.: Der HErr war \hebraeisch{עם} Joseph, mit Joseph. Hier müssen wir Deutschen wohl so sagen, aber es gibt den accusativum nicht wohl, und ist doch accusativi nota in Ebraeo; das soll es auch bleiben. Das sei genug von dem Spruch, da
\quellseite{978}
Eva, oder vielmehr Mose, mit dem Neuen Testament stimmt, daß dieser Same des Weibes Jehova sei, und von ihr und Mose also verstanden, sonst hätten sie es beide wohl anders reden können.

109. Hieher gehört auch der Spruch Mosis, 1 Mos. 22, 18., da GOtt Abraham mit einem Eide verheißt, und spricht: „In deinem Samen sollen alle \hebraeisch{גוים} (Heiden) auf Erden gesegnet werden.“ Hier steht das Wort \hebraeisch{גוים}, damit uns diejenigen Juden (sind sie anders Juden) schmähen und fluchen, eben darum, daß wir uns dieses Segens rühmen, den GOtt Abraham verheißt, und spricht: „Alle \hebraeisch{גוים} sollen gesegnet werden in deinem Samen“; sie aber, die beschnittenen Heiligen, wollen, daß wir Heiden sollen verflucht, und allein sie der Same Abrahä sein. Aber weil sie den Heiden fluchen, und ein solcher Same sind, durch welchen alle Heiden verflucht werden, ist es offenbar, daß sie nicht Abrahams Same, sondern des Teufels Same sind. Denn GOtt ist es, deß Urtheil recht und gewiß ist, der spricht, daß Abrahams Same soll den Heiden nicht fluchen, wie sie thun, sondern alle Heiden sollen in ihm gesegnet werden; wie denn bisher nun bei 1543 Jahren geschehen, und noch bis [in] Ewigkeit geschehen wird.

110. Nun, dieser Segen ist nicht ein menschlicher Segen, da man mit Worten segnet, oder guten Morgen oder Abend wünscht; denn anders vermag der Mensch nicht segnen. Auch ist es nicht ein teuflischer Segen, damit die Zauberinnen die Kinder, Vieh und dergleichen segnen, daß es solle gedeihen, und vor Unglück behütet werden. Auch nicht ein jüdischer Segen, der durchs Schamhaperees\footnote{Vergleiche Luther's Schrift vom Schem Hamphoras, St. Louis. Ausg., Bd. XX, 2051, z. 60. In § 146 dieser Schrift findet sich richtig „Schamhapereß“, während die Ausgaben hier „Schemhaparees“ bieten.} und ihre Zauberei, mit Buchstaben und Figuren, oder GOttes Namen Tetragrammaton, soll kräftig sein und Wunder thun; wie der Türken Segen auch des Teufels Segen und Abgötterei ist, da sie sich mit Briefen und Worten segnen im Streit, wider Eisen und alle Waffen. Ja, es ist auch nicht ein päbstlicher Segen, der das Wasser und Wachs bezaubert, daß sie Weihwasser und Agnus Dei werden, und viel Tugenden, über ihre natürliche Tugend, kriegen sollen und helfen: sondern es ist ein göttlicher Segen, den GOtt allein geben kann und will. Solcher Segen ist nicht ein bloß, ledig Wort, das uns guten Morgen gibt oder wünscht, und nichts daraus folgt, sondern gibt und schafft alles, das es spricht. Als 1 Mos. 1, 28. segnete GOtt alle Thiere und Menschen, und sprach: \hebraeisch{פרו ורבו}: „Seid fruchtbar, und mehret euch.“ Da bleibt es nicht bei dem bloßen Worte, sondern folgt daraus das Werk, nämlich, daß Thiere und Menschen fruchtbar wurden, und sich mehrten, bis sie die Welt erfüllten. Und solcher Segen steht und geht noch jetzt bis an der Welt Ende. Denn durch solchen sind wir, was wir sind und haben, an Leib, Seele, Gut und allem, was da ist oder wird.

111. Also ist dieser göttliche Segen, in Abrahams Samen verheißen, auch ein thätlicher, wirklicher, lebendiger Segen, der da schafft, was er sagt oder segnet. Denn er ist verheissen und gegeben wider den Fluch, darein uns die Schlange gefället hat durch Adams Ungehorsam und Sünde, und ist hiermit die Verheißung von des Weibes Samen verneuet, und soll nun Abrahams Same heißen; wie er weiter hernach Davids Same, und zuletzt der Jungfrau Same worden ist. Darum heißt hier „Segen in dem Samen Abraham“ eben so viel, als droben [1 Mos. 3, 15.]: „Der Same des Weibes soll der Schlange den Kopf zertreten“, das ist, er soll die Sünde und den Tod wegnehmen, und Unschuld und Leben wiederbringen. Denn Sünde und Tod ist der Fluch, darunter wir liegen müßten ewiglich, wo wir nicht durch diesen Samen wiederum gesegnet, das ist, lebendig und gerecht, heilig und selig gemacht würden. Ja, also werden wir in diesem Samen Abrahä gesegnet; ja, des Segens rühmen wir uns Gojim, und nehmen uns sein an durch den Glauben, sind sehr hoffärtig, stolz und prächtig wider den Teufel, seine Gewalt, Tod und Sünde, und alles, was deß mehr ist; singen und sagen also: In dem Samen Abrahams, Davids und Mariä des Weibes, haben wir Vergebung der Sünden, Abwaschung der Sünde, Erlösung von der Sünde, Erlösung vom Tode und allem Uebel: „Denn er ist uns von GOtt gemacht“, 1 Cor. 1, 30., „unsere Gerechtigkeit, unsere Weisheit, unsere Erlösung, unsere Heiligung“, unser Segen, unser Trost, Leben und Freude in Ewigkeit. Deß sei GOtt gelobt in Ewigkeit, Amen.

112. Soll nun dieser Same Abrahams solchen
\quellseite{979}
starken, thätlichen Segen geben und schaffen unter den Heiden, so wird er nicht müssen ein lauter Mensch sein, der nicht mehr könne, denn guten Morgen zu uns sagen, welches alle Menschen können, sondern muß der rechte, natürliche, einige GOtt sein, der solchen Segen gewaltiglich in der Hand habe. Denn, Sünde und Tod aufheben, Gerechtigkeit und Leben geben, sind nicht Menschen- noch Engel-Werke, sondern allein der einigen, ewigen, göttlichen Majestät, Schöpfers Himmels und der Erde. Wiederum: Soll er Abrahams Same, das ist, sein Kind und Sohn sein, so muß er nicht eitel lauter GOtt sein, sondern ein rechter, natürlicher Mensch, vom Fleisch und Blut Abrahams; das ist, er muß zugleich GOtt und Mensch sein, in Einer Person. Weiter: Weil er nicht die Person ist, die zu Abraham von dem Samen oder dieser Person spricht: „In deinem Samen sollen alle Heiden gesegnet sein“, so muß er eine andere, unterschiedliche Person sein. Denn der zu Abraham spricht: „In deinem Samen“ etc., ist nicht Abrahams Same, sondern redet als von einem andern, der Abrahams Same solle sein. Daraus schließet sich der Unterschied der zwo Personen, und bleibt doch der einige, ungetrennte GOtt in seinem einigen göttlichen Wesen. Hierbei findet sich die dritte Person zugleich, der solches durch Mose oder Engel mündlich ausspricht von diesen zwo Personen; wie droben [§ 9 ff.] gesagt ist, daß dem Heiligen Geist zugemessen wird das Aussprechen des mündlichen Worts, drinnen er uns unterscheidlich offenbart wird; gleichwie die Menschheit des Sohnes unterschiedliche und eigene Offenbarung ist.

113. Bei uns Christen schließet sich’s noch weiter heraus, daß dieses Samens Abrahams Mutter müsse eine Jungfrau sein, die ihn ohne Sünde vom Heiligen Geist empfangen und gebären sollte. Denn wo er von einem Manne empfangen sollte werden (wie andere Adamskinder), so müßte er in Sünden selbst auch empfangen werden, wie der 51. Psalm, V. 7., klagt von allen Menschen: „Siehe, in Sünden bin ich empfangen“ etc. Der Weise nach würde er selbst eines andern Samens bedürfen, in welchem er müßte gesegnet, das ist, von Sünden und Tode erlöst werden, und würde uns kein Segen sein, noch geben können. Aber von diesem Spruche haben wir St. Paulum einen reichen Prediger, sonderlich zu Römern und Galatern, da er Abrahams und Davids Samen auf diese Weise meisterlich lehrt, daß nicht noth ist diesmal uns Christen, weiter davon zu handeln; denn es ist unser täglich Brod und stetige Predigt, Lesen und Singen.

114. Nun siehe, ob Mose nicht ein guter Christ sei, der so trefflich lieblich stimmt mit St. Paulo und dem ganzen Neuen Testamente. Sollten die Fluchjuden und Teufelsnamen solchen Ketzer nicht steinigen, wie sie in der Wüste oft thun wollten? Sollte der ihr Prophet und Meister sein? Oder ist nicht werth, mit solcher Ketzerei, daß ein beschnittener Heiliger ihn sollte mit seinen allerheiligsten Ohren nennen hören; sein Name muß mit den verfluchten Gojim, denen er solchen herrlichen, seligen Segen verkündigt, verflucht sein. Wiewohl er sie nicht ausdrücklich mit dem Worte „alle Gojim“; denn auch das Volk Israel oft Goj in der Schrift genannt wird; sondern sie schließen sich selbst aus, wie David von ihnen weissagt Ps. 109, 17. 18.: „Er will des Segens nicht, so wird er auch ferne von ihm bleiben; er will den Fluch haben, der wird ihm auch kommen. Und zeucht den Fluch an, wie ein Hemd (das nächste Kleid am Leibe), und ist in sein Inwendiges gegangen, wie Wasser (durch Blut und Fleisch), und wie Oel in seine Gebeine“ (durch Mark und Beine). Wir Christen verstehen nun wohl das Wort Christi Joh. 5, 46.: „Wenn ihr Mose glaubtet, so glaubtet ihr auch mir, denn von mir hat er geschrieben.“ Freilich geschrieben, durch sein ganzes Buch, wo er von GOtt redet und von Messia. Item, das Wort Joh. 8, 56.: „Abraham, euer Vater, ward froh, daß er meinen Tag sehen sollte, und er sahe ihn, und freuete sich.“ Wo sahe er ihn? In diesem Spruch, da er hörte, wie sein Same solle GOtt und Mensch sein, der alle Heiden segnen, von Sünde und Tod erlösen, ewiglich lebendig, heilig und selig machen sollte; gleichwie droben David, 1 Chron. 18, 16, eben dieselbige Freude hatte, da ihm derselbe Sohn auch verheißen ward.

115. Noch einen Spruch wollen wir aus Mose hören. 2 Mos. 33, 19. 20., da GOtt über das Volk erzürnt war um des güldenen Kalbes willen, und schlecht nicht mit dem Volke ziehen, noch sein sich annehmen, sondern Mose alles befahl, und einen Engel ihm zuordnen wollte, wollte mit dem Volke nicht mehr reden, sondern redete allein mit Mose, der sprach: „Laß mich deine Herrlichkeit
\quellseite{980}
sehen. Der HErr antwortete: „Ich will vor deinem Angesicht vorüber lassen gehen alle mein Gutes, und ich will predigen in des HErrn Namen vor dir, und will gnädig sein, wem ich gnädig bin, und mich erbarmen, weß ich mich erbarme.“ Und sprach: „Mein Angesicht kannst du nicht sehen, denn es kann kein Mensch mich sehen, und lebendig bleiben.“ Siehe mir diesen Text an, ungeachtet was Raben\footnote{Erlanger: „Rabben.“ „Raben“ = Rabbinen; vgl. Walch, St. Louiser Ausg., Bd. XX, 2093. 2103.} oder Juden hierin deuteln, ob er sich auch nach der unverfälschten Art der Sprache mit dem Neuen Testamente reime. Hier antwortet der HErr Mosi, da er seine Herrlichkeit zu sehen begehrt, und spricht: Es könne nicht sein. Aber gleichwohl verheißt er Mosi, er wolle alle sein Gut vor ihm übergehen lassen. Das ist eine Person, der Vater, der von seinem Sohne (der alle sein Gut ist, durch welchen er alles gemacht hat) redet, den soll Mose, das ist, sein Regiment und Volk Israel, sehen, nicht in der Herrlichkeit, sondern in einem Uebergange, hier in zeitlichem Leben. Denn Mose ist in diesen Geschichten und Gesichten nicht der geborne Mose von seinem Vater Amram, ein privatus Mose, sondern der berufene Prophet und Haupt des Volks Israel, dem er das Gesetz gibt.

116. Flugs darauf spricht eine andere Person: „Und ich will predigen vor dir im Namen des HErrn.“ Hier hörst du, daß der HErr will predigen vor Mose, das ist, vor dem Volk Israel, im Namen des HErrn. Was ist das gesagt: Ich HErr will predigen im Namen des HErrn? Müssen es nicht zwo unterschiedliche Personen sein? Ein HErr, der da predigt, und ein HErr, in deß Namen der HErr predigt. Nun muß dieser Prediger, der ein HErr, gewißlich ein Mensch werden, soll er vor Mose und Israel predigen. Denn das Predigtamt hat GOtt den Menschen, als Propheten und Aposteln, befohlen, durch welche er sein Wort uns verkündigt. Was aber die Predigt sein soll in des HErrn Namen, folgt: „Ich will gnädig sein, wem ich gnädig bin, und mich erbarmen, weß ich mich erbarme“, das ist, ich will nicht predigen, wie du, Mose, predigen mußt. Denn du mußt predigen das Gesetz, also: Ich gebiete dir heute, daß du dies und das thuest und haltest; wo nicht, so soll dir es nicht wohl gehen. Ich will aber also predigen, daß vor GOtt dem HErrn kein Mensch durchs Gesetz fromm oder gerecht sei; denn keiner hält es, wie er es soll, und schuldig ist zu halten. Darum macht deine Predigt eitel elende Leute, zeigt ihnen an ihre Sünde, vor welcher sie nicht können das Gesetz halten; darum es St. Paulus nennt ein Sündenamt und Todesamt, 2 Cor. 3, 6. Gal. 2, 16.

117. Aber meine Predigt im Namen des HErrn heißt also: Der HErr will es selber thun, und soll eigen Verdienst und Gerechtigkeit nichts gelten, sondern wer es haben wird, der wird es allein aus Gnaden und Barmherzigkeit haben. Wer Gnade und Barmherzigkeit ohne Verdienst sucht, der soll es haben. Das heißt: „Wem ich gnädig bin, dem bin ich gnädig.“ Nicht heißt es also: Wer das Gesetz hat oder Verdienst rühmt, dem bin ich gnädig, sondern, der meiner Gnaden sich rühmt, des gnädiger Gott bin ich. Daß also dieser Spruch vornehmlich straft und gewaltig geht, nicht wider die elenden Sünder, die vom Gesetze gefangen sind, sondern wider den halsstarrigen, eisernen, ehernen Stolz eigener Gerechtigkeit. Siehe, das heißt im Namen des HErrn gepredigt, das ist, GOtt will das thun, das Christus predigt; der predigt aber eitel Gnade, und spricht Joh. 7, 19.: „Euer keiner hält das Gesetz; und Joh. 8, 24.: „Wenn ihr nicht glaubet, daß ich es bin (das ist, daß ich der Jehova, der Ewige, der Gott selbst sei), so müßt ihr in euren Sünden sterben.“ Joh. 1, 17.: „Das Gesetz ist durch Mosen gegeben, aber die Gnade und die Wahrheit ist durch Jesum Christum geschehen.“

118. Nun halte solche Meinung des Textes Mosis an diesem Ort gegen das Neue Testament, und sage mir, ob sie sich nicht mit demselben gar fein, ungezwungen und gar gerne reime, daß man kein Wort mit seltsamen Glossen wider seine natürliche Deutung nöthigen noch martern darf, sondern, wie sie lauten in der hebräischen Sprache, so stimmen sie mit unserm christlichen Glauben, der uns lehrt im Neuen Testament, daß Jesus Christus Jehova, GOtt und Mensch sei, und des Volks Israel Prediger gewesen ist; wie St. Paulus ihn Röm. 15, 8. nennt, minister circumcisionis, einen Prediger des beschnittenen Volks Israels; und er selbst Matth. 15, 24. sagt: „Ich bin nicht gesandt, denn zu den verlornen Schafen Israel“; verbot auch den Aposteln, sie sollten unter die Heiden nicht gehen;
\quellseite{981}
und hier zu Mose spricht er: „Ich will vor dir predigen.“ Als sollte er sagen: „Ich will persönlich allein in deinem Volke, und des beschnittenen Israel Prediger sein, sonderlich der Elenden, welche du mit deinem Gesetze demütiget hast; wie Jes. 61, 1. sagt: „Er hat mich gesandt, das Evangelium zu predigen den Elenden.“ Also ist das Evangelium nichts Anderes denn das Wort Christi, GOttes Sohnes, darin er uns eitel Gnade und Barmherzigkeit predigt, in GOtt des Vaters Namen, der ihn dazu gesandt hat, und alles selbst durch ihn uns thut.

119. Das ist der Uebergang vor Mose und seinem Volke, darin er uns erkennt, und alle GOttes Güte uns erzeigt ist, und doch seines göttlichen Wesens Herrlichkeit nicht bloß\footnote{bloß = unverhüllt. Cruciger: nuda.} gesehen ist noch wird. Denn das gehört nicht in dies Leben, sondern, wenn wir gestorben sind, alsdann wird es geschehen, wie er spricht 2 Mos. 33, 20.: „Kein lebendiger Mensch wird mich sehen.“ Hiermit ist nicht versagt, daß GOtt kein Mensch nimmermehr sehen werde, ja, es ist vielmehr damit die Auferstehung der Todten zugesagt, da wir ihn sehen werden. Allein ist es gesagt von diesem Leben. „Kein Mensch“, spricht er, „wird mich sehen und leben“, das ist, sehen mag mich wohl der Mensch, aber leben muß er nicht, muß zuvor sterben, und in ein ander Leben kommen; da soll es ihm unverjagt sein. Da wird er verstehen, daß ich ihm gnädig sei, denn ich gnädig gewesen bin; und daß ich ihm gar nicht gnädig gewesen sei, weder um seiner Gerechtigkeit noch um der Werke des Gesetzes willen.

120. Wohl weiß ich, daß dies hebräische Wort \hebraeisch{קרא}, predigen, kann auch heißen rufen, nennen, lesen; wie auch Lyra und Burgensis zeugen, darnach es einen Buchstaben anders und anders um sich hat; aber wie es hier steht, mit dem Wörtlein „im“ [\hebraeisch{ב}], heißt es gemeiniglich predigen, 1 Mos. 4, 26. 12, 8. 13, 4. Und ob das gleich die Rabbininen oder zünftige Hebräischen nicht annehmen, da liegt mir nicht Macht an; mir ist genug daran, wie ich oft gesagt, wenn sich Mosis Wort, nach guter Art der hebräischen Sprache, so fein und gerne ablöst und gibt von der Rabbininen genöthigter Deutung zum Neuen Testament, daß jedermann sagen muß (wer der Sprache kundig ist, ob er schon nicht ein Christ wäre): Wohl, wenn der Christen Glaube recht ist, so hat wahrlich Mose mit diesen Worten ihren Verstand gehabt; denn die Worte reimen sich doch nirgend so fein und gewiß, als zum Neuen Testament. Und auf solche Weise wollte ich gerne die ganze hebräische Bibel den Juden nehmen von ihren schändlichen, lästerlichen Glossen; aber es ist nicht Eines Mannes Arbeit, es ist genug, andern, die gelehrter sind denn ich, ein Exempel oder meinen guten Willen erzeigt, daß sie es mehr und besser machen.

121. Das aber folgt an diesem Orte [2 Mos. 33, 21–23.]: „Und der HErr sprach zu Mose: Siehe, es ist ein Raum bei mir, da sollst du auf dem Felsen stehen, und wenn meine Herrlichkeit wird vorüber gehen, will ich dich in die Felskluft thun, und meine Hand über dir halten, bis ich vorüber gehe. Und wenn ich meine Hand von dir thue, wirst du mein Hinterstes sehen, aber mein Angesicht kann man nicht sehen.“ Hier reden auch zwo Personen Jehova. Eine spricht: „Wenn meine Herrlichkeit wird vorüber gehen“; dies ist die Person des Vaters, der vom Uebergange seiner Herrlichkeit, das ist, seines Sohnes, redet; und er selbst der Sohn spricht: er sei es, der vorüber gehe. Das ist alles von Christo, GOtt und Menschen, hier auf Erden gegangen, gesagt; wie droben [§§ 115. 116] gehört ist.

122. Das andere, da er spricht: „Es ist ein Raum bei mir“, und was er von dem Felsen und seiner Hand über Mose redet, bis er vorüber gehe, verstehe ich also, daß GOtt um des künftigen Felsens, Christi, willen, das Volk des Gesetzes oder Israel dennoch geschützt und erhalten hat, unter seiner Geduld, weil sie das Gesetz nicht halten konnten. Auf welche Weise Paulus Röm. 3, 25. redet, daß die Sünde, so unter dem Gesetze blieb unter göttlicher Geduld, zu dieser Zeit sei vergeben, da Christus gekommen ist mit seinem Uebergange \&c. Aber nach diesem Uebergange hat GOtt die Hand solcher Geduld und Schutz des Felsens weggethan. Denn das Gesetz ist aus und erfüllt, [wir] dürfen seiner Geduld noch Schutzes des künftigen Christi. Ja, verdammt ist der, so Christum noch künftig glaubt, und mit Mose im Felsen unter GOttes Hand noch stehen will. Der Fels und Hand ist weg, wir haben den HErrn und seinen Uebergang, bis auf welchen der Schutz und Geduld GOttes währen sollte. Wir sehen
\quellseite{982}
nun ihm nach, was er gemacht hat vor uns,\footnote{Erlanger: ante nos. Erlanger: „für“.} das ist, was sein Hinterstes ist, was er uns hinter sich gelassen hat, nämlich, daß er (G)Ott und Mensch für uns gestorben und auferstanden ist; und möchte also die Menschheit sein Rücken oder Hinterstes heißen, darin wir ihn erkennen in diesem Leben, bis wir dorthin kommen, da wir sein Angesicht und Herrlichkeit auch sehen werden.

123. Eben zu solchem Prediger macht Mose den HErrn im folgenden 34. Capitel, V. 5. ff.: „Und der HErr (JEsus Christus) fuhr hernieder in einer Wolke, und trat bei ihm (Mosen) daselbst, und predigte im Namen des HErrn. Und der HErr ging vor ihm über und rief (predigte): HErr, HErr GOtt, barmherzig und gnädig, geduldig, und [von] großer\footnote{Erlanger: und große.} Güte und wahrhaftig, der seine Güte hält in tausend Glied, der da vergibt Sünde, Missethat und Uebertretung, vor welchem niemand unschuldig ist, der da heimsucht die Missethat der Väter über die Kinder und Kindes Kinder, bis ins dritte und vierte Glied.“ Es ist hier die lateinische Bibel zumal falsch, wer es auch gethan hat, die setzt „Mosen“, da „HErr“ stehen soll.\footnote{In der Vulgata lautet 2 Mos. 34, 5.: Cumque descendisset Dominus per nubem, stetit Moyses cum eo, invocans nomen Domini.} [Es] hat vielleicht solchen Meister hinbillig gedacht, daß der HErr vom HErrn predigen oder rufen sollte, welches besser Mosi anstände. Burgensis gefällt mir auch wohl, der hält, daß der Text also stehe im Hebräischen: Und der HErr ging vor ihm über und rief (oder predigte) den HErrn HErrn GOtt, den\footnote{„den“ fehlt in der Erlanger.} Barmherzigen, den Gnädigen; in accusativo, oder, wie wir es deutsch sagen: Er predigte von dem HErrn HErrn GOtt. Doch ist es und bleibt einerlei Meinung; denn es gleich viel ist im Deutschen, den HErrn, und von dem HErrn predigen.

124. Wohlan, das ist ja auch ein heller Text, daß der HErr ein Prediger ist, und predigt im Namen des HErrn; da sind zween HErren genannt, und doch nicht zween Götter noch HErren sind; ja, der HErr, spricht er, predigt von dem HErrn, von dem HErrn, von dem GOtt; da steht „HErr“ zweimal und GOtt dabei, deren sind drei; noch müssen es nicht drei Götter sein. Es ist droben [§ 116 ff.] gesagt, was da heißt, der HErr predigt im Namen des HErrn, nämlich, daß JEsus Christus sei dieser Prediger, GOtt und Mensch, der im Namen seines Vaters, vom Vater gesandt, und von seinem Vater dem Mosi, das ist, dem Volk Mosis gepredigt hat von eitel Gnade und Barmherzigkeit, daß niemand durchs Gesetz, weil es niemand hält, kann gerecht werden. Eben dieselbige Predigt thut er hier auch mit andern Worten, da er spricht: „Der HErr tritt bei Mosen und predigt.“ Warum tritt er bei Mosen, und nicht über ihn, oder ferne von ihm? Die zwei Predigtämter, Gesetz und Evangelium, müssen bei einander sein, obwohl die Geschäfte ungleich sind. Denn Mose predigt von Sünden, und tödtet dadurch; Christus predigt von Gnaden, und macht lebendig dadurch. Doch kann Gnade nichts schaffen, wo durchs Gesetz nicht zuvor die Sünde offenbart und erkannt ist. Der HErr Christus spricht selbst Matth. 11, 5. [15, 24.], er predige den Elenden das Evangelium, und den verlornen Schafen Israel, das ist, die sich durchs Gesetz verloren fühlen.

125. Was predigt nun der HErr bei Mose und vor Mose? Er predigt, spricht er, von HErrn HErrn GOtt, Gnädigem und Barmherzigen etc. Das ist, daß drei Personen Ein GOtt sind, vor welchem nichts gilt eigen Verdienst aus dem Gesetze, die da gar nichts vor ihm, und in der Wahrheit auch nichts sind; sondern eitel Gnade und Barmherzigkeit, Güte und Treue, der die Sünde, Uebertretung, Missethat vergibt, und niemand unschuldig vor ihm ist. Willst du nun GOtt halten und nennen bei seinem rechten Namen, wie hier geschrieben steht, daß er ein Vergeber der Sünden, gnädig und barmherzig, und niemand vor ihm unschuldig sei, so mußt du nicht viel Rühmens vor ihm machen, du seiest Mose, Johannes, oder wer du willst; sondern mußt mit St. Paulo, Röm. 3, 23. sagen: „Alle Welt ist vor GOtt schuldig, und mangeln des Ruhms an GOtt“, oder, daß sie nicht sagen dürfen, sie seien unschuldig und gerecht vor GOtt; wo anders, so werden sie dem letzten Stück unterworfen bleiben, da er spricht [V. 5.]: Er sucht heim die Missethat der Väter bis ins vierte Glied; wie Christus auch spricht Marc. 16, 16.: Wer nicht glaubt, wird verdammt.

126. Und ist zu merken, daß hier nicht steht, wie im ersten Gebot, 2 Mos. 20, 6.: „Der Barmherzigkeit thut bis ins tausendste Glied, denen,
\quellseite{983}
die mich lieben, und halten meine Gebote“; sondern dafür steht: „Es ist niemand unschuldig vor ihm“, das ist, niemand liebt ihn, noch hält seine Gebote, ohne die allein, so kein Verdienst rühmen; sondern GOtt gnädig und barmherzig, und einen Vergeber der Sünden glauben und nennen, sich selbst aber schuldig bekennen, und beten im Vater-Unser [Matth. 6, 12.]: „Vergib uns unsere Schuld, als wir vergeben unsern Schuldigern.“ Das heißt ja nichts anderes, denn eitel Gnade gepredigt; nicht, was wir thun sollen, wie die zehn Gebote fordern, und nicht geschieht; sondern was GOtt mit Gnaden an uns thun will, und gethan hat; wie das Neue Testament uns lehrt und zeugt. Da haben wir nun den Prediger, der sich Mosi offenbart, und weissagt, was er im Neuen Testament predigen wolle. Und wie er dazumal geweissagt hat bei Mose, also sehen wir es im Neuen Testament erfüllt, nämlich, daß niemand durch eigene Gerechtigkeit, sondern allein durch GOttes Gnade, die uns durch diesen Prediger, seinen lieben Sohn, gepredigt ist, soll gerecht und selig sein.

127. Daß nun hier weiter folgt, wie Moses bittet, daß der HErr ja mit ziehen wolle, und sie nicht lassen, und der HErr darauf antwortet, er wolle mit ziehen, große Wunder thun etc., da ist er wieder versöhnt mit dem Volke, erneuet den Bund, und schreibt andere Tafeln, und faßt in der Summa kurz des alten Testaments Gebote und Gottesdienst, wie sie leben und thun sollen. Aber von der Gnade und Vergebung, wie droben [V. 6. 7.], steht hier nichts. Und ist die Meinung: Mose hat nun die tröstliche Verheißung des neuen Testaments, da der HErr selber predigen und regieren will; nun bittet er, weil ihm dies Volk befohlen ist zu lehren und zu regieren bis aufs neue Testament, der HErr wolle bei und mit ihm sein. Denn was soll ich machen? [2 Mos. 34, 9.:] Es ist ein halsstarrig böse Volk. Wo du nicht mit uns bist, und die Sünde vergibst, und geduldiglich trägst, bis du selber kommst, und Prediger der Gnaden werdest, so sind wir alle Stunden verloren. Wir müssen deine göttliche Geduld und Schutz in diesem Regimente haben, da wir dein Gesetz predigen sollen, und doch nicht halten werden. Und ist eben das, das im 33. Capitel, V. 22., von dem Schutz und Hand GOttes in der Felsenkluft gesagt ist.

128. Denn GOtt antwortet und spricht [V. 10. f.]: Wohlan, ich will es thun, und will einen Bund machen mit deinem ganzen Volke, und Wunder thun, dergleichen nicht geschehen sind im ganzen Lande, und unter allen Heiden; und das ganze Volk, unter welchem du bist, soll sehen das Werk des HErrn, wie wunderlich es ist, das ich mit dir thun will. Siehe zu und halte, was ich dir heute gebiete etc. Dies ist alles vom alten Testament und Mosis Volke geredet; wie auch folgt von den Ausstößen der Amoriter, Cananiter, Hethiter etc., welches im alten Testament geschehen ist. Und mit Fleiß meidet er das Wort: „Mein Volk“; sondern nennt es Mosis Volk, dein Volk, und unter welchem du bist. Doch will ich, wie ich verheißen habe, meine Hand über ihm halten, im Felsen schützen, dazu große Wunder thun, die nie geschehen sind unter allen Heiden. Und wahr ist es, lies das Alte Testament durch und durch, so wirst du sehen, wie manche große Wunder GOtt in diesem Volke gethan hat, von Mose an bis auf Christum, ob es wohl nicht sein Volk, das ist, nicht der Gnaden, sondern des Gesetzes Volk ist; ausgenommen die, so Mosen verstanden, und auf Christum ihre Hoffnung gesetzt haben; der andere Haufe ist eitel Werktheilige, und halsstarrige, ruhmräthige\footnote{Wittenberger und Jenaer: „rhummetige“; Erlanger: ruhmredige. Cruciger: persuasione propriae justitiae sibi placentium.} Gesetztreiber gewesen.

129. Merke aber, wie der Text klar gibt, daß der HErr, so mit\footnote{„mit“ fehlt in der Erlanger.} Mose redet, ist JEsus Christus, der künftige Prediger des neuen Testaments; denn er hier auch unterschiedlic[h] redet von sich, und vom Vater, da er spricht 2 Mos. 34, 10.: „Alle dein Volk soll sehen das Wunderwerk des HErrn“, das ich thun will. Siehe, es ist des HErrn Wunderwerk, davon er jetzt redet, und er will es gleichwohl thun, was der HErr thut; gleich wie er Joh. 5, 19. auch redet: „Was der Vater thut, das thut auch der Sohn.“ Item, V. 17.: „Mein Vater wirket bisher, und ich wirke auch.“ Item, V. 21.: „Wie der Vater die Todten auferwecket und lebendig macht, also auch der Sohn macht lebendig, welche er will.“ Das heißt ja, Johanneum mit Mose, und Mosen mit Johanne gleich stimmen, und sich fast mit einerlei Wort reimen. Hier sind ja unterschiedliche zwo Personen, Vater und Sohn (wie Johannes
\quellseite{984}
sagt), und der HErr, der vom HErrn redet, und des HErrn Wunderwerk thut (wie Mose sagt), und ist doch einerlei, nicht zweierlei oder anderlei Werk; darum muß es auch nicht mehr, denn ein einiger HErr und GOtt sein.

130. Item, so redet weiter am selben Ort [V. 23.] der HErr mit Mose, und spricht: „Dreimal des Jahres soll ein jeglich Knäblein erscheinen vor dem Herrscher, dem HErrn GOtt Israel.“ Hier redet abermal der HErr vom Herrscher HErrn GOtt Israel. Denn es sind nicht Mosis Worte, sondern des HErrn, der mit Mose redet, und ihm das alte Testament befiehlt, daß er schützen und mit Geduld tragen will, bis auf seinen künftigen eigenen Uebergang, wie droben [§ 119] genug gesagt ist. Ob nun die Raben und Juden dies alles anders deuten, und unsern Verstand verachten, das ist recht; GOttes Feind soll GOttes Wort nicht sehen. Was sie aber hier über diesen Text speien, ist nicht werth, daß [es] eine Sau oder Esel lesen sollte, wenn sie gleich lesen könnten. Mosis Angesicht hat Hörner, und glänzt zu helle, daß sie nicht darein sehen können; wir aber haben Mosen, daß seine Worte ungezwungen, natürlicher Art der Sprache, so herzlich und lieblich sein stimmen mit dem Neuen Testament. Und ob er wohl muß das halsstarrige, böse Volk seiner Zeit regieren im alten Testament, so weissagt er doch daneben gewaltiglich von Jesu Christo, unserm HErrn, daß er ein wahrhaftiger Mensch, und mit dem Vater und Heiligem Geist, in unterschiedlicher Person, ein einiger wahrhaftiger GOtt sei, der alles thut, was der HErr thut. Das ist uns genug, wollen gerne Narren und Ungelehrte heißen in der Schrift, und den Juden und Türken ihre hohe Weisheit in ihrem Schlauraffenlande lassen.

131. Wohlan, glaube ein jeder für sich, was er will, ich glaube und weiß fürwahr, daß ich und alle Christen Mosen für uns haben, und daß er ein rechter Christ, ja, ein Lehrer der Christen ist. [Es] schadet nicht, daß er dazumal noch in der Kappe steckt, und im alten Testament gekleidet daher geht, als sei er nicht ein Christ, gleichwie ein frommer Mönch, als St. Bernhard, daher geht als ein Mönch, aber doch in seinem Glauben ein rechter ernstlicher Christ ist, der nicht auf seine Kappe noch Orden, wie der andere Haufe, bauet, pocht, und trotzt, sondern allein auf die Gnade Jesu Christi; wie er selbst oft zeuget. Also läßt Mose den andern Haufen sich des Gesetzes und Beschneidung rühmen, geht immer mit in solcher Kappe; aber sein Herz, Glaube und Bekenntniß ist JEsus Christus, GOttes Sohn etc. Haben wir nun Mosen, das ist, den Meister und Obersten, so werden seine Jünger, die Propheten, mit Haufen, ihm nach, zu uns fallen, denn sie nichts Anderes glauben, bekennen und lehren, weder Mose, ihr Präceptor, thut.

132. Aber wo wollen wir sie sehen, die lieben Gäste? Dies Büchlein ist zu enge, [ich] kann sie nicht alle hierein setzen, so Mose nicht ganz hierein kommen kann. Also wollen wir thun: Wir wollen zu ihnen gehen, und mit ihnen essen, sie haben die Küche und Keller besser bestellt denn wir, und können uns wohlauf vorsehen und herrlich speisen. Das ist, ein jeder nehme die Propheten vor sich, lese mit Fleiß darinnen, und merke, wo der HErr, Jehova, JEsus Christus, unterschiedlich redet, oder wo von ihm geredet wird. Denn du hast nun gehört, daß er es sei, der mit Mose auf dem Berge Sinai redet, Mosen und das Volk führt, und Wunder thut. Und ob er wohl nicht allein ist, der es thut, sondern der Vater und Heiliger Geist mit ihm alles und einerlei Werk thun, so offenbart er sich gleichwohl mit solchem Reden und Thun, daß er sei eine unterschiedliche Person vom Vater, in dem einigen göttlichen Wesen. Und zwar, wer so viel versteht in der Schrift (welches nicht ein jeder achtet), daß er merken kann, wo die Person eine von der andern redet, als nämen mehr denn eine da, der hat bald ersehen den Unterschied, welche des Vaters oder des Sohnes Person ist. Hast du aber den Unterschied des Vaters und des Sohnes, so ist des Heiligen Geistes Person gleich auch unterschiedlich mit da.

133. Als Psalm 2, 7.: „Der HErr sprach zu mir: Du bist mein Sohn, heute habe ich dich gezeuget.“ Und hier, im Mose, 2 Mos. 33, 19.: „Der HErr predigt im Namen des HErrn.“ 1 Mos. 19, 24.: „Der HErr ließ regnen Schwefel und Feuer vom HErrn.“ Hier siehst du klugs, daß der HErr, der regnen läßt, der Sohn ist, vom HErrn, das ist, vom Vater; denn der Sohn ist vom Vater, und nicht wiederum. Hos. 1, 7.: „So spricht der HErr: Ich will mich erbarmen des Hauses Juda, und ihnen helfen durch den HErrn, ihren GOtt, und nicht durch Schwert und Bogen“ etc. Zephanja 3, 9. spricht der HErr:

\quellseite{985}
„Ich will den Völkern zukehren eine reine Lippe, daß sie alle rufen oder predigen sollen im Namen des HErrn, und ihm dienen mit einerlei Schultern.“ Ps. 45, 8.: „Darum hat dich, GOtt, dein GOtt gesalbet mit Freudenöle“; und, V. 12.: „Der König wird Lust an deiner Schöne haben; denn er ist dein GOtt, und du sollst ihn anbeten.“ Jer. 23, 5. 6.: „Zu der Zeit will ich David ein gerecht Gewächs erwecken, und das wird sein Name sein, daß man ihn nennen wird, HErr, der unsere Gerechtigkeit ist.“

134. Wo aber die Person nicht unterschiedlich sich mit Reden offenbart, daß es scheint keine mehr, denn Eine Person [zu] sein, da magst du die Regel halten, droben [§ 132] gegeben, daß du nicht Unrecht thust, wo du den Namen Jehova deutest auf unsern HErrn Jesum Christum, GOttes Sohn. Des nimm ein fein Exempel Jes. 50, 1.: „So spricht der HErr: Wo ist der Scheidebrief eurer Mutter, damit ich sie gelassen habe?“ Hier heißt HErr die Person des Sohnes, ob er wohl nicht unterschiedlich von seiner Person redet; wie das Lyra und andere auch auslegen. Und hat mir vor vielen Jahren wohlgefallen, daß Lyra so frei heraus fährt: „So spricht der HErr, das ist, JEsus Christus.“ Doch, wenn man nach diesem Stücke das ganze Capitel liest (denn Jesaia redet hier kein Wort, sondern alles der HErr), so findet sich’s, daß die Person des Sohnes, JEsus Christus, hier redet, nicht allein nach der Gottheit, sondern auch nach der Menschheit. Denn er spricht Jes. 50, 6.: „Meinen Rücken habe ich dargehalten denen, die mich schlugen, und meine Wangen denen, die mich rauften, und mein Angesicht verbarg ich nicht vor Schmach und Speichel; denn der HErr HErr hilft mir“ etc. Lies das ganze Capitel, so findest du, daß GOtt der HErr ist, der da leidet, und vom HErrn HErrn Hülfe hat, das ist, Christus ist wahrer GOtt und Mensch.

135. Solch Exempel, da kein heller Unterschied der Person ist, nimmt die Epistel [an die] Hebräer Cap. 1, 6. und spricht aus dem 97. Psalm, V. 7.: „Da er einführet den Erstgebornen in die Welt, spricht er: Und es sollen ihn anbeten alle Engel GOttes.“ Hier sieht man kein sonderlich Zeichen, daß solcher Psalm von Jesu Christo, GOttes Sohne, zu verstehen sei, ohne daß er im Anfange spricht: „Der HErr ist König worden, des freuet sich das Erdreich“, welches doch ein Jude oder geistloser Mensch nicht würde von Christo lassen gesagt sein. Aber der Geist bekennt, daß keine Person König worden ist denn der Sohn; wie der andere Psalm zeuget, V. 6.: „Ich habe meinen König eingesetzt auf meinem heiligen Berge Zion.“ 1 Chron. 18, 14.: „Ich will ihn in mein Königreich setzen ewiglich.“ Daß er aber rechter GOtt sei, zeugt dieser Psalm, und spricht: „Alle Engel GOttes sollen ihn anbeten“; Hebräisch also: „Betet ihn an, alle Götter.“ „Götter“ kann hier nicht GOtt heißen, der nicht viel, sondern allein ein einiger GOtt ist; darum deutet er es „Engel“. Der ist aber GOtt, den sie anbeten, JEsus Christus, GOttes Sohn. Darum thun die nicht Unrecht, so diesen Psalm von Christo verstehen, wie die Epistel thut, ob sie gleich nicht setzen den Unterschied der Person.

136. Item, da dieselbe Epistel [Cap. 1, 10. 11.] aus dem 102. Psalm, V. 26. 27. führt: „Du hast im Anfang die Erde gegründet, und die Himmel sind deiner Hände Werk; sie werden vergehen, aber du bleibest“ etc. Hier findet die Vernunft auch kein Zeichen, daß dies solle von Jesu Christo unterschiedlich geredet und verstanden sein, wie die Epistel sagt, und hätte wohl andere hellere Sprüche im Psalter gefunden; aber er gibt uns (achte ich) solch Exempel, daß wir Christum in der Schrift suchen sollen, weil er gewißlich GOtt und Schöpfer ist, mit dem Vater und Heiligem Geiste; daß niemand daran Unrecht thun kann, wer von ihm sagt, er sei es, der Himmel und Erde geschaffen hat. Aber doch sollen wir fleißig suchen die unterschiedliche Offenbarung der Person des Sohnes, und die Worte ansehen, die seine Person unterschiedlich geben und offenbaren. Wer es nicht besser hat, der nehme dieses an, daß GOtt David verheißen hat, 1 Chron. 18, 11., daß sein Sohn Messia solle GOtt ein Haus bauen, und ewig darinnen HErr und König sein, davon viel Psalmen gemacht sind.

137. Von solchem Hause und Bau weissagt dieser Psalm [Ps. 102, 14. f.] auch, und bittet ängstlich, daß er kommen wolle, und Zion bauen. Das kann nicht vom leiblichen Zion gesagt sein, welches dazumal gebauet da stand, welches auch nicht das Gotteshaus noch Zion war, das Davids Sohn bauen sollte, und König darinnen sein. So haben wir droben [§ 23 f.] gehört, daß dieser Bauherr und Hausherr müsse GOtt
\quellseite{986}
sein, und doch Davids Sohn. Darum führt die Epistel diesen Psalm recht auf die Person Jesu Christi, der wohl ist Ein GOtt und Schöpfer mit GOtt dem Vater und Heiligem Geist, aber mit dem Bauen GOttes Hauses und Königreich unterschiedlicher eine andere Person denn der Vater offenbart wird. So ist auch das ein unterschiedliches Zeichen, daß Christus, GOttes Sohn, der Same Abrahams ist, durch welchen alle Heiden sollen gesegnet werden, und ihm gehorchen. 1 Mos. 22, 18. 49, 10.: „Dem Schilo werden die Heiden gehorchen“, und Ps. 2, 8.: „Heische von mir, so will ich dir die Heiden zum Erbe geben.“ Von solchem Reiche unter den Heiden redet dieser Psalm auch, und malt damit die Person des Sohnes etc.

138. Summa, es ist Ein GOtt, Ein HErr, Eine göttliche Majestät, Natur und Wesen aller drei Personen; aber es offenbart sich zuweilen die Person des Vaters, zuweilen des Sohnes, zuweilen des Heiligen Geistes. Welche sich nun offenbart, so ist es der einige GOtt in drei Personen, auf daß wir die göttliche Majestät recht erkennen, und nicht, wie Juden, Ketzer, Mahomet, blind feldeinhin glauben, als sei GOtt nichts mehr denn eine einige Person; das will GOtt nicht haben, sondern will erkannt sein, wie er sich uns offenbart. Und sonderlich ist es GOtt zu thun um die Offenbarung und Erkenntniß seines Sohnes, durch die ganze Schrift Altes und Neuen Testaments; alles geht’s auf den Sohn. Denn die Schrift ist gegeben um des Messias oder Weibes Samens willen, der alles wieder zurecht bringen soll, was die Schlange verderbt hat: Sünde, Tod, Zorn wegnehmen; Unschuld, Leben, Paradies und Himmelreich wiederbringen. Und gleichwie der Fall Adams ist geschehen, daß er sich unterschiedlich an den Sohn stieß, und doch damit zugleich wider alle drei Personen, als einigen GOtt, sündigte, also hat GOtt wiederum seinen Sohn allein und unterschiedlich lassen Mensch werden, daß Adam durch den Sohn unterschiedlich, an welchem er sich gestoßen und gefallen war, wieder aufgerichtet würde, und doch solche Aufrichtung oder Erlösung aller dreier Personen einerlei, als des einigen GOttes, Werk ist.

139. Denn da die Schlange, der Teufel, Adam fällen wollte, hetzte er ihn wider den Sohn unterschiedlich, da er sprach: „Ihr werdet GOtt gleich werden.“ Das war dem Sohne GOttes zu nahe, und wollten beide, der Teufel und Adam, den Sohn entsetzen seiner Ehren; denn allein der Sohn ist gleich, oder das Ebenbild des unsichtbaren GOttes, Col. 1, 15., und das Ebenbild seines göttlichen Wesens, der dem Vater gleich ist in Einer Gottheit, Hebr. 1, 3. Und der Teufel hatte zuvor denselben Fall im Himmel gethan, und gelernet mit seinen Engeln eben an demselben Ebenbilde, dem Sohne GOttes; er ließ ihm nicht genügen, daß er das allerschönste Bild GOttes unter allen Engeln war (doch nicht geboren von Ewigkeit, sondern anfänglich geschaffen), sondern wollte auch das inwendige, natürliche GOttes Ebenbild sein, dem Sohne gleich; wie die Väter den Spruch Jes. 14, 12–14. unter des Königs zu Babel Namen auslegen: „Wie bist du gefallen, du schöner Morgenstern? der du sprachest in deinem Herzen: Ich will in den Himmel, und gleich sein dem höchsten GOtt.“ Darum ist die Person des Sohnes unterschiedlich offenbart, und Mensch worden, daß wir an ihm durch seine Menschheit wieder aufstünden, die wir uns an seiner Gottheit gestoßen hatten und gefallen waren. Also ist nun die ganze Schrift, wie gesagt, alles theils auf Christum, Gottes und Marien Sohn; alles ist es zu thun um denselben Sohn, daß wir ihn unterschiedlich erkennen, und also den Vater und den Heiligen Geist, Einen GOtt, ewiglich sehen mögen. Wer den Sohn hat, dem steht die Schrift offen, und je größer und größer sein Glaube an Christum wird, je heller die Schrift ihm scheint.

\subsection[Von beiden Naturen in Christo]{[Von beiden Naturen in Christo.]\footnote{Diese Ueberschrift findet sich in der Wittenberger und in der Jenaer.}}

140. Wenn du nun glaubst und verstehst, daß Christus sei wahrhaftiger GOtt und Mensch, wie die Schrift uns lehrt, so siehe darnach zu, und lerne weiter gewiß werden, daß du die Person Christi nicht trennest, noch die zwo Naturen, oder das göttliche und menschliche Wesen, nicht mengenest in Ein Wesen, sondern die Natur hier unterscheidest, und die Person einig behaltest.

141. Denn hier haben sich auch viel Klüglinge gestoßen, daß sie entweder Gottheit und Menschheit Eine Natur, oder aber zwo Personen haben machen wollen, wie Nestorius und Eutyches mit ihren Gleichen; die Juden aber
\quellseite{987}
und Fürsten trefflich hohes und überhohes Geistes hier sind, und uns Christen für große Narren halten. Ist er GOtt (sprechen sie), wie kann er als ein Mensch sterben? denn GOtt ist unsterblich. Ist er Mensch, wie kann er GOttes Sohn sein? denn GOtt hat kein Weib. Hier geht es, wie man spricht: Geld hat Ehre, sprach der Frosch, und saß auf einem Heller. Hier lehren uns die hoch, hoch, noch höher, und allerhöchst verständigen Leute, die Türken, Juden, daß GOtt nicht sterben kann, und kein Eheweib habe; wie wollten oder könnten wir armen Christen solch hoch Ding immermehr wissen, wenn es uns tollen Gänsen und Eseln nicht solche überaus überhohen Meister zeigten, daß GOtt kein Weib habe, und nicht sterben könnte? Es wäre nicht Wunder, daß, wo ein Jude oder Türke geht, die\footnote{Erlanger: daß die.} Erde unter ihm vor solchem hohen Geiste so fröhlich würde, daß sie mit ihm über den Himmel spränge, und der Himmel desgleichen vor solcher großer Weisheit, mit Sternen, Sonnen und Monden herunter fiele, den Juden, Türken zum Füßen, oder in Abgrund der Hölle, denn es ist, fürwahr, eine unbegreifliche Weisheit, daß GOtt kein Weib hat, noch sterben könne. O HErr Gott, da müssen die Christen nichts von! Wer wollte GOtt eine Amme bestellen? Wo wollte er eine Kindermagd kriegen? Wer wollte ihn begraben? Wer wollte ihm zur Hochzeit pfeifen, und tanzen? Wer wollte ihm Seelmessen nachhalten?\footnote{Wittenberger und Erlanger: noch halten.} Ei, pfui die Christen, daß sie einen sterblichen GOtt anbeten, und einen ehelichen Mann draus machen! Selig, selig sind Mahomet und Rabbinen, die uns viel Besseres lehren. Ei, pfui dich unsinnigen Mahomet! sollst du ein Prophet heißen, der du ein solcher grober Tölpel und Esel bist?

142. Wohlan, diese elenden Narren laß fahren, und sich klug dünken, bis sie es genug haben; du aber halte fest am christlichen Glauben, der uns durch die Schrift lehrt, daß Jesus Christus sei wahrhaftiger GOtt, und GOttes Sohn, dazu auch wahrhaftiger Mensch, Davids und Marien Sohn; doch nicht zween Söhne, zween Männer, oder zwo Personen, sondern ein einiger Sohn, und einige Person, aus und in zwo unterschiedlichen Naturen, der Gottheit und Menschheit. Denn gleichwie droben, in dem Artikel von der Gottheit, du dich hüten mußt, daß du drei Personen nicht mengest in Eine Person, noch das Wesen oder Natur trennest in drei Götter, sondern drei unterschiedliche Personen in Einem göttlichen Wesen haltest: also mußt du dich wiederum hier hüten, daß du die einige Person nicht theilest oder trennest in zwo Personen, oder die zwo Naturen in Eine Natur mengest, sondern zwo unterschiedliche Naturen, in einer einigen Person, haltest. Und gleichwie die zwo Naturen sich in Eine Person vereinigen: also vereinigen sich auch die Namen beider Naturen in den Namen der einigen Person; welches man heißt Latein communicatio idiomatum vel proprietatum. Als, der Mensch heißt und ist geboren von der Jungfrau Maria, und von den Juden gekreuzigt; denselben Namen soll man auch GOttes Sohn geben, und sagen: GOtt ist geboren von Maria, und gekreuzigt von den Juden; denn GOtt und Mensch ist Eine Person, und nicht zween Söhne, einer GOttes, der andere Mariä; sondern ist ein einiger Sohn, GOttes und Marien.

143. Wenn du nun wolltest sagen, wie Nestorius, GOtt, oder Jesus, GOttes Sohn, sei nicht geboren von Maria, noch gekreuzigt von den Juden, sondern allein der Mensch, Marien Sohn, siehe, da machst du zwo Personen, und trennst die einige Person, daß eine andere Person sei, die geboren und gekreuzigt ist, und eine andere Person sei, die nicht geboren noch gekreuzigt ist, und wird also eine jegliche Natur für sich selbst eine sonderliche Person, und zween unterschiedliche Söhne; welches ist ebenso viel, als, daß GOtt sei nicht Mensch worden, sondern GOtt bleibe für sich selbst, eine abgesonderte Person von dem Menschen, und der Mensch für sich selbst, eine abgesonderte Person von GOtt. Das taugt nicht, das leidet die Schrift nicht, die da spricht Joh. 1, 14.: „Und das Wort ist Fleisch worden.“ Luc. 1, 35.: „Das in dir geboren wird, soll heißen des Allerhöchsten, GOttes Sohn.“ Und der Kinderglaube spricht: „Ich glaube an Jesum Christ, GOttes Sohn, der empfangen ist vom Heiligen Geist, geboren von Maria der Jungfrauen“ etc., spricht nicht, daß GOttes Sohn ein anderer sei, sondern derselbige, der von Marien geboren, und ihr Sohn wird.

144. Wiederum: Wenn du wolltest sagen wie Eutyches, daß der Mensch Jesus, Marien
\quellseite{988}
Sohn, sei nicht Schöpfer Himmels und der Erde, oder sei nicht GOttes Sohn, den man anbeten solle; wie auch neulich ein toller Geist große Grumpen vorgab, wie gefährlich wir Christen lebten, daß wir eine Creatur für GOtt anbeteten; der unsinnige Narr keine Schrift noch Bücher liest, sondern aus seinem eigenen tollen Kopfe von solchen hohen Sachen träumt, und ein selbstmachender Meister Klügel ist: Siehe, hier wird die Person abermal zertrennt, und zwo Personen aus Einer gemacht. Nestorius trennt die Personen also, daß er die Menschheit von der Gottheit reißt, und eine jegliche Natur eine sonder Person macht, daß allein der Mensch besonders gekreuziget sei. Eutyches wiederum reißt die Gottheit von der Menschheit, also, daß er auch eine jegliche Natur zur sonderlichen Person macht, daß GOtt solle angebetet werden, als von der Menschheit abgesondert. Aber die Schrift und der Glaube sprechen also: Wenn wir den Menschen, von Maria geboren, anbeten, so beten wir nicht einen abgesonderten Menschen an, der für sich selbst, ohne GOtt und außer GOtt, eine sonderliche eigene Person sei, sondern wir beten den einigen rechten GOtt an, der mit dem Vater und Heiligen Geiste ein einiger GOtt, und mit der Menschheit eine einige Person ist.

145. Wer nun solchen Verstand nicht hat, der muß in der Schrift irre werden, und kann sich nirgend drein schicken. Denn in der Schrift heißt Messia GOttes Knecht, Jes. 42, 1. „Siehe, das ist mein Knecht, an dem ich Wohlgefallen habe“; und am 53., V. 1. [Cap. 52, 13.]: „Siehe, mein Knecht wird klüglich fahren“; ja, er heißt „ein Wurm, und nicht ein Mensch“, Ps. 22, 7. Und, das wohl greulicher ist, er heißt ein armer Sünder, Ps. 41, 5.: „Ich sprach: HErr, sei mir gnädig, heile meine Seele, denn ich habe dir gesündigt.“ Ps. 69, 6.: „GOtt, du weißest meine Thorheit, und meine Schulden sind dir nicht verborgen.“ Item, V. 10.: „Die Schmach derer, die dich schmähen, fallen auf mich.“ Ps. 40, 13.: „Es hat mich Leiden ohne Zahl umgeben, mich haben meine Sünden ergriffen, daß ich nicht sehen kann, ihrer ist mehr, denn Haare auf meinem Haupte, und mein Herz hat mich verlassen.“ Hier schreiet Vernunft, Jude und Mahomet über uns Christen: Wie kann solches von GOtt verstanden werden? Wie kann GOtt ein Knecht sein? Wie kann er ein elender Sünder sein? Hilf GOtt, welch unsinnige, tolle, ungeheure Leute sind wir Christen vor solchen hohen, weisen, heiligen Menschen, die keine Creatur anbeten, sondern allein den einigen GOtt.

146. Wahr ist es, die Vernunft findet solches nicht in ihrer Bibel, das ist, im Rauchloch, oder im Schlauraffenlande. So finden es die Juden in ihrer Bibel auch nicht, das ist, im Talmud, unter den Säupizel, da sie ihr Schamhaperes innen studiren. So findet es Mahomet in seiner Bibel auch nicht, das ist, im Hurenbette; denn darin hat er am meisten studirt; wie er sich rühmt, der schändliche Unflat, daß ihm Gott (der Teufel) so viel Stärke seines Leibes gegeben habe, daß ihm vierzig Weiber nicht genug sein mögen zu Bette.\footnote{Vergleiche „Bruder Richards Verlegung des Alkoran“. Walch, St. Louiser Ausgabe, Bd. XX, 2289, § 72.} Ja, eben wie er studirt hat in derselben Bibel, im Hurenfleisch, so reucht und schmeckt auch sein teufisches Buch, der Alkoran; er hat den Geist seiner Prophetie am rechten Ort, im Venusberge, gesucht und gefunden. Wer nun in solchen Büchern studirt, was ist es Wunder, daß der nichts wisse von GOtt, noch Messia; so sie auch nicht wissen, was sie reden oder thun?

147. Wir Christen wissen (GOtt Lob und Dank in Ewigkeit!), daß Messia ist GOttes einiger ewiger Sohn, den er gesandt hat in die Welt, unsere Sünde auf sich zu nehmen, für uns zu sterben, und den Tod für uns zu überwinden; wie Jes. 53, 6. 10. klärlich sagt: „Wir gingen alle irre, GOtt aber hat unser aller Sünde auf ihn gelegt, und er hat sein Leben zum Schuldopfer gegeben“ etc. Daher singen und rühmen wir mit allen Freuden, daß GOttes Sohn, der rechte einige GOtt mit dem Vater und Heiligem Geist, sei für uns Mensch, ein Knecht, ein Sünder, ein Wurm worden; GOtt sei gestorben, GOtt trage unsere Sünde am Kreuz, in seinem eigenen Leibe; GOtt hat uns erlöst durch sein eigen Blut. Denn GOtt und Mensch ist Eine Person; was der Mensch hier thut, leidet und redet, das thut, leidet und redet GOtt; und was GOtt thut und redet, das thut und redet der Mensch, welcher ist einerlei Sohn, GOttes und Marien, in Einer unzertrennten Person, und zwo unterschiedlicher Naturen. Der Teufel, und sein Hurenwirth und Hurenjäger Mahomet, und seine Schamhaperisten, die Juden,
\quellseite{989}
mögen sich hieran ärgern, lästern, fluchen (wer es nicht lassen kann); aber allesamt sollen sie im Abgrund der Hölle ewiglich dafür zittern, zähnklappern und heulen, da nicht lang (ob GOtt will) hingehören soll, Amen.

148. Hier will ich aufhören diesmal von diesen hohen Artikeln zu reden aus dem Alten Testament, denn ich hoffe, es sollen hiermit genug vermahnt sein unsere Hebräisten, daß sie das Alte Testament sollen den Rabbinen nehmen, wo sie immer können, unangesehen ihre Deutung, Glossen, oder Grammatiken, weil die Rabbinen selbst unter einander oft hierin nicht wissen, wo sie daheim sind, und die vocabula und Sentenz gerne äquivociren auf ihren tollen Verstand, da doch der Buchstabe gerne mit dem Neuen Testament sich reimt, und gewiß ist, daß Jesus Christus der HErr über alles ist, dem die Schrift soll Zeugnis geben, als die allein um seinetwillen gegeben ist. Aus dem Neuen Testament habe ich diesmal nicht viel wollen führen, weil darinnen solches alles klärlich (darum es die Juden auch nicht annehmen) überzeugt ist, nun bei 1500 Jahren; und sonderlich im Evangelio Johannis, da schier über das andere Wort (wie man spricht) Jesus wird GOtt und Mensch, in Einer Person, gepredigt. Derselbe Johannes, sammt andern Aposteln, Evangelisten, und viel tausend ihrer Jünger, sind auch Juden, oder Israel und Abrahams Samen gewesen nach dem Geblüte, so wohl und viel reiner und gewisser, denn diese jetzigen Juden oder Israel sind, die niemand weiß, wer sie sind, oder woher sie kommen.

149. Wollen wir nun den Juden oder Israel glauben, so glauben wir billiger diesen Juden und Israel, welche nun bei 1500 Jahren öffentlich in aller Welt durch ihr Evangelium die Kirche regiert, Teufel, Tod und Sünde überwunden, die Schrift der Propheten ausgelegt, immer für und für durch ihre Jünger Wunder gethan haben. Billiger, sage ich, glauben wir solchen rechten, bekannten Juden und Israeliten, denn daß wir sollten glauben den falschen, unbekannten Juden oder Israeliten, die diese 1500 Jahre kein Wunder gethan, keine Schrift der Propheten ausgelegt, alles verkehrt, und im Lichte öffentlich nichts gethan, sondern in ihrem Winkel menschlings, wie die Kinder der Finsterniß, das ist, des Teufels, eitel Lästern, Fluchen, Morden und Lügen wider die rechten Juden und Israel (das ist, wider die Apostel und Propheten) geübt haben, und noch üben täglich: damit sie überwiesen sind, daß sie nicht Israel noch Abrahams Same, sondern giftige, teuflische Feinde sind des rechten Israels und Abrahams Kinder, dazu der heiligen Schrift Diebe, Räuber und Verführer. Darum man, als von öffentlichen Dieben, wieder nehmen soll die Schrift, wo es die Grammatica gerne gibt, und sich mit dem Neuen Testament reimt; wie die Apostel uns Exempel reichlich genug geben.

150. [Ich] will also wieder zu den letzten Worten Davids kommen, wie ich dies Büchlein anfing, und das Kränzlein mit dem Ende und Anfang zusammen binden, denn ich genug herumgeschweift habe. Andere können es und werden es (hoffe ich) besser machen, und den HErrn Christum in dem Alten hebräischen Testamente fleißig suchen, denn er läßt sich gerne drinnen finden, sonderlich in dem Psalter und Jesaia. Versuche es nach der Regel droben [§ 32 f.] gegeben, so wirst du mir wohl glauben, und GOtt danken. Nun, ich habe dies Büchlein also angefangen, daß Davids letzte Worte sollen vernommen und verstanden werden nach christlichem Verstande, auf diese Weise:

\begin{schriftzitat}
V. 1–3. Es sprach David, der Sohn Isai; es sprach der Mann, der von dem Messia des GOttes Jakob versichert ist, lieblich mit Psalmen Israel: Der Geist des HErrn hat durch mich geredet, und seine Rede ist durch meine Zunge geschehen; es hat der GOtt Israel zu mir gesprochen, der Hort Israel hat geredet, der gerechte Herrscher unter den Menschen, der Herrscher in der Furcht Gottes.
\end{schriftzitat}

151. Drei Redner sind hier (wie droben [§ 11] gesagt), der Geist des HErrn, der GOtt Israel, der Hort Israel; und ist doch ein einiger Redner. Aber bei dem dritten, das ist, bei dem „Hort Israel“, steht: „Der Herrscher unter den Menschen, der Herrscher in der Furcht GOttes.“ Dieser Herrscher ist Messia, wie auch der chaldäische Text gibt. Nun hangt es im Hebräischen an einander, nämlich: „Der Hort Israel, der gerechte Herrscher, der Herrscher in der Furcht GOttes.“ Gewiß ist es aber, daß \hebraeisch{יהוה} der Hort Israel heiße, und sei GOtt selbst; und ist doch auch Messia, der Mensch und Herrscher in der Furcht GOttes. „Herrscher“ ist hier hebräisch \hebraeisch{נגיד}, das heißt nicht HErr, wie GOtt HErr
\quellseite{990}
heißt; sondern, wie Menschen Herren sind und herrschen; und wo GOtt also genannt wird, magst du künstlich Jesum Christum daselbst verstehen. Als, da Gideon spricht Richt. 8, 23.: „Ich will nicht euer Herr sein, und mein Sohn soll auch nicht euer Herr sein, sondern der HErr soll euer Herr sein.“ Ps. 22, 29.: „Der HErr hat ein Reich, und er herrschet unter den Heiden.“ Ps. 59, 14.: „Der GOtt Jakob ist Herrscher in aller Welt.“ Also redet auch Psalm 8, 7. von Christo: „Du wirst ihn zum Herrscher machen über deiner Hände Werk, alles hast du unter seine Füße gethan“, das ist eben so viel, als hier David sagt, er sei „der Hort Israel“ (das ist GOtt) und „gerechter Herrscher“ (das ist Mensch) über alles, was GOtt gemacht hat. Das heißt, GOtt gleich sein, und doch auch Mensch sein.

152. Er nennt ihn einen „gerechten Herrscher“. Das ist nicht von weltlicher, zeitlicher Gerechtigkeit geredet, davon David sonst einen schönen Psalm gemacht hat, nämlich den 101. Psalm, „Von Gnade und Recht will ich singen“, sondern von der ewigen Gerechtigkeit, die Messia in die Welt gebracht, und uns von Sünden erlöst, und gerecht gemacht hat. Denn, wie folgt, redet er von dem ewigen Bunde, so GOtt mit Davids Hause gestiftet hat, wie es Jesaias auch verstehet am 55. Capitel, V. 3.: „Ich will euch die Gnade, David verheißen, treulich halten“; und Ps. 89, 3. 5.: „Es wird eine ewige Gnade gebaut werden, und du wirst deine Wahrheit treulich halten im Himmel.“ Hierzu ist weltliche Gerechtigkeit viel zu geringe, welche, wo sie am besten steht (das selten geschiehet), schwerlich den äußerlichen Frieden erhält, dem Mord, Raub, Ehebruch, Diebstahl \&c. wehrt, denn damit ist man noch nirgend vor GOtt gerecht, ob er sie wohl zeitlich und herrlich belohnt mit Reichthum, Ehre, Gewalt, Glück \&c., welches vor GOtt schlechte, geringe, vergängliche Parteken sind, die er auch seinen Feinden reichlicher gibt weder seinen lieben Kindern, welche einen bessern Lohn zu hoffen haben, davon die Welt nichts weiß.

153. Darum ist der Rabbinen und ihrer Nachfolger Verstand nichts, da sie meinen, es sei von David gesagt, daß er solle gerecht sein, und gottesfürchtig leben, weil er ein König und Herrscher gesetzt sei. Nein, es ist ein anderer Mann, dieser Herrscher in Gerechtigkeit und Gottesfurcht. David hat durch sein Regiment nicht einen einigen Menschen gerecht und gottesfürchtig gemacht, sich selbst auch nicht, wie mit seinem Gesetz auch nicht, Röm. 3. 10. 24., sondern sind alle gerecht und gottesfürchtig gemacht durch diesen Herrscher Messia und Hort Israel, Jesum Christum; wie auch Sacharja Cap. 9, 9. sagt: „Freue dich, du Tochter Zion, siehe, dein König kommt dir sanftmüthig, ein Gerechter und Heiland, und reitet auf einem Esel“ \&c. Und Paulus 1 Cor. 1, 30. 31.: „Er ist uns von GOtt gemacht zur Weisheit, Gerechtigkeit, Heiligung und Erlösung, auf daß, wer sich rühmet (wie geschrieben steht), der soll sich des HErrn rühmen“, und nicht unserer eigenen Gerechtigkeit, Weisheit \&c. Denn das ist seine Herrschaft. Darum ist er zum HErrn gesetzt, daß er solche Werke soll unter den Menschen thun, sie gerecht machen, und wieder zu Gottesfurcht, in Zucht und Gehorsam bringen, davon wir im Paradies gefallen sind durch der Schlange List. Von dieser Gerechtigkeit und Gottesfurcht ist hier nicht noch zu reden, es ist unsere tägliche Predigt, wie wir in Christo allein, aus lauter Gnaden, gerecht und selig werden. Folgt:\footnote{„Folgt“ fehlt in der Erlanger.}

\begin{schriftzitat}
V. 4. Und wie das Licht des Morgens, wenn die Sonne aufgehet des Morgens ohne Wolken, da vom Glanz nach dem Regen das Gras aus der Erde wächst.
\end{schriftzitat}

154. Er vergleicht die Herrschaft oder das Reich Messia, der die Gerechtigkeit und Gottesfurcht soll wieder bringen und aufrichten, dem schönen, lieblichen Wesen des Lenzens; denn der Winter, weil die Sonne weg von uns zeucht, die Erde mit Frost, Eis, Schnee \&c. zuschleußt, daß alle Bäume kahl, alle Gewächse matt werden, und nichts grünt, noch blüht, noch Frucht bringt, und alles\footnote{Wittenberger und Erlanger: als. Bei dieser Lesart würde das Subject des Satzes fehlen, und es müßte „die Welt“ ergänzt werden.} eine tote Welt anzusehen ist. Wenn aber gegen [den] Sommer die liebe Sonne wieder zu uns naht, so thut sich die Erde auf, grünt, blüht, reucht alles schön, wird alles neu, und die Welt gleichsam wieder lebendig und fröhlich. Denn alle Menschen, auch die Heiden, den Lenzens für das lustigste Theil des Jahres halten, wie Virgilius schreibt: Tunc formosissimus annus;\footnote{Virg. Bucolica, Ecl. III, 57: Nunc frondent silvae, nunc formosissimus annus.} und hält es dafür, daß die Welt
\quellseite{991}
im Lenzen habe angefangen, welches mit der heiligen Schrift stimmt, denn Moses den April zum ersten Monden des Jahres setzt. Also ist die Herrschaft und das Reich der Gnaden auch eine fröhliche, lustige Zeit, darinnen Messia uns gerecht und gottesfürchtig macht, daß wir grünen, blühen, wohlriechen, wachsen und fruchtbar werden. Denn er ist die Sonne der Gerechtigkeit, die wieder zu uns naht, wie Maleachi Cap. 4, 2. sagt: „Es soll euch, die ihr meinen Namen fürchtet, die Sonne der Gerechtigkeit aufgehen, und Heil unter seinen Flügeln.“ Darum er auch hat wollen leiblich im Lenzen oder April, in der luftigen Zeit, von den Todten auferstehen, und seine Herrschaft anfangen, ob er wohl im Winter geboren, das ist, um unsertwillen unter die Sünde, allerlei Jammer und Tod Adams sich gegeben, und also den harten Winter über 33 Jahre ausgestanden hat.

155. Denn gleichwie der Prophet David hier durch den Lenzen bedeutet die selige Zeit der Gnaden, so durch Messia, seinen Sohn, uns scheint: also gibt er damit zu verstehen, daß durch den Winter das Widerspiel, nämlich, die Zeit der Ungnaden unter der Erbsünde, die wir durch Adams Fall haben, bedeutet sei. Und hat GOtt also in seiner Creatur uns zum ewigen Gedächtniß bis an den jüngsten Tag (da andere Jahre, Erde und Himmel sein werden) vorgebildet die Sünde und Gnade, daß sie uns täglich und jährlich durch den Winter und Sommer vorgepredigt werden, wenn wir Ohren zu hören, und Augen zu sehen hätten. Solcher geistlicher Deutung nach ist Adam im schönen Lenzen erstlich gewesen (da er auch leiblich\footnote{Erlanger: leiblich. Wittenberger und Jenaer: leiblich. Cruciger: corporaliter.} im Lenzen, im Anfang des Jahrs, geschaffen ist), bald aber durch die Sünde den geistlichen Winter über sich gebracht hat, welchen Christus, die liebe Sonne, wiederum vertrieben, und den Lenzen angefangen. Und geht nun also: Wer im Lenzen lebt, der stirbt nimmermehr; wer im Winter stirbt, der lebt nimmermehr; das ist: „Wer glaubet und getauft wird, der wird selig, wer nicht glaubet, wird verdammt“ [Marc. 16, 16.]. Denn diesem entgeht die Sonne, jenem geht die Sonne auf, von welcher David hier weissagt.

156. Nicht allein das vom Winter und Sommer meint David, sondern vielmehr und näher dieses Geheimniß, daß Messia Herrschaft nicht solle sein, wie Mosis Herrschaft. Mosis Herrschaft ist des Gesetzes Herrschaft, welches nicht allein die Sünde nicht wegnimmt, sondern auch mehrt, das ist, offenbart sie, wie groß und greulich sie sei, und straft sie, dadurch der Mensch erschreckt, und (so zu reden) GOttes Gerichte und seinem Gesetze feind wird, durch welches er in seiner Sünde verdammt und getödtet wird; wie St. Paulus hiervon herrlich disputirt zum Römern [Cap. 7, 13.] und Galatern [Cap. 3, 21.]. Das ist der Berg Sinai, darauf es donnert, blitzt, regnet, erdenbebet, als wollte Himmel und Erde eingehen, und die Sonne viel tiefer hinter den finstern Wolken verborgen ist, ob es wohl in der Zeit des Lenzen ist, denn im Winter, da es zuweilen helle scheint, aber doch der Sonne Kraft zu ferne von uns ist. Denn die Heiden, so ohne Gesetz in den unbekannten Sünden des Winters leben, viel sicherer sind weder GOttes Volk, das auch zur Zeit des Lenzen des Gesetzes Donnern und Blitzen leiden muß. Denn wo die Sonne, Christus, nicht helle leuchtet, da ist auch der Lenz nicht lustig, sondern Mose macht es mit des Gesetzes Donner alles erschreckt und ganz tödtlich. Also sind die Wetter am Himmel auch ewige Propheten, daß uns zuweilen auch das Gesetz übereilt im Gewissen, die wir doch in der Zeit der Gnaden sind.

157. Aber hier, zu Messia Zeiten (spricht David), wenn der \hebraeisch{יהוה} selbst herrschen wird, mit Gnaden uns gerecht und selig zu machen, wird es so lieblich\footnote{Im Original: leiblich.} sein, als die beste Zeit im Lenzen, da es vor Tages einen lieblichen, warmen Regen gethan hat, das ist, das tröstliche Evangelium gepredigt ist; und flugs darauf die Sonne, Christus, aufgeht in unserm Herzen, durch rechten Glauben, ohne des Mosis Wolken, und Donner, und Blitzen. Da wächst es, grünt es, blüht es alles nach, und ist der Tag freudenreich und friedenreich, dergleichen das ganze Jahr nicht hat. Denn hier heißt es: Winter, Wolken, Donner, Sünde, Tod und alles Schrecken überwunden, und einen schönen, fröhlichen Ostertag gehalten bis in Ewigkeit. Siehe, das heißt David, seines Sohnes Messia Herrschaft gleich sein einem Tage im Lenzen, da es frühe geregnet, darauf die Sonne aufgeht aufs allerlieblichste, und macht grün, blühend,
\quellseite{992}
riechend, und alles lebendig und fröhlich. Trage dich selber darum, ob es nicht die beste und fröhlichste Zeit im Jahr sei. Folgt:

\begin{schriftzitat}
V. 5. Denn mein Haus ist nicht also bei GOtt; denn er hat mir einen Bund gesetzt, der ewig, und alles wohl geordnet und gehalten wird.
\end{schriftzitat}

158. Droben ist gesagt aus 1 Chron. 18, 16., daß diese Worte: „Allein Haus ist nicht also bei GOtt“, solle so viel heißen: „Ach, was bin ich? was ist mein Haus, gegen GOtt?“ Es ist ja nicht ein solch Haus, das bei GOtt solcher unaussprechlichen Ehren würdig sei, von welchem Messia, der \hebraeisch{זרע}, Gottes Sohn, der gerechte Herrscher unter den Menschen, soll geboren werden. Und fällt hiermit David herunter in große Demuth und Verwundern, daß von seinem Fleisch und Blut solche großen Dinge kommen sollen.

159. Das andere Stück, vom ewigen Bunde und Hause Davids, habe ich in dem Büchlein von den Juden ziemlich gering gehandelt,\footnote{Walch, St. Louiser Ausgabe, Bd. XX, 2015 ff.} und damit andern Ursache gegeben, weiter und besser zu handeln. Die folgenden zwei Worte, \hebraeisch{צדקה} und \hebraeisch{שכרה}, wohl geordnet und gehalten, sind mit Fleiß gesetzt, zur Lehre und Trost. Denn, so du ansiehst die Historien, wird dich dünken, GOtt habe seines Bundes vergessen, und den nicht gehalten. So gar wüst, unordig und seltsam geht es zu in Davids Hause und Nachkommen. Noch ist es bis auf Messia nicht allein gehalten, sondern ist alles in seiner Ordnung wohl und fein geblieben, wider alle Teufel und Menschen. Und hat es niemand können ändern noch dämpfen, sondern haben es müssen lassen gehen und stehen, ein Scepter Juda, wie es verheißen ist, bis auf Messiam.

160. Aber nach Messia ist sein Reich, die Kirche, äußerlich anzusehen, viel wüster und unordiger, daß kein zerrissener, elender, nichtiger Regiment oder Herrschaft ist, denn die christliche Kirche, Christi Herrschaft. Hier zerreißen es und [ver]wüsten es mit Feuer, Wasser, Schwert und aller Macht die Tyrannen. Hier verwühlen und verderben es die Rottengeister und Ketzerien. So machen es die falschen Christen mit ihrem bösen Leben auch also, als sei keine schändlichere, ungeordnetere Herrschaft auf Erden. Und diese alle arbeiten dahin, oder vielmehr der böse Geist durch sie, daß Christi Herrschaft solle nichts, oder je ein elend, unordig Ding sein. Und Summa, Christus stellt sich, als habe er seiner Herrschaft vergessen, und sei nirgend heime, daß hier weder \hebraeisch{עריה} noch \hebraeisch{שכרה} von der Vernunft gesehen wird; noch heißt sie \hebraeisch{צדקה בכל} und \hebraeisch{שכרה}, alles wohl geordnet und gehalten. Ob wir es nicht sehen, so sieht er es, der da spricht Hohel., Cap. 8, 12.: „Mein Weinberg ist vor mir.“ Matth. 28, 20.: „Siehe, ich bin bei euch bis an der Welt Ende.“ Joh. 16, 33.: „Seid getrost, ich habe die Welt überwunden.“ Gleichwohl sehen wir, daß immer geblieben ist und bleibt ein Volk, das den Namen Christi ehrt, sein Wort, Taufe, Sacrament, Schlüssel und Geist hat, auch wider alle Pforten der Hölle. Folgt:

\begin{quote}
Denn alle mein Heil und Thun ist, daß nichts wächst.
\end{quote}

161. Er will sagen: Ich bin auch ein König und Herr, vor allen Königen, von GOtt geordnet, habe viel Kriege geführt, groß Glück, Sieg und Heil gehabt durch Gottes Hülffe und Wundertat, auch viel gethan im Regiment, wohl regiert, das Reich wohl angerichtet und gestellt, das Recht gehandhabt, darüber auch viel erlitten; aber solch mein Reich, vielmehr aller Könige auf Erden Reich, gegen dieser Herrschaft meines Sohnes Messia, des \hebraeisch{זרע}, ist nichts denn ein dürrer Zweig, der nie einmal gewachsen oder gegrünt zu rechnen ist. Denn ich und kein König hat den Sieg wider den Tod, Sünde, Hölle, Teufel und Welt erobert; so hat auch unser keiner das gethan in seinem Regimente, daß er die Leute gerecht, gottessürchtig und ewig selig machte. Wir sind arme, dürftige, dürre Herren in unserm Regimente; aber mein Sohn Messia, der \hebraeisch{זרע}, das ist der Mann, der hat den Sieg über Sünde, Tod, Teufel, Hölle, Welt und alles erhalten, der hat ein Regiment, darinnen er das thut und anrichtet, daß er alle die Seinen ewiglich gerecht und selig macht; das heißt gegrünt, geblüht, fruchtbar sein, und daß nimmermehr verdorren kann.

162. Ich habe das Wort \hebraeisch{חפץ} verdeutscht „thun“, nach dem Exempel Salomo, Pred. 3, 1.: „Ein jegliches hat seine Zeit, und alles Vornehmen (das ist, Thun) hat seine Stunde.“ Denn so redet man deutsch: Du mußt ja etwas
\quellseite{993}
vornehmen, daß du nicht müßig gehest, sondern etwas thuest, damit du dich nährest. Also heißt „Thun“ allerlei Stand, darin sich einer übt in diesem Leben; und die Philosophi heißen es auch placita, proposita, instituta, darum, daß einem dies, dem andern das gefällt zu thun oder vorzunehmen. Denn \hebraeisch{רצון} heißt eigentlich, ein Gefallen oder Willen, Lust, Neigung zu einem Dinge haben; denn wer nicht Willen dazu hat, der thut es nicht, oder thut es, daß [es] wohl so gut, als ungethan wäre.

\begin{schriftzitat}
V. 6. 7. Aber Belial sind allesamt wie die ausgeworfenen Disteln, die man nicht mit Händen fassen kann; sondern wer sie angreifen soll, muß Eisen und Spießstangen in der Hand haben, und werden mit Feuer verbrannt in der Wohnung.
\end{schriftzitat}

163. Hier weissagt er von den Juden, die solchen HErrn und Messia nicht würden annehmen, und nennt sie \hebraeisch{בליעל}, oder, wie wir es gewohnt, Belial, das heißt auf deutsch: unnütz oder schädlich. Nach dem äußerlichen Regiment heißt man es unnütze, böse Buben, die gerne Schaden thun. Aber David redet hier im Geiste vom Reiche Christi, da hat es diese Weise, daß die, so den Reiche Christi feind sind, als Juden, Ketzer, Heiden, für die Allernützigsten gehalten werden. Denn auch noch jetzt die Juden, Mahomet, Pabst, Rottengeister sich dünken lassen, sie thun eitel Gottesdienst daran, wo sie den rechten Christen Schaden thun können. Diese wollen nicht \hebraeisch{בליעל} heißen, sondern allein die Allernützigsten sein. Also sagt Jeremia Cap. 23, 32. von den falschen Propheten: „Sie sind mit ihrem Nutzen kein nütze diesem Volk“, das ist, sie\footnote{„sie“ fehlt in der Wittenberger und in der Jenaer.} sind die Schädlichsten, eben da sie die Nütze sten sein wollen. Summa, die Christen sind Belial und Teufels Kinder, diese aber sind allein GOttes Kinder; was sie thun, das ist recht, bis sie GOtt ausstößt, und mit Feuer seines Zorns verbrennt; wie wir an den jetzigen Juden sehen, welch ein greuliches Feuer göttliches Zorns über sie gekommen ist.

164. Er vergleicht sie den Disteln im Korn auf dem Felde, welche, meines Achtens, Christus der HErr Matth. 13, 25. zizania nennt, welches wir „Unkraut“ verdeutscht haben. St. Ambrosius in Hexa. spricht: Ex tritici semine degenerans in proprium genus, das heißen wir Trespen. Aber Christus redet von einem Aergern, daß man in der Ernte aussondert vom Getreide, und mit Feuer verbrennt, und schier mit David einerlei Worte braucht, der seine Disteln auch aussondert und mit Feuer verbrennt. Darum wird zizania hier sein, das David \hebraeisch{קוץ} heißt, die großen, bösen, stachligen Disteln, oder die andern Disteln, so unsere Bauern toll Graet heißen, die man mit Sensen, Sicheln und Harken oder Spießholz aussondert\footnote{Wittenberger und Jenaer: aussondern muß.} in der Ernte, denn mit Händen kann sie alle beide niemand angreifen, und dienen nirgendhin, denn ins Feuer; aber Trespen braucht man für das Vieh. Also sind die verstockten Juden solche bösen, stachligen Disteln und toll Graete, daß sie mit keiner Wohlthat noch Wunder GOttes zu bekehren gewesen und noch sind, sondern durch der Römer Eisen und Spieße ausgetilget, und mit ihrer Stadt in ihrer eigenen Wohnung verbrannt, auch mit leiblichem Feuer. Ueber das brennen sie noch jetzt in ihnen selber, wo sie im Elende sind, mit geistlichem Feuer göttliches Zorns. Also hat David diesem Volke seine Verstörung und endliches Verderben verkündigt, darum, daß sie diesen König nicht haben wollten, wie der HErr Luc. 19, 43. 44. auch davon weissagt, und Dan. 9, 26. Sach. 14, 2.

165. Hiemit will ich die letzten Worte Davids verdeutscht und ausgelegt haben nach meinem eigenen Sinne. GOtt gebe, daß unsere Theologen getrost Hebräisch studiren, und die Bibel uns wieder heimholen von den muthwilligen Dieben, und alles besser machen, denn ich’s gemacht habe; das ist, daß sie den Rabbinen sich nicht gefangen geben in ihre genarrte Grammatica und falsche Auslegung, damit wir den lieben HErrn und Heiland hell und klar in der Schrift finden und erkennen.

\begin{quote}
Dem sei Lob und Ehre sammt dem Vater und Heiligen Geist in Ewigkeit, Amen.
\end{quote}

Ende des dritten Theils.
